% 本文件由 LaTeX 排版 Agent 根据有效 translation.json 自动生成。
% author=Augustine of Hippo; work_id=1512; title=论圣人的预定

\chapter{论圣人的预定}

% structure: p0001 source=h1 target=omitted reason=page-title-already-rendered-by-parent

\SourceRecord{Translated by Peter Holmes and Robert Ernest Wallis, and revised by Benjamin B. Warfield.；Nicene and Post-Nicene Fathers, First Series；Vol. 5.；Edited by Philip Schaff.；Buffalo, NY: Christian Literature Publishing Co.,；1887.；<http://www.newadvent.org/fathers/1512.htm>.}

% structure: p0001 source=h1 target=section reason=page-title

\section{论圣人的预定（卷一）}

\Emph{致普罗斯柏和希拉里。}\par

\Emph{本书旨在为预定和恩宠的真理辩护，反对半白拉奇派——即那些绝不完全脱离白拉奇异端的人，因为他们主张救恩和信德的开端来自我们自己；因此，仿佛基于此前的功绩，天主其他美好的恩赐才得以获得。奥古斯丁指出，不仅信德的长进，就连信德的开始也是天主的恩赐。关于此事，他不否认自己曾有过不同想法，并且在主教任期前所写的一些小作品中存在错误，例如他们在反驳他时所指出的那篇关于《罗马书》段落的解释。但他指出，后来他主要是被这段证词所说服：\Quote{你有什么不是领受的呢？}他证明这证词也应被视为关于信德本身的证词。他说信德应被算作其他行为之一，而宗徒在说\Quote{不是出于行为}时否认它先于天主的恩宠。他宣告，心的刚硬被恩宠除去，凡受父教导而来的都来到基督面前；但祂以慈悲教导那些人，而祂所不教导的人，则以审判不教导。至于他从第一百零二封信第二问\Quote{关于基督宗教时期}中引用的段落——半白拉奇派以此指控——可以正确地解释，而不损害恩宠和预定的道理。他教导恩宠和预定之间的区别。此外，他说天主在祂的预定中预知了祂定意要做的事。他极为惊奇，预定的反对者们，据说他们不愿依赖天主旨意的不确定性，却宁愿将自己的软弱置于天主应许的力量之上。他清楚指出，他们滥用了\Quote{你若相信，就必得救}这个权威。在得救的婴孩身上，恩宠和恆心的真理大放光明，这些婴孩与那些丧亡的婴孩并无任何自己的功绩可以区分。因为在他们之间，没有任何区别是来自他们若活得更长可能拥有的功绩的预知。反对者们错误地拒绝接受他为讨论而引用的证词，即\Quote{他已被取去，免得邪恶……}等，认为它不是正典。预定和恩宠最光辉的例子就是救主本身，在其中一个人获得了作为救主和天主独生子的特权，这是由于被与圣父同永恒的圣言所取纳为同一性体，而非基于任何行为或信德的先有功绩。预定者被某种特殊选民的呼召所召唤，他们在创世之前就被拣选；不是因为预知他们将要相信并成为圣洁，而是为了借着那拣选的恩宠本身，他们成为这样的人，等等。}\par

% structure: p0004 source=h2 target=subsection reason=relative-heading-rank; base=h2

\subsection{第一章 引言}

我们知道宗徒在致斐理伯人书中有话说：\Quote{我把同样的事写给你们，在我并不厌烦，在你们却是稳妥的。}（\BibleRef{斐 3:1}）然而，同一位宗徒在致迦拉达人书时，当他看到自己通过宣讲的职务在他们中已做了他认为为他们所必需的事，就说：\Quote{从今以后，不要有人使我劳苦。}（\BibleRef{迦 6:17}）或如许多抄本所读的：\Quote{不要有人麻烦我。}但我承认，当神的话语——其中宣讲天主的恩宠（若按我们的功绩赐予，则绝非恩宠）——虽如此伟大而显明，却仍不被接纳时，这确实使我劳苦。然而，我最亲爱的儿子们，普罗斯柏和希拉里，你们的热情和弟兄之爱——这使你们如此不愿看到任何弟兄在错误中，以至于希望我在写了那么多关于这主题的书信之后，仍从远处再次写信——我爱之不尽，尽管我不敢说我爱它如同应该的那样。所以，看哪，我再给你们写信。虽然不与你们同在，却藉着你们，我仍在做我以为已做得足够的事。\par

% structure: p0006 source=h2 target=subsection reason=relative-heading-rank; base=h2

\subsection{第二章 马西利亚人如何脱离白拉奇派}

因为考虑到你的信件，我似乎看到那些弟兄们——你对他们怀有敬虔的关怀，以免他们持有那种诗意的观点（即断言\Quote{每个人都为自己存有希望}），从而陷入那非诗意而是先知所宣告的谴责——\BibleQuote{凡以人为希望的人，是可咒骂的}\BibleRef{耶 17:5}——必须按照宗徒对待那些他对他们说\BibleQuote{如果你们在什么事上另有想法，天主也会将这事启示给你们}\BibleRef{斐 3:15}的人的方式来处理。因为目前他们在关于圣人预定的问题上仍处于黑暗之中，但他们拥有那种根据——\BibleQuote{如果他们在什么事上另有想法，天主也会将这事启示给他们}——只要他们行走在他们已达到的地步上。因此，当宗徒说了\BibleQuote{如果你们在什么事上另有想法，天主也会将这事启示给你们}之后，他又说：\BibleQuote{然而，我们已达到什么地步，就当照此而行}\BibleRef{斐 3:16}。而我们的那些弟兄们——你虔诚的爱意对他们挂念——已与基督的教会达到了这样的信仰：人类生来就负有第一个人的罪的罪责，并且除非通过第二个人的正义，无人能从那邪恶中得救。此外，他们也达到了这样的宣认：人的意志被天主的恩宠所抢先；也同意了无人能凭自己起始或完成任何善工。因此，他们所已达到的这些事，如果牢牢持守，就能充分地将他们与白拉奇主义者的错误区分开来。此外，如果他们行走在其中，并恳求赐予领悟的那一位，如果他们在关于预定的事上另有想法，祂也会将这事启示给他们。然而，让我们也按照祂的恩赐——我们曾请求祂在这封信中让我们说出对他们适宜和有益的话——向他们倾注我们的爱心和话语的侍奉。因为，我们如何知道，通过我们这项事奉——我们在基督的自由之爱中服侍他们——我们的天主或许不会愿意实现那个目的呢？\par

% structure: p0008 source=h2 target=subsection reason=relative-heading-rank; base=h2

\subsection{第三章 —— 甚至信德的开始也是天主的恩赐。}

因此，我应当首先表明，我们成为基督徒的信德是天主的恩赐，如果我能比我已在这么多、这么大的著作中已经做的更详尽的话。但我看出，我现在必须回答那些说我所引用的有关此事的圣经证据只足以保证我们确信我们自己拥有信德本身，但其增加是出于天主；仿佛信德不是由祂赐给我们，而是仅凭我们从自己开始的信德的功劳由祂在我们内增加似的。因此，这里并没有背离白拉奇本人在巴勒斯坦主教们的审判中被谴责的观点，正如同一会议记录所证实的：\Quote{天主的恩宠是按我们的功劳赐给的}，如果信德的开始不是来自天主的恩宠，而是因着这个开始，我们才被加添更圆满、更完全的信德；这样，我们首先将我们信德的开始献给天主，以便祂的补充也再次赐给我们，以及我们虔诚祈求的任何其他事物。\par

% structure: p0010 source=h2 target=subsection reason=relative-heading-rank; base=h2

\subsection{第四章 —— 接续上文。}

但为什么我们反而没有听到这些话：\BibleQuote{谁先给了祂，以致祂要偿还呢？因为一切都是出于祂，借着祂，并在祂内。}\BibleRef{罗 11:35} 那么，我们信德的那个开始，若不是出于祂，又是出于谁呢？因为当其他事物被说是出于祂时，这并不例外；而是\Emph{一切}都是出于祂、借着祂、并在祂内。然而，谁能说，那已经开始相信的人，在他所相信的那一位面前，毫无功劳呢？由此得出，对于已经有功劳的人，其他事物被认为是出于神圣的回报而添加的，因此天主的恩宠是按我们的功劳赐给的。白拉奇本人听到这话后，就谴责了它，以免自己被定罪。因此，凡要从各方面避免这种可谴责之观点的人，就当明白宗徒所说的话是完全真实的：\BibleQuote{因为赐给你们的恩宠，为了基督，不仅是要你们相信祂，而且还要为祂受苦。}\BibleRef{斐 1:29} 他表明两者都是天主的恩赐，因为他说了两者都是被赐予的。他并没有说\Quote{更圆满、更完全地相信祂}，而是说\Quote{相信祂}。他也没有说他自己蒙受了仁慈是为了成为更忠信的，而是说\BibleQuote{成为忠信的}\BibleRef{格前 7:25}，因为他知道，他并没有首先将自己信德的开始献给天主，然后从祂那里得到增长的回报；而是他被天主造就成了忠信的，这位天主也使他成为了宗徒。因为他信德的开始是有记载的，并且通过阅读教会中的经文（一个适合区分它们的场合）而非常有名：他如何背离了他所毁灭的信德，并强烈反对它，却突然被一种更强大的恩宠（借由祂的转化，关于这位要行此事的天主，先知曾说过：\Quote{祢将转化我们，使我们获得生命}）所转变，以至于不仅从一个不信者成为一个自愿相信的人，而且从一个迫害者，反而为捍卫他所迫害的信德而受苦。因为基督赐给他\BibleQuote{不仅是要相信祂，而且还要为祂受苦}。\par

% structure: p0012 source=h2 target=subsection reason=relative-heading-rank; base=h2

\subsection{第五章 —— 相信就是带着赞同的思考。}

因此，祂赞扬那并非按照任何功绩赐予、却是一切善功之原因的天主恩宠，说：\BibleQuote{这并不是说我们凭自己能承担什么事，因为我们所能承担的，乃是出于天主。}\BibleRef{格后 3:5} 愿那些认为信德的开端出于我们自己、而信德的成全出于天主的人，留意这话并深思之。因为谁看不出\Quote{思想}先于\Quote{相信}？因为除非先思想那该信的，没有人会相信任何事。因为无论思想如何在意志相信之前多么突然、多么迅速地掠过，并且相信紧随其后，仿佛与之紧密相连，但一切被相信的事必然是先有思想后才相信的；然而，连相信本身也无非是带着赞同的思想。因为并非所有思想的人都相信，因为许多人思想是为了不相信；但每个相信的人都思想——他在相信中思想，也在思想中相信。因此，在关乎宗教与虔敬的事上（宗徒所谈论的），如果我们不能凭自己承担什么思想，而是出于天主，我们当然也不能凭自己承担什么相信，因为我们不能在没有思想的情况下做到这一点；但我们开始相信的能力是出于天主。所以，正如没有人凭自己足以开始或完成任何善工——你们那些弟兄，如你们来信所表明的，已同意这是真的，因此，在每一善工的开始和进行中，我们的能力是出于天主——同样，没有人凭自己足以开始或成全信德；但我们的能力是出于天主。因为如果信德不是思想之事，它就毫无价值；而我们不能凭自己承担什么思想，而是出于天主。\par

% structure: p0014 source=h2 target=subsection reason=relative-heading-rank; base=h2

\subsection{第六章——应避免僭越与傲慢}

蒙天主所爱的弟兄们，务必注意：当一个人说他在做天主所应许的事时，切不可高举自己反对天主。外邦人的信德岂不正是向亚巴郎应许的吗？ \BibleQuote{他因信德而坚强，归光荣于天主，且满心相信天主所应许的，也必能成就。}\BibleRef{罗 4:20} 因此，那能成就祂所应许的，正是祂自己，祂使外邦人的信德得以实现。再者，如果天主在我们心中以一种奇妙的方式运作我们的信德，使我们相信，那么有什么理由担心祂不能完成全部呢？而人是否会因此将最初的因素归功于自己，以配得从天主那里领受最后的因素？请想一想，若这样行事，除了使天主的恩宠以某种方式按照我们的功绩赐予、因而恩宠不再是恩宠之外，还能得到什么结果？因为按此原则，恩宠被当作债务偿还，而非白白赐予；因为相信者理应让主增加其信德，而增加的信德是已开始之信德的报酬；这样说的时侯，并未注意到这报酬是作为债务而非恩宠给予相信者的。我完全看不出为什么不应将全部归功于人——因为那能为自己创造先前所没有的，也能增加他所创造的——除非无法抵挡最明显的神圣证言，这些证言表明信德（虔敬由此开始）也是天主的恩赐：例如这样的证言：\BibleQuote{天主按各人信德的大小，分给了各人。}\BibleRef{罗 12:3} 以及另一个：\BibleQuote{愿平安与爱德连同信德，由天主父和主耶稣基督赐与众弟兄。}\BibleRef{弗 6:23} 以及其他类似经文。因此，人既不愿抗拒这些如此清楚的证言，又渴望自己拥有相信的功绩，便与天主讨价还价：他为自己保留一份信德，给天主留下一份；而更傲慢的是，他把最初的一份归给自己，把后续的一份归给天主；这样，在他所说的两者共有的信德中，他使自己为首，使天主为次！\par

% structure: p0016 source=h2 target=subsection reason=relative-heading-rank; base=h2

\subsection{第七章 [III.]——奥思定承认他从前关于天主恩宠的见解有误}

那位虔诚谦逊的导师——我指的是\Person{圣西彼廉}——并非如此思想，他曾说：\Quote{我们不可在任何事上夸耀，因为没有什么属于我们自己。}为了证明此点，他引证宗徒的话说：\BibleQuote{你有什么不是领受的呢？既然是领受的，为什么你夸耀，好像不是领受的呢？}\BibleRef{格前 4:7}正是主要凭此见证，我本人也深信不疑——当时我处于类似错误之中，以为我们信天主的信德并非天主的恩赐，而是来自我们自己，并藉此获得天主的恩赐，使我们能在此世度节制、正义和虔诚的生活。因为我并不以为信德是以天主的恩宠为先决条件的，以致藉此恩宠赐予我们可有益地祈求的事物，除非真理的宣讲先行，我们无法相信；但我以为，当福音传给我们时，我们应同意乃是我们自己的行为，来自我们自身。我的这一错误在我担任主教前所写的一些小作品中已充分表明。其中包括你们在信中所提到的那一本，其中诠释了\BibleBook{}{罗马书}中的某些命题。最终，当我撤回我所有的小作品，并将那撤回之辞付之文字时——在着手你们更长的书信之前，我已完成了两部著作——当在第一卷中我论及此书之撤回时，我这样说：——我也讨论说：\Quote{天主在尚未出生的人中可能选择了谁，祂说年长的要服事年幼的；又在同样尚未出生的年长者中，祂可能拒绝了谁；关于此人，因此有先知证词记载，尽管是后来才宣布的：\BibleQuote{我爱了雅各伯，恨了厄撒乌。}}我推理到如此地步，说：\Quote{因此天主并非预知祂自己要赐予他们什么而选择了任何人的工作，而是预选了信德，预知祂会选择那个祂预知会相信祂的人——祂将赐予圣神给他，使他因行善工而也能获得永生。}我尚未仔细探求，也尚未找到恩宠的拣选的性质，宗徒说：\BibleQuote{按照恩宠的拣选，残余的人得蒙拯救。}\BibleRef{罗 11:5}这无疑不是恩宠，若任何功绩先于它，免得现在所赐予的，不是按照恩宠，而是按照债，与其说是白白赐予，不如说是酬报功绩。而此后我又附加说：\Quote{因为同一位宗徒说：\BibleQuote{同一天主在一切人身上行一切事。}\BibleRef{格前 12:6}但从未说过：天主在一切人身上相信一切事。}然后又说：\Quote{所以，我们所信的是我们自己的，但我们所行的善事，是出于那赐圣神给相信者的一位。}若我已知道信德本身也属于那由同一圣神所赐的天主恩赐，我肯定不能说这样的话。因此，这两者都属于我们，是由于意志的选择，而两者又都是由信德与爱德之神所赐予的。因为信德并非单独存在，而是如经上所载：\BibleQuote{爱德连同信德，来自天主父及我们的主耶稣基督。}\BibleRef{弗 6:23}而我在稍后所说的话：\Quote{因为相信和愿意是我们的，但赐给那些相信和愿意者行善工的能力，是祂的事，藉着圣神，爱德倾注在我们心中。}——这固然是真实的；但按同一法则，两者也都是天主的，因为天主预备意志；两者也是我们的，因为它们只有在我们良好的意志下才得以实现。于是我又接着说：\Quote{因为我们除非被召唤，便不能愿意；而当我们被召后愿意时，我们的愿意和奔跑都不足够，除非天主赐力量给我们奔跑者，并引领我们往祂召唤我们去的地方。}遂又加上：\Quote{所以，显然不是出于愿意者，也不是出于奔跑者，而是出于显示仁慈的天主，我们才能行善工。}——这绝对是至真之理。但关于那按天主旨意而来的召唤本身，我发现的很少；因为并非所有被召者都如此受召，而只有选民是如此。因此我稍后所说的话：\Quote{因为在天主所拣选的人中，不是工作而是信德开始功绩，以藉天主的恩赐行善工；同样，在祂所定罪的人中，不信和邪恶开始惩罚的功绩，以致甚至藉惩罚本身他们行恶事。}——我这话说得极其正确。但连信德的功绩本身也是天主的恩赐，我未曾考虑去探究，也没有说。而在另一处我说：\Quote{因为祂对谁施予仁慈，就使他行善工；祂使谁心硬，就任由他行恶事；但那仁慈是赐予在先的信德功绩上，而那心硬是加于先前的邪恶上。}这固然也是真实的；但还应进一步探问，是否连信德的功绩也并非来自天主的仁慈——即是说，那仁慈只因为那人已是信徒而显明在他身上，还是也为了使他成为信徒而显明？因为我们读到宗徒的话：\BibleQuote{我蒙受了仁慈，为成为信徒。}\BibleRef{格前 7:25}他并没有说：\Quote{因为我原是信徒。}因此，虽然它赐给了信徒，但也是为使他能成为信徒而赐予的。因此，在同一书中的另一处，我也极真实地说：\Quote{因为，若是出于天主的仁慈，而不是出于工作，我们甚至被召叫，为能相信，并且我们相信者获赐行善工，那么那仁慈也不应吝啬于外邦人。}——虽然我在那里对那按天主旨意而赐予的召唤论述得不够谨慎。\par

% structure: p0018 source=h2 target=subsection reason=relative-heading-rank; base=h2

\subsection{第八章 [第四部分]——奥古斯丁致米兰主教盎博罗削的继任者辛普利齐亚努斯的信。}

你清楚看到，当时我对信德和行为的看法是什么，尽管我那时致力于赞扬天主的恩宠；而现在我看到我们的那些弟兄们正是持这种看法，因为他们没有像仔细阅读我的著作那样小心地在其中与我一同进步。因为他们若是那样仔细，便会发现那问题已按照天主圣经的真理得到解决，就在我牧职之初写给\Place{米兰}教会主教、可敬的\Person{辛普利西安}的书中，那两部书的第一部中；\Person{辛普利西安}是\Person{圣盎博罗削}的继任者。除非，也许他们不知道这些书；既然如此，要留心使他们知道。我首先是在\WorkTitle{更正篇}第二卷中讲论了这两部书中的第一部；我当时所说的是如下：\Quote{我说，关于我作为主教所劳苦的那些书，有两部是写给\Place{米兰}教会的主席\Person{辛普利西安}的，他继承了最蒙福的\Person{盎博罗削}，讨论了若干问题，其中两个问题我收集在第一卷中，取自宗徒保禄致罗马人书。第一个问题是关于所记载的：\BibleQuote{那么，我们可说什么呢？法律是罪吗？绝对不是！}\BibleRef{罗 7:7}直到他说：\BibleQuote{谁能救我脱离这死亡的身体呢？天主的恩宠藉我们的主耶稣基督。}\BibleRef{罗 7:24}在其中我解释了宗徒的那些话：\BibleQuote{法律是属神的，但我是属肉体的，}\BibleRef{罗 7:14}以及其他的话，在那些话中，指明肉体与圣神相争，好像一个人仍处于法律之下，尚未建于恩宠之下。因为很久以后，我觉察到那些话甚至可能（并且大概是）一个属神的人的言论。本书中的第二个问题取自宗徒所说的那段话：\BibleQuote{不仅如此，而且当黎贝加由一次同房怀孕，即由我们的父依撒格怀孕后，}\BibleRef{罗 9:10}直到他说：\BibleQuote{若非万军的上主给我们留下了苗裔，我们早已像索多玛，相似哈摩辣了。}\BibleRef{罗 9:29}在解决这个问题时，我的确为人的自由意志而辛劳，但天主的恩宠得胜了，我最终只能达到宗徒以最明显真理所说的那句话：\BibleQuote{谁使你与众不同呢？你有什么不是领受的呢？既然是领受的，为什么你还要夸耀，好像不是领受的呢？}\BibleRef{格前 4:7}殉道者\Person{西彼廉}也渴望阐释这一点，他在那个标题中概括了全部：\InnerQuote{我们不可在何事上夸耀，因为没有什么是我们自己的。}}这就是为什么我先前说过，正是在这方面我曾持有不同看法时，主要藉着这宗徒的见证使我信服；而这是天主在我寻求解答此问题时启示给我的，正如我所说，是在我写信给\Person{辛普利西安}主教的时候。因此，宗徒为了抑制人的自负而说的\BibleQuote{你有什么不是领受的呢？}\BibleRef{格前 4:7}这见证，不容许任何一个信徒说：\Quote{我有信德，是我没有领受的。}这种回答的所有傲慢都被这些宗徒的话完全压制了。而且，甚至不能说：\Quote{虽然我没有完美的信德，但我有它的开端，藉此我首先相信了基督。}因为在此也同样回答了：\BibleQuote{你有什么不是领受的呢？既然是领受的，为什么你还要夸耀，好像不是领受的呢？}\par

% structure: p0020 source=h2 target=subsection reason=relative-heading-rank; base=h2

\subsection{第九章 [V.]—— 宗徒这些话的目的。}

然而，他们所持的看法，即\Quote{即\InnerQuote{你有什么不是领受的呢？}这句话不能用于这个信德，因为信德虽然已经败坏，却仍留在同一本性内，这本性起初被赋予了健全与完美}，若考虑到宗徒为何说这些话，便可见这看法在他们所期望的目的上毫无力量。因为他关心的是，没有人应在人身上夸耀，因为在\Place{格林多}的基督徒中发生了分裂，以致每个人都说：\BibleQuote{我确实是属\Person{保禄}的，另一人说我属\Person{阿颇罗}，另一人说我属\Person{刻法}。}\BibleRef{格前 1:12}然后他接着说：\BibleQuote{天主拣选了世上愚拙的，叫有智慧的羞愧；又拣选了世上软弱的，叫那强壮的羞愧；天主也拣选了世上卑贱的，被人轻视的，以及那无有的，为要废掉那有的，使一切有血气的，在天主面前一个也不能自夸。}\BibleRef{格前 1:27}这里宗徒的意图显然是足够清楚的，反对人的骄傲，叫没有人应在人身上夸耀；同样，也不应在自己身上夸耀。最后，当他说\BibleQuote{使一切有血气的，在天主面前一个也不能自夸}之后，为了指明人应在何事上夸耀，他立刻加上：\BibleQuote{但你们得以在基督耶稣内，是由于天主，他使我们成为从天主而来的智慧、正义、圣化与救赎：如经上记载：夸耀的，应因主而夸耀。}\BibleRef{格前 1:30}然后他的这个意图继续发展，直到后来责备他们时说：\BibleQuote{你们仍是属肉体的；因为在你们中间有嫉妒和纷争，你们岂不是属肉体，并按照人的方式行事吗？因为有人说我是属\Person{保禄}的，另一人说我是属\Person{阿颇罗}的，你们岂不是人吗？那么，\Person{阿颇罗}是什么？\Person{保禄}是什么？都是仆役，你们藉着他们而相信，且照主所赐予的每人。我栽种了，\Person{阿颇罗}浇灌了，但使之生长的是天主。所以，栽种的算不得什么，浇灌的也算不得什么，只有那使之生长的天主。}难道你们看不见，宗徒的唯一目的是要人谦卑，只高举天主吗？因为在那一切栽种和浇灌的事上，他说栽种的和浇灌的都算不得什么，只有使之生长的天主；而且，一个人栽种、另一个人浇灌这一事实，他也不归功于他们自己，而是归功于天主，当他说：\BibleQuote{照主所赐予的每人；我栽种了，\Person{阿颇罗}浇灌了。}因此，他坚持同一意图，进而说：\BibleQuote{所以，任何人都不可在人身上夸耀}\BibleRef{格前 3:21}，因为他已经说过：\BibleQuote{夸耀的，应因主而夸耀。}在此之后，连同其他一些相关的事，同一个意图继续在他以下的言语中：\BibleQuote{弟兄们，为了你们的缘故，我把这些事转用到我自己和\Person{阿颇罗}身上，叫你们从我们学习，不可超过圣经所记载的，以致为了一个而膨胀起来反对另一个。谁使你与众不同呢？你有什么不是领受的呢？既然领受了，为什么你夸耀，好像没有领受呢？}\BibleRef{格前 4:6}\par

% structure: p0022 source=h2 target=subsection reason=relative-heading-rank; base=h2

\subsection{第十章——是天主的恩宠特别区分一人与另一人。}

宗徒在此處的最明显意圖，是他反對人的驕傲，好叫無人因人的緣故而誇耀，只誇耀於天主。依我看，若認為他所指的便是天主的自然恩賜，不論是人在最初狀態中獲得的完整完善人性本身，或是墮落後人性所殘留的一切，這都太過荒謬。因為，難道是這類人人共有的恩賜，將人與人區分開來嗎？然而，他首先說：\Quote{誰使你與眾不同？}然後又說：\Quote{你有什麼不是領受的呢？}因為一個人若驕傲自大，輕視他人，可能說：\Quote{我的信德使我與眾不同}，或\Quote{我的義德}，或其他此類的話。為回應這類想法，這位良善的導師說：\Quote{你有什麼不是領受的呢？}這不是從那使你和他人不同、沒有賜給他人卻賜給了你的一位那裡領受的，又是從誰呢？他說：\Quote{既然你領受了，為什麼還要誇耀，好像沒有領受似的？}請問，他關心的，豈不是要讓那誇耀的，只誇耀於主嗎？但是，再沒有比一個人誇耀自己的功績，好像是他自己為自己賺得的，而非天主的恩寵（而且此恩寵使善人與惡人不同，並非善人與惡人所共有），更違背這種心態的了。因此，使我們成為有靈的生命體、與牲畜有別的恩寵，應歸於本性；同樣，在人中間，使俊美與醜陋有別，或聰明與愚笨有別，或諸如此類的恩寵，也應歸於本性。但宗徒所責備的人，並不因與牲畜相比，也不因與任何其他人在任何自然稟賦上（即使最惡劣的人也可能擁有）而自高自大。他卻將屬於聖善生活的一些美好恩賜歸於自己，而不歸於天主，並因此自高自大，理當聽到責備說：\Quote{誰使你與眾不同？你有什麼不是領受的呢？}因為，雖然擁有信德的能力是出於本性，但擁有信德本身也出於本性嗎？\Quote{因為並非眾人都擁有信德}（\BibleRef{得后 3:2}），雖然所有人都擁有信德的能力。但宗徒不是說：\Quote{你有什麼能力，這能力不是你領受的呢？}而是說：\Quote{你有什麼不是你領受的呢？}因此，擁有信德的能力，如同擁有愛德的能力一樣，屬於人的本性；但擁有信德，如同擁有愛德一樣，屬於信徒的恩寵。所以，那賦予我們擁有信德能力的本性，並不將人與人區分開來；而是信德本身使信徒與不信者不同。因此，當經上說：\Quote{誰使你與眾不同？你有什麼不是你領受的呢？}如果有人敢說：\Quote{我憑自己擁有信德，所以我並沒有領受它}，他便直接反駁了這最明顯的真理——這並非因為人的意志沒有選擇相信或不相信的自由，而是因為在蒙選者中，意志是由主預備的。如此，那經文\Quote{誰使你與眾不同？你有什麼不是你領受的呢？}正是指向那存在於人的意志中的信德本身。\par

% structure: p0024 source=h2 target=subsection reason=relative-heading-rank; base=h2

\subsection{第十一章【六】——有些人蒙選，乃是出於天主的仁慈。}

\Quote{許多人聽到了真理之言；但有些人相信，有些人反對。因此，前者願意相信，後者不願意。}誰不知道這事？誰能否認這事？但既然在某些人身上，意志是由主預備的，在另一些人身上卻沒有預備，我們必須能夠分辨什麼來自天主的仁慈，什麼來自祂的審判。\Quote{以色列所追求的}宗徒說，\Quote{他們沒有得到，但蒙選者卻得到了；其餘的都頑固了，正如經上所載：天主賜給了他們麻木的神——眼睛看不見，耳朵聽不見，直到今日。達味也說：願他們的筵席成為他們的網羅、報應和絆腳石；願他們的眼睛昏暗，以致看不見；願你時常彎曲他們的腰。\BibleRef{罗 11:7}}這裡有仁慈和審判：仁慈是對那得到了天主義德的蒙選者；審判是對其餘頑固的人。然而，前者因為願意相信而相信，後者因為不願意相信而不信。因此，在這些意志本身中，仁慈和審判已被顯明。無疑，這樣的揀選是出於恩寵，絕非出於功績。因為祂先前曾說：\Quote{所以，同樣，在現今的時期，也按照恩寵的揀選，留下了遺民。如果是出於恩寵，就不再是由於行為；不然，恩寵就不再是恩寵了。\BibleRef{罗 11:5}}因此，蒙選者白白地得到了他們所得的；沒有任何事情預先存在，使他們可以首先給予，然後才得到回報。祂白白地救了他們。至於其餘頑固的人，如同那裡清楚說明的，那是出於報應。\Quote{上主的一切途徑都是仁慈和真理。}但祂的道路是不可追尋的。因此，祂白白施予的仁慈和祂公正審判的真理，同樣是無可追尋的。\par

% structure: p0026 source=h2 target=subsection reason=relative-heading-rank; base=h2

\subsection{第十二章【七】——為何宗徒說我們是因信德成義，而非因行為。}

但或许可以说：\Quote{宗徒将信德与行为区分开来；他确实说恩宠不是由于行为，但他并没有说恩宠不是由于信德。}这的确是真的。但耶稣说信德本身也是天主的工作，并命令我们作这工作。因为犹太人对祂说：\BibleQuote{我们该作什么，才算作天主的工作呢？耶稣回答他们说：天主的工作，就是你们相信祂所派遣的那一位。}\BibleRef{若 6:28}因此，宗徒将信德与行为区分开来，就如希伯来两个王国中犹大与以色列被区分，尽管犹大本身就是以色列一样。他说人因信德成义，而非因行为，因为信德本身是首先赐下的，由此可以获得其他特别称为行为的事物，人可由此正直地生活。因为他自己也说：\BibleQuote{你们得救是由于恩宠，借着信德，而且这不是出于你们自己，而是天主的恩赐。}\BibleRef{弗 2:8}——这就是说：\Quote{而当他说\InnerQuote{借着信德}时，连信德本身也不是出于你们自己，而是天主的恩赐。}\Quote{不是出于行为，}他说，\Quote{免得有人自夸。}因为常有人说：\Quote{他配得相信，因为他在未信之前就是一个好人。}这事可用于\Person{科尔乃略}\EditorialNote{宗徒大事录·第十章}，因为他在相信基督之前，他的施舍已蒙悦纳，祈祷已蒙垂听；然而，没有某种信德，他既不会施舍，也不会祈祷。因为他怎能呼求他所不信的那一位呢？但如果他能没有对基督的信德而得救，\Person{宗徒伯多禄}就不会被派去作建筑者来建造他了；尽管：\Quote{若不是主建造房屋，建造的人就枉然劳力。}并且我们被告知，信德是出于我们自己，其他属于正义行为的事物是出于主；好像信德不属于建筑——好像，我说，基础不属于建筑一样。但如果信德首要且特别地属于建筑，那么凡试图通过宣讲来建造信德的人，除非主在祂的仁慈中从内里建造它，否则就是枉然劳力。因此，\Person{科尔乃略}所行的一切善工，无论是他信基督之前，还是信时或信后，都应归于天主，免得有人自夸。\par

% structure: p0028 source=h2 target=subsection reason=relative-heading-rank; base=h2

\subsection{第十三章 [VIII.]——天主恩宠的效力。}

因此，我们唯一的主和师傅自身，在祂说了我上面提到的话——\Quote{天主的工作，就是你们相信祂所派遣的那一位。}——之后不久，在同一篇讲话中说：\BibleQuote{我对你们说过：你们看见了我，仍然不信。凡父赐给我的人，必到我这里来。}\BibleRef{若 6:36}\Quote{必到这里来}是什么意思，不就是\Quote{必信我}吗？但这是父的恩赐才能如此。此外，过了一会儿祂说：\Quote{你们不要彼此窃窃私议。除非派遣我的父吸引他，没有人能到我这里来；在末日我要使他复活。先知书上记载：\InnerQuote{他们都要蒙天主训诲。}凡听了父所教导而学习的，必到我这里来。}\Quote{凡听了父所教导而学习的，必到我这里来}是什么意思？除非是说，凡听了父而学习的，没有不到我这里来的？因为如果每一个听了父而学习的都来，那么每一个不来的肯定没有听父的；因为如果他有听和有学习，他就会来。因为没有人听了和学习而不来；但每一个人，如真理所宣告的，凡听了父而学习的，都来。这种教导，其中父被听见并教导人来归向圣子，远远超越肉体的感觉。圣子本身也参与其中，因为祂是祂的圣言，借此祂如此教导；祂这样做不是通过肉体的耳朵，而是通过心。同时，圣父与圣子的圣神也参与其中；祂也教导，并且不是分开教导，因为我们知道\OriginalTerm{天主圣三}{Trinity}的工作是不可分割的。那确实是同一圣神，宗徒说：\BibleQuote{但我们既然具有同样信德的神。}\BibleRef{格后 4:13}但这尤其归因于圣父，因为独生子是由祂所生，圣神由祂所发，对此再详细论述就繁琐了；我想我已经写了一部十五卷的关于天主圣三的著作，应该已经送到你手中了。我说，这种教导，其中天主被听见并教导，远远超越肉体的感觉。我们看到许多人来到圣子面前，因为我们看到许多人相信基督；但他们在何时何地听到了父的教导并学习了，我们却看不见。的确，那恩宠极其隐秘，但谁怀疑它是恩宠呢？因此，这恩宠，由天主的恩赐隐秘地赋予人心，没有一个铁石心肠会拒绝，因为它正是为了先除去心的铁石。因此，当父在里面被听见并教导，使人来到圣子面前时，祂除去石心，赐给肉心，正如先知所宣告的祂所应许的。因为祂如此使他们成为祂所预备的荣耀的慈悲之子和容器。\par

% structure: p0030 source=h2 target=subsection reason=relative-heading-rank; base=h2

\subsection{第十四章——为什么父不教导所有的人，使他们来归向基督。}

那么，为什么祂不教导所有的人，使他们来到基督面前呢？岂不是因为凡祂所教导的人，祂是怀着慈悲教导；凡祂不教导的人，祂是出于审判而不教导？正如经上记载：\BibleQuote{祂愿意怜悯谁，就怜悯谁；祂愿意使谁硬心，就使谁硬心。}\BibleRef{罗 9:18} 但祂施恩时是怜悯，祂报应应得的事时是硬心。或者，若有人愿意这样区分，那些话也是那位宗徒对其说\Quote{你必对我说}之人的话，这样他可被视为说了\Quote{因此，祂愿意怜悯谁，就怜悯谁；祂愿意使谁硬心，就使谁硬心}，以及随后的内容——即\Quote{这样还有什么可抱怨的呢？谁能抗拒祂的旨意呢？}——宗徒回答\Quote{人啊，你说的是假话吗？}不，他反而说：\Quote{人啊，你是谁，竟敢向天主强辩？受造之物岂能对造它的说：你为什么这样造了我？窑匠难道没有权柄用同一团泥…}以及随后你所熟知的那些话。然而在某种意义上，圣父教导所有的人来到祂的圣子面前。因为先知书上所写的话不是徒然的：\BibleQuote{他们都要蒙天主教导。}\BibleRef{若 6:45} 当祂也先引用了这见证后，祂接着说道：\BibleQuote{因此，凡听了父的教导而学习的，都到我这里来。}因此，正如我们公正地说到某城唯一的语文教师：他向所有人教授语文——并非所有人都学，而是凡在那里学语文的，没有不从他而学的——同样，我们公正地说：天主教导所有人来到基督面前，不是因为所有的人都来，而是因为没有人以其它方式来到。至于祂为什么不教导所有的人，宗徒按他认为可以解释的方式解释了，因为：\BibleQuote{祂愿意显示祂的义怒，并彰显祂的大能，就多多忍耐了那些可怒、预备遭毁灭的器皿；并且要将祂荣耀的丰富，彰显在那蒙怜悯、早预备得荣耀的器皿上。}\BibleRef{罗 9:22} 因此，\BibleQuote{十字架的道理，在灭亡的人看来是愚拙；在我们得救的人看来，却是天主的德能。}\BibleRef{格前 1:18} 天主教导所有这样的人来到基督面前，因为祂愿意所有这样的人得救，并认识真理。如果祂曾愿意教导那些视十字架的道理为愚拙的人来到基督面前，毫无疑问，这些人也会来。因为当祂说\BibleQuote{凡听了父的教导而学习的，都到我这里来}时，祂既不欺骗也不受骗。那么，千万别认为有谁听了父的教导而学习却不来的。\par

% structure: p0032 source=h2 target=subsection reason=relative-heading-rank; base=h2

\subsection{第十五章　凡相信的人都是由天主所教导的。}

\Quote{那么，}他们说，\Quote{祂为什么不教导所有的人？}若我们说，那些祂不教导的人是不愿学习，我们就会遇到这样的回答：那么，向祂所说的话\Quote{天主啊，求祢使我们回转，并使我们复活}又算什么呢？如果天主不使那些不愿的人成为愿意，那么教会根据主的命令为迫害她的人祈祷，又是什么原则呢？因为真福西彼连也这样理解我们所念的\Quote{愿祢的旨意奉行在人间，如同在天上}——即如同在那些已经相信的人身上，他们好比是\Emph{天}，同样也在那些不信的人身上，他们因此仍是\Emph{地}。那么，我们为那些不愿相信的人祈祷，岂不是祈求天主也在他们身上工作，使他们也愿意吗？当然，宗徒说：\BibleQuote{弟兄们，我心里所愿的，向天主所求的，是要他们得救。}\BibleRef{罗 10:1} 他为那些不信的人祈祷——为的是什么呢？岂不就是为了使他们相信？因为除此以外，他们没有别的得救方式。那么，如果祈求者的信德先于天主的恩宠，那么为他们祈求得以相信的人的信德难道也先于天主的恩宠？——因为为他们所求的正是这一点：使那些不信的人——即没有信德的人——能被赐予信德本身。因此，当福音被宣讲时，有些人相信，有些人不相信；但那些因外在宣讲者的声音而相信的人，从内在听到了父的教导并学习了；而那些不相信的人，外在听到了，但内在既没有听见也没有学习——也就是说，前者被赐予了相信，后者没有被赐予。因为祂说：\BibleQuote{没有人能到我这里来，除非派遣我的父吸引他。}\BibleRef{若 6:44} 随后这话说得更明白。因为过了一会儿，当祂谈到吃祂的肉、喝祂的血时，连祂的一些门徒也说：\Quote{这话甚难，谁能听呢？}耶稣心里知道门徒为这话议论，就对他们说：\Quote{这话使你们厌烦吗？}稍后祂又说：\Quote{我对你们所说的话就是神，就是生命；只是你们中间有不信的人。}接着福音作者立即说：\Quote{因为耶稣从起头就知道谁是不信的人，谁要出卖祂；所以祂说：因此我对你们说过，若不是蒙我父的恩赐，没有人能到我这里来。}因此，被父吸引到基督面前，以及为了来到基督面前而听父的教导并学习，无非就是从父领受相信基督的恩赐。因为那位说\BibleQuote{若不是蒙我父的恩赐，没有人能到我这里来}的，并不是区别听福音的人与未听福音的人，而是区别相信的人与不相信的人。\par

% structure: p0034 source=h2 target=subsection reason=relative-heading-rank; base=h2

\subsection{第十六章　为何信德的恩赐没有赐给所有人。}

信德，无论是在其开端还是在其完成，都是天主的恩宠；任何人都不可有任何怀疑，除非他愿意抵抗最清晰的圣经，即这恩宠赐给了某些人，而另一些人则未蒙赐予。但为何不赐予所有人，不应困扰信徒，因为他相信从一人众人都进入了那无疑是极其公义的定罪之中；因此，即使没有人从中获得拯救，也没有正当理由指责天主。由此显而易见，许多人获得拯救乃是莫大的恩宠，并在那些未蒙拯救的人身上认识到自己本应承受什么；这样，那夸耀的，应当夸耀于主，而非夸耀于自己的功绩，因为他看到这些功绩在那些被定罪的人中也同样存在。但为何祂拯救一人而非另一人——\BibleQuote{祂的判断不可猜测，祂的踪迹不可追寻}\BibleRef{罗 11:33} 在这种情况下，我们最好听见或说：\BibleQuote{人哪，你是谁，竟敢向天主抗辩？}\BibleRef{罗 9:20} 而不应胆敢开口，仿佛我们能知道祂选择保密的事。此外，祂不可能愿意任何不义的事。\par

% structure: p0036 source=h2 target=subsection reason=relative-heading-rank; base=h2

\subsection{第十七章〔IX.〕——他在反对\Person{波菲利}的信函中的论证，论福音为何如此晚才来到世界。}

但你在我的某篇反驳\Person{波菲利}的小论著中记得我曾说过的话，题为\WorkTitle{基督宗教的时代}，我那样说是为了避开关于恩宠的这个更为仔细和精妙的论证；虽然其意义（可以在别处或由他人阐述）并未完全略去，尽管我不愿意在那个地方解释它。因为，在其他问题中，我这样回答所提出的问题，即为什么基督是在如此长久的时间之后才来的：\Quote{因此，我说，既然他们并不反对基督，认为所有人都没有跟随祂的教导（因为他们自己也觉得这一点完全不能公正地反对哲学家的智慧，甚至反对他们自己神祇的神性），那么他们会如何回答呢？如果我们——暂不考虑天主智慧和知识的深奥，在那里或许某种更隐秘地隐藏的另一神圣计划，也不预判其他原因，这些是智者无法探究的——只为了在论证此问题时保持简洁，对他们说这一点：基督愿意向人显现，并且祂的教义应当在人们中宣讲，是在祂知晓的时间、在祂知晓的地点，那里有些将会信祂的人。因为在那些时间、那些地点，祂的福音没有被宣讲，祂预知在祂的宣讲中，所有人都会像——并非所有人，而是许多人——在祂肉身临在时那样，即使祂使死人复活，他们仍不会信祂；正如我们现在看到许多人，虽然先知关于祂的宣告借着这些显现得以应验，但他们仍然不愿相信，宁愿以人间的狡诈来抵抗，而不屈服于如此清晰易懂、如此崇高、如此卓越显明的神圣权威，只要人的理解在接近神圣真理时仍渺小软弱。那么，如果基督知道古时的世界如此充满不信者，以至于祂理所当然地拒绝向他们显现或被宣讲，因为他们（祂预知）既不信祂的话也不信祂的奇迹，这有什么奇怪呢？因为并非难以置信，那时所有人都像从祂来临直到现在，我们惊叹有那么多人曾经是、现在也是这样的人。然而从人类之初，有时更隐秘，有时更显明，甚至就天主眷顾而言，时机似乎合适，预言从未中断，信祂的人也不缺乏；从亚当到梅瑟，在以色列民族本身（它因某种特殊奥秘而成为先知性的人民）中，以及在祂肉身来临之前的其他民族中。因为在神圣的希伯来经卷中提到，早在亚伯拉罕时代，有些人——既不属于他的血统，也不属于以色列民族，也不属于以色列人中的外邦团体——却分享了他们的圣事，我们为何不能相信在其他民族中也有这样的人，各处皆有，尽管我们在同样的权威著作中没有读到任何提及？因此，这宗教的救恩（唯有通过它才真正应许了唯一真实的救恩）从未辜负配得的人；而无论谁辜负了它，都是不配的。并且从人类繁衍之初，直到终结，福音被宣讲，对一些人是为了奖赏，对另一些人是为了审判；因此，那些从未听到信德宣讲的人，被预知为不会相信的人；而那些听到宣讲的人，即使他们并非会相信的人，也被树立为前者的榜样；而那些听到宣讲并将会相信的人，则为天国和圣天使的团契而预备。}\par

% structure: p0038 source=h2 target=subsection reason=relative-heading-rank; base=h2

\subsection{第十八章——前述论证应用于当下。}

岂不见我的愿望是，在不预先判断天主奥秘的旨意和其他理由的情况下，说出关于基督预知似乎足以说服提出此问题的外邦人的不信吗？因为有什么比以下更真实：基督预知谁将信祂，以及他们在何时何地信祂？但那些人是否因向他们自己宣讲基督而拥有信德，还是他们将藉天主的恩赐领受信德——即，天主是否仅预知他们，还是也预定了他们——我当时认为不必询问或讨论。因此，我当时所说的：\Quote{基督愿意在那个时代向人显现，并且祂的教训当在祂知道和祂知道那里有将要信祂的人中间被宣讲}，也可以这样说：\Quote{基督愿意在那个时代向人显现，并且祂的福音当在祂知道和祂知道那里有在创世以前在祂内被拣选的人中间被宣讲}。但是，既然若是这样说了，会使读者想要询问那些现在因白拉奇错误的警告而必须更丰盛、更仔细地讨论的事，那么在我看来，当时所足够的事应当简要地说出，如同我所说的，把天主智慧与知识的深奥搁置一旁，并且不预先判断其他理由，关于这些理由，我认为我们更适合在另一时候，而不是当时，进行辩论。\par

% structure: p0040 source=h2 target=subsection reason=relative-heading-rank; base=h2

\subsection{第十九章［X］——预定与恩宠的区别何在。}

此外，我所说的话：\Quote{这个宗教的救恩从未缺乏给那堪当的人，而那缺乏的人是不堪当的}——如果讨论并追问任何人如何能堪当，不乏有人说——凭人的意志。但我说，凭天主的恩宠或预定。此外，恩宠与预定之间仅有这一区别：预定是恩宠的准备，而恩宠是这赠与本身。因此，当宗徒说：\BibleQuote{不是出于行为，免得有人自夸。因为我们原是祂的化工，在基督耶稣内受造，为行善事}，\BibleRef{弗 2:9}这是恩宠；但随后的话——\BibleQuote{就是天主所预备叫我们行在其中}——是预定，这不能没有预知而存在，尽管预知可以没有预定而存在；因为天主藉预定预知祂将要作的事，因此经上说：\BibleQuote{祂造了将要成的事。}\BibleRef{依 45:11}此外，祂也能预知祂自己不亲自作的事——例如一切罪。因为，尽管有些罪同时也是罪的惩罚，经上说：\BibleQuote{天主就任凭他们存邪僻的心，行那些不合理的事，}\BibleRef{罗 1:28}在这种情况下，罪不是天主的，而是审判。因此，天主对善的预定，如前所述，是恩宠的准备；这恩宠就是那预定的效果。因此，当天主向亚巴郎应许他后裔中的万民信德，说：\BibleQuote{我已立你为万民之父，}\BibleRef{创 17:5}宗徒据此说：\BibleQuote{所以，由于信德，为的是恩宠，使应许坚定不移地归于一切后裔，}\BibleRef{罗 4:16}祂应许的不是出于我们意志的能力，而是出于祂自己的预定。因为祂应许了祂自己要作的事，而不是人要作的事。因为，虽然人行那些属于敬拜天主的善事，但祂亲自使他们作祂所命令的；不是他们促使祂作祂所应许的。否则，天主应许的实现就不在天主的能力内，而在人的能力内了；这样，天主向亚巴郎所应许的，就会由人自己赐给亚巴郎。然而，亚巴郎并不这样信，而是\BibleQuote{他怀着信德，将光荣归于天主，且深信祂所应许的，也必能成就。}\BibleRef{罗 4:21}他没有说\Quote{预言}——他没有说\Quote{预知}；因为祂也能预言并预知别人的行为；而是说\Quote{祂也必能成就}；这样，他所说的不是别人的行为，而是祂自己的。\par

% structure: p0042 source=h2 target=subsection reason=relative-heading-rank; base=h2

\subsection{第二十章——天主是向亚巴郎应许了万民的善工，而不是他们的信德吗？}

难道天主向亚巴郎应许他后裔中的万民善工，以致应许那祂自己所作的事，而并未应许外邦人的信德——那是人自己作的；但为了应许祂自己所作的事，祂预知人会成就那信德？宗徒并非如此说，因为天主应许了亚巴郎儿女，他们要跟随他的信德的足迹，正如祂非常清楚地说的。但是，如果祂应许了善工，而不是外邦人的信德——当然，因为善工除非从信德来就不是善工（因为\BibleQuote{义人因信德而生活，}\BibleRef{哈 2:4}并且\BibleQuote{凡不出于信德的都是罪，}\BibleRef{罗 14:23}又\BibleQuote{没有信德，是不可能中悦天主的，}\BibleRef{希 11:6}），然而，天主成就祂所应许的事仍在于人的能力。因为除非人作那没有天主恩赐而属于人的事——即，除非人自己有信德——他就不会使天主赐予；这样，天主所应许的——即由天主赐予正义的善工——就不实现。因此，天主能够成就祂的应许，就不在天主的能力内，而在人的能力内。如果真理与虔诚不禁止我们相信这一点，那么让我们与亚巴郎一同相信，祂所应许的，也必能成就。但祂向亚巴郎应许了儿女；这些人除非有信德就不能存在，因此祂也赐予信德。\par

% structure: p0044 source=h2 target=subsection reason=relative-heading-rank; base=h2

\subsection{第二十一章——人竟宁愿信赖自己的软弱而胜过天主的刚强，真是可惊奇。}

诚然，宗徒说：\BibleQuote{所以是人出于信德，为使恩许有保证，按照恩宠，}\BibleRef{罗 4:16}，我惊异人们宁愿托付于自己的软弱，而不托付于天主应许的坚强。但你说：天主对我自己的旨意，对我是不确定的吗？那么，你对你自己的旨意是确定的吗？你难道不害怕—\BibleQuote{凡自以为站得稳的，务要小心，免得跌倒}\BibleRef{格前 10:12}？既然如此，两者都不确定，人为何不将自己的信德、望德和爱德托付给更强的意志，却托付给更弱的呢？\par

% structure: p0046 source=h2 target=subsection reason=relative-heading-rank; base=h2

\subsection{第22章——天主的应许是确实的。}

\Quote{但是，}他们说，\Quote{当说\InnerQuote{你若相信，便得救}时，这两者中一个是要求的，另一个是赐予的。所要求的是在人能力范围内的，所赐予的则是在天主内的。}为什么这不都属于天主呢？无论祂所命令的或祂所赐予的？因为祂被祈求赐予祂所命令的。信徒祈求他们的信德得以增长；他们替那些不信的人祈求，为使他们得到信德；因此，无论是在其增长中还是在其开端，信德都是天主的恩赐。但这样说：\Quote{你若相信，便得救}，就像说\BibleQuote{若依赖圣神，你们去克制肉身的行为，必得生活}\BibleRef{罗 8:13}一样，因为在这段经文里，这两件事中一个是要求的，另一个是赐予的。经文说：\BibleQuote{若依赖圣神，你们去克制肉身的行为，必得生活}。因此，克制肉身的行为是被要求的，而得以生活是赐予的。那么，我们能否说，克制肉身的行为不是天主的恩赐，并因我们听到它被要求，而且若实行便获得生活作为赏报，就不承认它是天主的恩赐？这决不能被恩宠的分享者与维护者所赞同！这是白拉奇派的可责的错误，宗徒立刻堵住了他们的口，接着加上：\BibleQuote{因为凡受天主圣神引导的，都是天主子。}\BibleRef{罗 8:14}，免得我们相信克制肉身的行为是出于我们自己的，而不是出于天主圣神。而且，关于这圣神，他在另一处说：\BibleQuote{这一切无非是同一圣神所行，随祂的心愿分给各人。}\BibleRef{格前 12:11}。在这些恩赐中，如你所知，他也提到了信德。因此，虽然克制肉身的行为是天主的恩赐，却仍被要求于我们，且生命作为赏报摆在我们面前；同样，信德也是天主的恩赐，尽管当说\Quote{你若相信，便得救}时，信德被要求于我们，而救恩作为赏报被提供给我们。因为这一切既被命令给我们，又被显示为天主的恩赐，为使我们明白，我们既行这些事，又是由天主使我们行事，正如祂通过厄则克耳先知最清楚地说的。有什么比祂说\BibleQuote{我必要使你们遵行}\BibleRef{则 36:27}更清楚的呢？请留意那段圣经，你将会看到天主应许祂将使他们遵行祂所命令的事。祂实在没有沉默不提那些人的功过，而是提他们的恶行，祂表明自己是以善报恶，正因祂此后使他们有善行，使他们遵行天主的诫命。\par

% structure: p0048 source=h2 target=subsection reason=relative-heading-rank; base=h2

\subsection{第23章 [XII.]——恩宠与预定在婴孩身上及在基督身上的卓越例证。}

但所有这些论证，我们藉以主张我们主耶稣基督的天主恩宠是真实的恩宠，即不是按我们的功过赐予的，虽然它由神圣宣言的见证最明白地宣告，但在那些认为除非他们将某些东西归功于自己，即他们先给予好让回报给他们，否则他们就被剥夺了所有虔敬热忱的人中，对于已行使意志选择的成年人而言，确实有些困难。但当我们来到婴孩的情况，以及天主与人之间的中保，那人基督耶稣时，所有关于人的功德先于天主恩宠的主张都付之阙如，因为前者并非因任何先行的善功而分别出来，使他们属于人类的救主；同样，祂自己作为人，也并非因任何先行的属人功德而成为人类的救主。\par

% structure: p0050 source=h2 target=subsection reason=relative-heading-rank; base=h2

\subsection{第24章——没有人按他若活得更久会做的事来受审判。}

因为谁能听信说，尚在襁褓中领受圣洗的婴儿，是因他们未来的功绩而离开此世，而其他未领洗的婴儿，在同一年龄去世，是因他们的未来功绩已被预知——却是邪恶的；以致天主赏报或惩罚他们，并非根据他们或善或恶的生活，而是根本毫无生活？宗徒确实设下了一个限度，人的轻率猜疑（说得温和些）不应逾越，因为他说：\Quote{我们众人都要站在基督的审判台前，好使每人按他藉身体所作的事，或善或恶，领取报应。} \BibleRef{格后 5:10} 他说\Quote{所作的}，并未加上\Quote{或将要作的}。但我不知这种思想如何进入这些人的头脑，认为婴儿的未来功绩（将不会存在）应受惩罚或尊荣。但为何说人将按他藉身体所作的事受审判，尽管许多事是单凭心灵所作，而非藉身体或任何肢体；且多半是极其重要的事，以致如此最公义的惩罚应归于这种思想，例如——姑且不提其他——\Quote{愚妄的人心里说：没有天主}？那么，\Quote{按他藉身体所作的事}是什么意思呢？除非是指他在身体内期间所作的事，这样我们可以理解\Quote{藉身体}就是\Quote{在整个肉身生命期间}。但身体之后，没有人再会在身体内，除非在最后的复活——不是为了建立任何功绩，而是为了领受善功的报赏，忍受恶罚。但在脱去身体与再取身体之间的这段中间时期，灵魂要么受折磨，要么得安息，根据他们在肉身生命期间所作的事。此外，白拉奇派否认、而基督的教会承认的\OriginalTerm{原罪}{original sin}，也属于肉身生命期间；根据这原罪是被天主的恩宠解除，还是被天主的审判不解除，当婴儿去世时，他们或藉\OriginalTerm{重生}{regeneration}的功绩从恶迁善，或藉本源的功绩从恶迁恶。公教信仰承认这一点，甚至某些异端也毫无异议地同意。但我深感惊奇和诧异，无法发现那些才智（你的书信显示决非平庸）的人，何以会持有这样的观点：人应受审判，不是根据他仍在身体内时所拥有的功绩，而是根据他若在身体内活得更久就会拥有的功绩；若非我敢于不信你，我是不敢相信有这样的人的。但我希望天主会干预，使他们一经提醒就立刻明白，如果那些据说是将会有的罪，能藉天主的审判正当地惩罚未领洗的人，那么它们也能藉天主的恩惠正当地赦免领洗的人。因为凡说未来的罪只能藉天主的审判惩罚，却不能藉天主的慈悲赦免的人，应当思考他对天主及其恩宠造成了多大的伤害；好像未来的罪能被预知，却不能不被追溯。如果这是荒谬的，那么更有理由当那些若活得更久就会成为罪人的婴儿在幼年去世时，藉那能洗除罪恶的\OriginalTerm{圣洗}{laver}，给予他们帮助。\par

% structure: p0052 source=h2 target=subsection reason=relative-heading-rank; base=h2

\subsection{第二十五章 [XIII.]— 或许领受圣洗的婴儿若存活会悔改，而未领洗的则不会。}

但是，或许他们会说，罪过被重新赦免给忏悔的人，而那些在婴儿期死去的人之所以没有受洗，是因为天主预知他们如果活下来也不会忏悔；而天主预知那些受洗并在婴儿期死去的人如果活下来则会忏悔。那么，请他们观察并看到：如果是这样，那么在此情况下，受惩罚的并非在未受洗而死去的婴儿身上的原罪，而是如果他们活下来每人可能犯的罪；同样，在受洗的婴儿身上，被洗净的也不是原罪，而是他们如果活下来可能犯的未来之罪，因为他们不可能在更成熟的年龄之前犯罪；但有些人被预知为将要忏悔，而另一些人被预知为不会忏悔，因此有些人受了洗，另一些人则未受洗就离开了今生。如果佩拉基乌派胆敢这样说，那么通过否认原罪，他们就可以免除为那些在天国之外的婴儿寻找某种我不知道的幸福之地的必要——尤其是因为他们确信，这些婴儿不能拥有永生，因为他们没有吃基督的肉，也没有喝基督的血；而且，在他们这些根本没有罪过的人身上，为赦免罪过而设立的圣洗圣事就成了虚假的。因为他们会继续说，没有原罪，但那些作为婴儿离世的人，要么受洗，要么未受洗，是根据他们如果活下来将来会有的功过；并且根据他们将来会有的功过，他们要么领受，要么不领受基督的身体和血，而没有这些，他们绝对无法获得生命；而且，他们受洗是为了真正赦免罪过，尽管他们从亚当那里没有继承任何罪过，因为那些罪过——天主预知他们会忏悔——被赦免给了他们。这样，他们就可以极其轻松地辩护并赢得他们的论点，即否认有任何原罪，并坚持天主的恩宠只按照我们的功过赐予。但是，人的未来功过——这些功过绝不会成为现实——无疑是毫无功过可言的，这当然是最容易看出的：因此，连佩拉基乌派自己也无法这样说；而这些人更不应该说。因为我能以何等痛苦说出：我发现那些与我们一同在天主教权威下谴责那些异端者错误的人，竟然没有看到这一点，而佩拉基乌派自己却已经看到这是极其虚假和荒谬的。\par

% structure: p0054 source=h2 target=subsection reason=relative-heading-rank; base=h2

\subsection{第二十六章 [XIV]—— 引用\Person{西彼廉}的著作\WorkTitle{论死亡}。}

\Person{西彼廉}写了一部著作\WorkTitle{论死亡}，被许多几乎全部热爱教会文学的人熟知并认可，其中他说死亡不仅对信徒无害，甚至被认为是有益的，因为它使人脱离犯罪的风险，并使他们进入不再犯罪的稳妥之中。但是，如果连尚未犯下的未来之罪也要受惩罚，那么这又有什么益处呢？然而，他极为充分且出色地论证了：犯罪的风险在今生并不缺乏，而在今生结束后也不再继续；他在其中还引用了\BibleBook{}{智慧篇}中的那段证言：\BibleQuote{他被接去，免得邪恶改变他的悟性。}\BibleRef{智 4:11} 我也引用了这同一段话，虽然你曾说，你的那些弟兄们以它并非出自正典经卷为由拒绝了它；仿佛即使撇开这部书的见证，我所希望从中得到教导的事情本身也不清楚似的。因为哪个基督徒敢否认：义人若被死亡过早地抓住，必得安息？无论谁这样说，哪个信仰纯正的人会认为自己可以反对呢？此外，如果他说：义人若从他长久生活其中的正义中堕落，并在那邪恶中——我且不说一年，而是一天——生活之后死去，将进入为恶人所准备的惩罚，他的正义在未来将毫无效力——任何信徒会反驳这显而易见的真理吗？再者，如果问我们：他若在还是义人的时候死去，是遭受惩罚还是得享安息？我们会犹豫回答\Quote{安息}吗？这正是为何那句话——无论谁说——\BibleQuote{他被接去，免得邪恶改变他的悟性}。因为这话是就今生的风险而言的，而不是就天主的预知而言——天主预知将要发生的事，而非不会发生的事——也就是说，天主赏赐他早逝，为使他脱离诱惑的不确定性；并非因为他会犯罪，因为他本不会留在诱惑中。因为关于今生，我们在\BibleBook{}{约伯传}中读到：\BibleQuote{人的生命在世上不是一种考验吗？}\BibleRef{约 7:1} 但是，为何有些人蒙恩在尚为义人之际就被接离今生的危险，而另一些直到从正义堕落之前都保持义人的人，却在更长的生命中留在同样的危险中——谁能知道上主的心意？然而，由此可以明白：即使是那些持守善良虔诚品格直至成熟的老年和今生最后一天的义人，也不该夸耀自己的功过，而应夸耀主自己——因那位将义人从短暂生命中接去以免邪恶改变他悟性的天主，也亲自守护着义人在任何长久生命中，使邪恶不能改变他的悟性。但为什么天主却把义人留在此地以致堕落，而本来可以在此之前将他接走？祂的审判虽然绝对公义，却是不可探究的。\par

% structure: p0056 source=h2 target=subsection reason=relative-heading-rank; base=h2

\subsection{第二十七章——\BibleBook{}{智慧篇}在教会中获得正典圣经的权威。}

既然这些事如此，就不应拒绝智慧篇的论断，因为许多年来，这部书在基督的教会中，从读经员的席位被诵读，并由所有基督徒——从主教直到最卑微的平信徒、补赎者和慕道者——怀着对神圣权威的敬意聆听。因为确实，如果我从那些在我之前注释圣经的人中，提出为这个论断辩护——我们现在比平时更仔细、更详尽地被召来为它辩护，以对抗白拉奇主义者的新谬误——也就是说，天主的恩宠不是按照我们的功绩赐予的，而是白施予它所赐予的人，因为\BibleQuote{不在乎愿意的，也不在乎奔跑的，只在乎显示仁慈的天主}；但凭公义的审判，它不赐予它所不赐予的人，因为在天主没有不义——因此，如果我将我们之前的天主教注释者对神圣训示的辩护提出来，那么这些我现在为之发言的弟兄们肯定会同意，因为你们在信中暗示了这一点。那么，我们有何必要去查阅那些在此异端兴起之前无需深究如此难解问题之人的著作呢？毫无疑问，如果他们被迫回答此类问题，他们一定会这样做。因此，他们在著作中只是偶尔简短地提及他们对天主恩宠的看法；但他们在反驳教会敌人以及劝勉各种美德以事奉永生的真天主、获得永生和真福的事上，却详细论述。但是，天主的恩宠的能力，通过祈祷中频繁的提及而质朴地显明出来；因为天主命令应做的事，若非能由祂赐予使之成就，就不会向天主祈求。\par

% structure: p0058 source=h2 target=subsection reason=relative-heading-rank; base=h2

\subsection{第28章。—— \Person{西彼廉}的论著\WorkTitle{论死亡}。}

但如果有人希望了解那些处理过此主题之人的意见，他应当将\BibleBook{}{智慧篇}置于所有注释者之上，其中读到：\BibleQuote{他被夺去，免得邪恶改变他的悟性}；因为杰出的注释者，甚至在宗徒时代最接近的时期，也将其置于自身之上，因为他们使用时，相信自己使用的是神圣的证言；确实，最蒙福的西彼廉为了推崇早逝的好处，主张那些结束此生（其中可能犯罪）的人被夺去，免于罪的风险。在同一篇论著中，他说道：\Quote{当你们即将与基督同在，并确知天主的许诺时，为何不拥抱被召到基督那里，并因脱离魔鬼而欢喜呢？}在另一处他说：\Quote{孩童逃脱了他们不稳定年龄的危险。}又在另一处他说：\Quote{我们为何不赶快奔跑，好能看见我们的家乡，向我们的亲人致意？许多我们亲爱的人在那里等着我们——一大群密集的父母、兄弟、儿子渴望我们；他们已确知自己的安全，但仍担忧我们的得救。}通过这些和类似的情感，那位导师在天主教信仰最清晰的光照下充分而明确地见证：罪恶的危险和试探甚至直到脱去这身体都应当畏惧，但此后没有人会再遭受任何这类事情。即使他没有这样见证，何时有任何基督徒会对此事有疑问呢？那么，对于一个已经失足、并在同一失足状态中悲惨结束生命、并进入这样之人当受的惩罚的人——我说，这样的人若在跌倒之前被死亡从这诱惑领域夺去，怎能不是最大和最高的益处呢？\par

% structure: p0060 source=h2 target=subsection reason=relative-heading-rank; base=h2

\subsection{第29章。——天主的作为不依赖于人任何偶然的功绩。}

因此，除非我们沉溺于轻率的争论，关于那被夺去免得邪恶改变其悟性之人的整个问题就解决了。而\BibleBook{}{智慧篇}，如此多年来在基督的教会中被诵读，其中读到这句，不应遭受不公，因为它抵挡那些为人的功绩辩护而反对天主最明显恩宠的人；这恩宠主要显现在婴儿身上，当一些受洗、一些未受洗的婴儿结束此生时，他们充分显明天主的仁慈和祂的审判——祂的仁慈确实是白施的，祂的审判是应受的。因为如果人应按他们生活的功绩受审判——这些功绩他们因死亡而未能实际拥有，但若他们活着本会拥有——那么对被夺去免得邪恶改变其悟性的人就没有益处；对在失足状态中死去的人，若他们早死，也没有益处。这一点没有基督徒敢说。因此，我们的弟兄们，他们同我们一起为天主教信仰攻击白拉奇谬误的灾害，不应如此偏袒那认为天主的恩宠是按我们的功绩赐予的白拉奇意见，以至于企图（这是他们不敢的）否定那真实、明显且自古以来的基督徒情感——\BibleQuote{他被夺去，免得邪恶改变他的悟性}；并建立那种我们认为——我不说无人相信，而是无人梦想——即任何死者将根据他若活更长时期会做的事受审判。诚然，这样我们所说的就显明是无可争辩的：天主的恩宠不是按照我们的功绩赐予的；因此，那些反驳这一真理的聪明人被迫说出必须被所有人从耳中和思想中拒绝的话。\par

% structure: p0062 source=h2 target=subsection reason=relative-heading-rank; base=h2

\subsection{第30章 [XV]——预定的最光辉榜样是基督耶稣。}

此外，预定的最光辉明灯和恩宠的明灯就是救主本身——天主与人之间的中保本身，那人基督耶稣。请问，祂内的人性——无论就行为或信德而言——是凭自己什么先行的功绩，才得以成为这样的呢？我恳求，让这有一个回答。那个人，祂从哪里配得这事——被与圣父同永恒的圣言所接纳，进入位格的合一，并成为天主的独生子？是因为祂内有什么良善先行吗？祂先前做了什么？祂信了什么？祂求了什么，以致获得这不可言喻的卓越？岂不是因圣言的行动和接纳，那人从祂开始存在的那一刻起，就开始成为天主的独生子吗？难道那位充满恩宠的女子，不是怀孕了天主的独生子吗？祂不是由圣神和童贞玛利亚所生，成为天主的独生子——不是出于肉欲，而是出于天主的特殊恩赐吗？难道要担心随着年龄增长，这人会因自由意志而犯罪吗？或者说，祂内的意志因此而不自由吗？岂不是按祂更不可能成为罪的奴仆的程度，意志也更为自由吗？的确，在祂内，人性——也就是说，我们的人性——特别领受了所有那些特别可钦羡的恩赐，以及任何其他最可真实地说是祂所独有的恩赐，而这些并非凭其任何先行的功绩。让人在这里如果敢的话，回答天主说：为什么我不是祂？如果他听到：\Quote{人啊，你是谁，竟敢反驳天主？}\BibleRef{罗 9:10}，愿他此时不要克制自己，反而更加放肆地说：\Quote{我怎会听到\InnerQuote{人啊，你是谁？}既然我就是我所听到的——也就是说，一个人，而我所谈论的祂也不过是同样的人？为什么我不也是祂所是的呢？因为祂之所以如此伟大，是因恩宠；既然本性相同，恩宠为何不同？的确，天主不以貌取人。}我说，这不是什么基督徒，而是什么狂人会这样说呢？\par

% structure: p0064 source=h2 target=subsection reason=relative-heading-rank; base=h2

\subsection{第31章——基督被预定为天主子。}

因此，在祂内，我们的元首，愿恩宠的泉源彰显，这恩宠按各人的尺度，通过所有肢体分施给每个人。正是藉此恩宠，每个人从信德的开端成为基督徒，而正是藉此恩宠，那独一人从祂的开端成为基督。同一圣神也重生前者，正如后者由同一圣神所生。同一圣神在我们内成就罪的赦免，正如同一圣神使祂没有罪。天主当然预知祂要成就这些事。因此，这就是圣徒的预定，它在圣中之圣身上特别显明；凡是正确理解真理宣言的人，谁能否认这预定呢？因为我们得知，光荣的主本身被预定，只要那人是成为了天主子。外邦人的导师在他书信的开端呼喊说：\Quote{保禄，耶稣基督的仆人，蒙召为宗徒，被选拔去传天主的福音——这福音是天主曾藉其先知在圣经上所预许的——论及他的儿子，按肉身是生于达味的后裔，按圣德的神性，由于从死者中复活，被预定为具有大能的天主子。}因此，耶稣被预定，使那按肉身将成为达味后裔的祂，按圣德的神性，却成为大能的天主子，因为祂是由圣神和童贞玛利亚所生。这就是天主圣言对那唯一的人的提拔，以不可言喻的方式成就，使祂可以同时真实而恰当地被称为天主子与人之子——因所取的人性而为人子，因取人的独一神而为天主子，以致相信的是圣三而非圣四。人性的这种提升被预定，如此伟大、崇高、超越，无法再提高——就如同为我们之故，那神性不可能更谦卑地贬抑自己，比取人性连同肉身的软弱，甚至到十字架上的死亡。因此，正如那独一人被预定为我们的元首，同样，我们众多人被预定为祂的肢体。在此，让那些因亚当而失落的人类的功绩静默，让天主的恩宠作王，这恩宠藉我们的主耶稣基督、天主的独生子、唯一的主作王。凡能在我们的元首中找到那独特出生之先行功绩的人，就让他也在我们这些肢体中寻找那多次重生之先行功绩吧。因为那出生并非作为功劳报酬给基督，而是赐予的——也就是，祂由圣神和童贞女所生，脱离一切罪的纠缠。同样，我们由水和圣神的重生，也不是因任何功绩报酬给我们，而是白白赐予的；如果信德引领我们到达重生的洗礼，我们不应因此认为我们先行付出了什么，以致救恩的重生再次报酬给我们；因为那使我们相信基督的，正是那为我们造了一位我们信仰的基督。祂在人内开始并完成对耶稣的信德，而祂曾使那人耶稣成为信德的创始者和完成者；\BibleRef{希 12:2}因为，如你所知，祂在致希伯来人的书信中被如此称呼。\par

% structure: p0066 source=h2 target=subsection reason=relative-heading-rank; base=h2

\subsection{第32章 [XVI]——双重的召叫。}

天主确实召叫了许多祂预定得救的子女，使他们成为祂唯一预定得救的圣子的肢体，——并不是用那种召叫，即那些不愿来赴婚宴的人所受的召叫，因为那种召叫也临到了犹太人，对他们来说，被钉在十字架上的基督是绊脚石，也临到了外邦人，对他们来说，被钉在十字架上的基督是愚妄；而是用另一种召叫，即宗徒所区别出来的，当他宣讲基督——天主的智慧和天主的德能——给那些蒙召的人，无论是犹太人或希腊人时。故此他说：\BibleQuote{但蒙召的}（\BibleRef{格前 1:24}），这是为了表明，有些人并没有蒙召；因为他知道有一种确定的召叫，是那些按天主旨意蒙召之人所受的，天主已预知并预定他们效法祂圣子的形象。他所指的正是这种召叫，当他说：\BibleQuote{不是由于行为，而是出于召叫人的那一位；有话对她说：\InnerQuote{年长的要服事年幼的。}}（\BibleRef{罗 9:12}）难道他说的是\Quote{不是由于行为，而是出于相信的人}吗？实际上，他把这从人身上拿走了，好将一切都归于天主。因此他说\Quote{而是出于召叫人的那一位}——不是任何一种召叫，而是使人成为信徒的那种召叫。\par

% structure: p0068 source=h2 target=subsection reason=relative-heading-rank; base=h2

\subsection{三十三、恶人虽有能力犯罪，但藉此恶行做这事或那事，唯有天主有权。}

\Quote{天主的恩赐和召叫是不可撤回的。}\BibleRef{罗 11:29} 在那一句话中，也要稍微思考一下它的含义。因为当他说：\Quote{弟兄们，我不愿你们对此奥秘一无所知，免得你们自以为聪明：以色列人一部分变得顽梗，直到外邦人的全数进入，这样全以色列就得救；正如经上所记载：从熙雍将出来一位拯救者，从雅各伯中转离不敬；这就是我从我所立的与他们订立的盟约，那时我将除去他们的罪孽；}他立刻加上了一句需要非常谨慎理解的话：\Quote{论福音，他们因你们而成为仇敌；论拣选，他们因祖先而成为蒙爱者。}\BibleRef{罗 11:28} \Quote{论福音，他们因你们而成为仇敌}是什么意思？不就是说他们杀害基督的敌意，无疑正如我们所见的，对福音有益吗？他表明这是出于天主的安排，他知道如何善用恶事；不是为了使\OriginalTerm{义怒之器}{vessels of wrath}有利于祂，而是通过祂对它们的善用，它们可能有利于\OriginalTerm{怜悯之器}{vessels of mercy}。还有什么比实际所说的话更清楚的呢？\Quote{论福音，他们因你们而成为仇敌。}因此，行恶的能力在于恶人；但在行恶时他们通过这种恶行做这件事或那件事，不在于他们的能力，而在于天主，祂分开黑暗并调节它；以至于连他们所做的违背天主旨意的事，除非是天主的旨意，否则不会实现。我们在宗徒大事录中读到，当宗徒们被犹太人遣走，来到自己的朋友那里，向他们展示司祭和长老对他们所说的伟大事情时，他们全体一心向主扬声说：\Quote{主！你是创造天地、海洋和其中万物的天主；你藉着我们的祖先、你的圣仆达味的口曾说：外邦人为什么怒，人民为什么谋虚妄的事？地上的君王起来，首领聚在一起，反对主，反对他的受傅者。确实，在这座城里，黑落德和比拉多，以及以色列人，聚集起来反对你的圣仆耶稣，你所傅油的，要实行你的手和你的旨意所预定要实行的事。}看所说的：\Quote{论福音，他们因你们而成为仇敌。}因为天主的手和旨意预定这些事由敌对犹太人实行，为福音所必需，为了我们。但接着是什么？\Quote{论拣选，他们因祖先而成为蒙爱者。}因为那些在敌意中灭亡的\OriginalTerm{仇敌}{enemies}，以及同一民族中仍在对基督的反对中灭亡的人——他们是蒙拣选和\OriginalTerm{蒙爱者}{beloved}的吗？绝不！谁如此完全愚昧而说这话？但这两个表述，尽管彼此相反——即\Quote{仇敌}和\Quote{蒙爱者}——虽然不适用于同一个人，却适用于同一个犹太民族，同一个以色列的肉身子孙，其中有些人属于堕落，有些人属于以色列本身的祝福。因为使徒之前更清楚地解释了这个含义，当他说：\Quote{以色列所求的，没有得到；但蒙拣选的得到了，其余的都变得顽梗了？}\BibleRef{罗 11:7} 然而在这两种情况下，都是同一个以色列。因此，当我们听到\Quote{以色列没有得到}，或\Quote{其余的都变得顽梗了}，那些人应理解为因我们而成为仇敌的人；但当我们听到\Quote{蒙拣选的得到了}，那些人应理解为因祖先而成为蒙爱者，这些祖先确实得到了应许；因为\Quote{应许是对亚巴郎和他的后裔说的}\BibleRef{迦 3:16}，因此在那橄榄树上嫁接上了外邦人的野橄榄枝。而现在，我们当然应该认同那\OriginalTerm{拣选}{election}，关于它他说是按恩宠，不是按债，因为\Quote{凭着恩宠的拣选，有了残余}\BibleRef{罗 11:5}。这蒙拣选的得到了，其余的都变得顽梗。关于这拣选，以色列人因祖先的缘故而成为蒙爱者。因为他们不是被那种\Quote{许多人蒙召}的\OriginalTerm{召叫}{calling}所召，而是被那种蒙选者被召的召叫所召。因此，在他说了\Quote{论拣选，他们因祖先而成为蒙爱者}之后，他接着加上了那些话，由此引发了这个讨论：\Quote{因为天主的恩赐和召叫是不可撤回的}——就是说，它们坚定地确立，不可改变。属于这召叫的人都受天主教导；他们中没有人可以说：\Quote{我信了，是为了被这样召叫}，因为天主的怜悯先于他，因为他正是为此被召叫，以便他相信。因为所有受天主教导的人都来到圣子面前，因为他们从圣父那里藉着圣子听见并学到了，圣子最清楚地宣告：\Quote{凡从圣父听见并学习的，都到我这里来}\BibleRef{若 6:45}。但这样的人没有一个丧亡，因为\Quote{凡父赐给祂的，祂一个也不失落}\BibleRef{若 6:39}。因此，凡是属于这类的，都不会丧亡；凡丧亡的，也从未属于这这类。为此，经上说：\Quote{他们从我们中间出去，但他们不属于我们；因为如果他们属于我们，他们必会留在我们中间}\BibleRef{若 2:19}。\par

% structure: p0070 source=h2 target=subsection reason=relative-heading-rank; base=h2

\subsection{第三十四章 [XVII.]——蒙拣选者的特殊召叫不是因为他们相信了，而是为使他们相信。}

那么，让我们来理解那使他们成为被选者的\OriginalTerm{召唤}{vocatio}——不是那些因为他们相信而被选的人，而是那些被选好使他们相信的人。因为主自己也充分解释了这召唤，当祂说：\BibleQuote{不是你们拣选了我，而是我拣选了你们。} \BibleRef{若 15:16} 因为如果他们是因为相信而被选，那么他们自己必定首先藉着相信祂而拣选了祂，这样他们才配得被选。但祂完全消除了这假设，当祂说：\BibleQuote{不是你们拣选了我，而是我拣选了你们。} 然而他们自己无疑在相信祂时拣选了祂。因此祂说：\BibleQuote{不是你们拣选了我，而是我拣选了你们}，这并非由于其他原因，而是因为他们拣选祂并非为了祂拣选他们，而是祂拣选了他们好使他们能拣选祂；因为祂的仁慈按恩宠而非按债先于他们。因此，当祂披戴肉身时，祂从世界中拣选了他们，但他们是作为在创世以前已在祂内被拣选的人而被拣选的。这是关于\OriginalTerm{预定}{praedestinatio}和恩宠的不变真理。因为宗徒所说的\BibleQuote{就如祂在创世以前，在祂内拣选了我们} \BibleRef{弗 1:4} 又是什么呢？确实，如果这是说因为天主预知他们会相信，而不是因为祂自己要使他们成为信徒，那么圣子就反对这种预知，当祂说：\BibleQuote{不是你们拣选了我，而是我拣选了你们}；而天主更应该预知这件事本身，即他们自己会拣选祂，好使他们配得被祂拣选。因此，他们是在创世以前，按照那种预定而被选的，天主在这预定中预知祂自己要做什么；但他们是按照那种召唤而从世界中被选的，天主藉此召唤实现了祂所预定的。因为祂所预定的，祂也召唤了他们，即那按照旨意的召唤。因此，不是别人，而是祂所预定的那些人，祂也召唤了他们；也不是别人，而是祂如此召唤的人，祂也使他们成义；也不是别人，而是祂所预定、召唤、成义的人，祂也使他们得光荣；确实，直到那没有终结的终结。因此，天主拣选了信徒；但祂拣选他们是为了使他们如此，而不是因为他们已经如此。宗徒\Person{雅各伯}说：\BibleQuote{天主岂不是拣选了世上的贫穷人，使他们富于信德，并继承祂所应许给爱祂之人的国度吗？} \BibleRef{雅 2:5} 因此，藉着拣选他们，祂使他们富于信德，正如祂使他们成为国度的承继者；因为祂正确地说要在他们内拣选那为了在他们内造成而拣选他们的事。我问，谁能听见主说\BibleQuote{不是你们拣选了我，而是我拣选了你们}，竟敢说人相信是为了被选，而他们倒被选是为了相信，免得违反真理的判决，被人发现他们先拣选了基督，而基督对他们说\BibleQuote{不是你们拣选了我，而是我拣选了你们}？ \BibleRef{若 16:16}\par

% structure: p0072 source=h2 target=subsection reason=relative-heading-rank; base=h2

\subsection{拣选是为了成圣。}

谁能听见宗徒说：\BibleQuote{愿我们的主耶稣基督的天主和父亲受颂赞，祂在基督内，以天上各种属神的祝福，祝福了我们；就如祂在创世以前，在祂内拣选了我们，为使我们成为圣洁无瑕的；由于爱，预定我们藉着耶稣基督获得义子的名分，归于祂自己，按照祂旨意的喜悦，使祂恩宠的光荣得到颂赞，这恩宠是祂在爱子内赐予我们的；我们藉着祂的血获得救赎，罪过的赦免，是按照祂恩宠的丰富，这恩宠祂以各种智慧和明达赐予了我们，使我们知道祂旨意的奥秘，这旨意是祂在基督内所定的，为在时期圆满时，使天上和地上的万有，总归于基督元首；在祂内我们也获得了基业，因为我们是被预定的，按照那运行万有的旨意所定的计划，为使我们成为颂扬祂光荣的人。} — 我且说，谁能认真而智慧地听见这些话，还敢对我们所辩护的如此清楚的真理有任何怀疑呢？天主在创世以前，在基督内拣选了祂的肢体；祂怎能拣选那些尚未存在的人，除非藉着预定他们？因此，祂藉着预定我们而拣选了我们。祂会拣选不圣洁不洁净的人吗？如果问题被提出：祂是拣选这样的人，还是拣选圣洁无瑕的人？谁能在回答时犹豫，而不立即赞成圣洁纯洁的人呢？\par

% structure: p0074 source=h2 target=subsection reason=relative-heading-rank; base=h2

\subsection{天主拣选了义人；不是那些祂预知凭自身为义的人，而是那些祂预定使之成为义的人。}

\Quote{因此，}白拉奇派说，\Quote{祂预知谁将会因自由意志的选择而成为圣洁无瑕的，并基于这预知，在创世以前就拣选了他们。因此祂拣选了他们，}他说，\Quote{在他们存在之前，就预定他们成为祂预知将会成为圣洁无瑕的子女。祂当然没有使他们如此；祂也没有预见祂会使他们如此，而是预见他们会如此。}让我们来看看宗徒的话，看看祂是否因为我们将要成为圣洁无瑕而拣选了我们，还是为了使我们成为圣洁无瑕。\Quote{愿我们主耶稣基督的天主和父受赞美，祂在基督内，以各种属神的祝福，在天上祝福了我们，就如祂在创世以前，在祂内拣选了我们，为使我们成为圣洁无瑕的。}\BibleRef{弗 1:3} 那么，不是因为我们将会如此，而是为了使我们如此。这确实是确定的——确实是明显的。我们之所以会成为如此，当然是因为祂拣选了我们，预定我们因祂的恩宠而成为如此。因此，祂在基督耶稣内，以天上的各种属神祝福祝福了我们，就如祂在创世以前在祂内拣选了我们，为使我们成为圣洁无瑕的，好叫我们在这极大的恩宠恩赐中，不会因我们意志的善意而自夸。\Quote{在祂内，}他说，\Quote{祂在爱子内恩待了我们}——这当然是祂自己的旨意，祂恩待了我们。因此，可以说，祂以恩宠恩待了我们，正如说祂以正义使我们成为义人一样。\Quote{在祂内，}他说，\Quote{我们藉祂的血获得了救赎，罪过得赦免，这是按祂恩宠的丰富，这恩宠曾以各种智慧和明达丰富地赐予我们，使祂向我们揭示祂旨意的奥秘，这是按祂的善意。}在这旨意的奥秘中，祂按自己的善意，而非按我们的善意（除非祂自己按自己的善意帮助它成为善，否则我们的善意不可能成为善），放置了祂恩宠的丰富。但当祂说了\Quote{按祂的善意}后，祂又加上\Quote{这善意是祂在祂内所预定的}，即在祂爱子内，\Quote{到了时期圆满的职分，使天上和地上的万有，总归于基督元首；在祂内，我们也获得了基业，因为我们按那按祂旨意计划施行万有的祂的旨意，被预定成为祂荣耀的赞美。}\par

% structure: p0076 source=h2 target=subsection reason=relative-heading-rank; base=h2

\subsection{第三十七章——我们被拣选和预定，不是因为我们将要成为圣洁，而是为了使我们成为圣洁。}

若要对各个论点进行辩论，那就太冗长了。但你们无疑看见，你们看见宗徒的宣告何等明确地辩护了这恩宠，反对那将人的功德放在首位（好像人应先给予什么，以便得到回报）的主张。因此，天主在基督内，在创世以前拣选了我们，预定我们成为子女，不是因为我们将要依靠自己成为圣洁无瑕的，而是祂拣选并预定我们，为使我们成为如此。而且，祂这样做是按祂旨意的善意，好使没有人能因自己的意志而自夸，而要夸耀天主对他的旨意。祂这样做是按祂恩宠的丰富，按祂的善意，这善意是祂在爱子内所预定的；在祂内，我们获得了基业，是按那旨意——不是我们的，而是那万事都工作到甚至也在我们内工作愿意的那一位的旨意——被预定的。而且，祂是按祂旨意的计画行事，为使我们成为祂荣耀的赞美。\BibleRef{斐 2:13} 正是为此，我们呼喊：没有人应夸耀人，因此也不应夸耀自己；但谁若夸耀，当夸耀主，好使他成为祂荣耀的赞美。因为正是祂自己按祂的旨意工作，为使我们成为祂荣耀的赞美，当然也成为圣洁无瑕的，正是为此祂召叫了我们，在创世以前预定我们。出自这祂的旨意，就是那特别的召叫，为被拣选者，祂使万事互相效力，使他们获得益处，因为他们都是按祂的旨意蒙召的，并且\Quote{天主的恩赐和召叫是不会撤回的。}\BibleRef{罗 11:29}\par

% structure: p0078 source=h2 target=subsection reason=relative-heading-rank; base=h2

\subsection{第三十八章（第十九章）——白拉奇派和半白拉奇派关于预定的观点是什么？}

但我们现在所讨论、并为之代辩的这些弟兄们，也许会说，白拉奇派被这段证言驳倒了，因为其中说，我们在基督内被拣选、在创立世界以前被预定，为使我们因爱而在祂面前成为圣洁无瑕的。因为他们认为，\\Quote{我们接受了天主的诫命，就凭自由意志的选择在爱中使自己成为圣洁无瑕，在祂面前；天主预见了这事，}他们说，\\Quote{因此祂在创立世界以前就在基督内拣选并预定了我们。}虽然宗徒说，不是因为祂预知我们会成为那样，而是为了使我们能由同样的恩宠拣选成为那样，祂在祂的爱子内恩待了我们。因此，当祂预定我们时，祂预知了祂自己的工程，藉此工程祂使我们成为圣洁无瑕的。由此，白拉奇派的错误被这段证言正确地驳倒了。\\Quote{但我们说，}他们说，\\Quote{天主除了我们开始相信的信德之外，没有预知任何属于我们的东西；祂在创立世界以前就拣选并预定了我们，为使我们藉祂的恩宠和祂的工程成为圣洁无瑕的。}但也让他们在这段证言中听到他的话：\\BibleQuote{我们得了基业，是按那位成就万事的旨意被预定的}\\BibleRef{弗 1:11}因此，成就万事的祂，也成就我们相信的开端；因为信德本身并不先于那个召叫，关于这召叫，经上说：\\BibleQuote{因为天主的恩赐和召叫是不会反悔的}\\BibleRef{罗 11:29}并且经上说：\\BibleQuote{不在乎行为，而在乎那召人的}\\BibleRef{罗 9:12}（尽管祂也许可以说：\\Quote{在乎那相信的}）；以及主所表明的拣选，祂说：\\BibleQuote{不是你们拣选了我，而是我拣选了你们。}\\BibleRef{若 15:16}因为祂拣选了我们，不是因为我们相信了，而是为了我们可以相信，免得我们被认为先拣选了祂，以致祂的话落空（我们断乎不可这样想），\\BibleQuote{不是你们拣选了我，而是我拣选了你们。}我们被召也不是因为我们相信了，而是为了我们可以相信；藉着那不会反悔的召叫，我们得以相信的事实被成就并贯彻始终。但关于这事我们已经说了许多，无需重复。\par

% structure: p0080 source=h2 target=subsection reason=relative-heading-rank; base=h2

\subsection{第三十九章——信德的开端是天主的恩赐。}

最后，在这段证言后面，宗徒为那些相信的人感谢天主——当然不是因为福音已向他们宣讲，而是因为他们相信了。因为他说：\\BibleQuote{在祂内，你们听了真理之言，就是你们得救的福音；在祂内，你们相信之后，就受了所应许的圣神的印记，这圣神是我们得基业的凭据，直到天主所置的产业得赎，使祂的荣耀受到赞美。为此，我听见了你们在基督耶稣内的信德，并想念众圣徒，就不断为你们感谢天主。}他们的信德，在向他们宣讲福音时是新的、近期的；宗徒听见这信德，就为他们感谢天主。如果他为那他所认为或知道不是人给予的东西而感谢人，那该被称为奉承或嘲笑，而不是感谢。\\BibleQuote{不要错误，因为天主是轻慢不得的；}\\BibleRef{迦 6:7}因为信德的开端也是祂的恩赐，否则宗徒的感谢就会被视为错误或虚假。那么如何？难道得撒洛尼人的信德开端不也是这样吗？然而，同一位宗徒为他们感谢天主，说：\\BibleQuote{为此，我们也不断感谢天主，因为你们从我们接受了天主的言语，你们接受了，不以为是人的言语，而实在是天主的言语，这言语在你们信的人身上运行。}\\BibleRef{得前 2:13}他在这里为什么事感谢天主呢？倘若他所感谢的那位并没有亲自做这事，那感谢便是空洞无益的。但既然这不是空洞无益的，那么他所为这事感谢的天主，确实是亲自做了这事：使他们听见了天主的言语时，不以为是人的言语，而实在是天主的言语。因此，天主是按祂的旨意在人心中运行那召叫，关于这召叫我们已说了许多，使他们不至于徒然听福音，而是听见时，就悔改并相信，接受它不以为是人的言语，而实在是天主的言语。\par

% structure: p0082 source=h2 target=subsection reason=relative-heading-rank; base=h2

\subsection{第四十章（第二十节）——宗徒对信德开端是天主恩赐的见证。}

此外，我们受提醒：人信德的开端是天主的恩赐，因为宗徒在\BibleBook{哥罗森}{书}中指明这一点，他说：\BibleQuote{你们要恒心祈祷，在此儆醒感恩，也为我们祈求，使天主给我们开传道的门，能以讲基督的奥秘（我为此被捆绑），叫我按着所该说的话，将这奥秘发明出来。} 天主的道之门如何被打开？除非听者的心智被打开，使他能够相信，并且当他已开始了信德，就能接受那些为建立纯正道理而宣讲和论述的事，免得因不信而封闭的心拒绝并排斥所说的话。 因此，他也对格林多人说：\BibleQuote{但我要在厄弗所待到五旬节，因为有宽大又有功效的门为我开了，并且反对的人也多。}\BibleRef{格前 16:8} 这里还能有什么别的意思，无非就是：当福音首先由他在那里被宣讲时，许多人已经相信，而且出现了许多那同一个信德的反对者，这与主的话相符：\BibleQuote{除非蒙我父的恩赐，没有人能到我这里来；}\BibleRef{若 6:66} 以及\BibleQuote{天国的奥秘只给你们知道，但给别人，不给他们。}\BibleRef{路 8:10} 所以，门在那些蒙赐予的人那里是敞开的，但在那些未蒙赐予的人中间，却有许多反对者。\par

% structure: p0084 source=h2 target=subsection reason=relative-heading-rank; base=h2

\subsection{第四十一章——进一步的宗徒证言}

而且，同一宗徒在他的第二封书信中向同一群人说：\BibleQuote{我从前为基督的福音到了特洛阿，主也给我开了门，那时，因为没有遇见我的兄弟弟铎，我心神不安，便辞别他们，往马其顿去了。}\BibleRef{格后 2:12} 但请注意他接着说：\BibleQuote{感谢天主，祂时常使我们在基督内夸胜，并藉着我们在各处彰显那因认识基督而有的芬芳；因为我们是献给天主的基督的馨香，在得救的人中，在灭亡的人中：对后者，是死的香气使人死；对前者，是活的香气使人活。} 看看这位最热忱的战士、恩宠的不可战胜的捍卫者，因何事而感恩。看看他因何事而感恩——即宗徒们在那些蒙祂恩宠得救的人中，以及在那些因祂审判而灭亡的人中，都是献给天主的基督的馨香。但为使那些不太明白这些事的人不至于恼怒，他自己在加上这句话时给予警告：\BibleQuote{这样的事，谁能当得起呢？}\BibleRef{格后 2:16} 但让我们回到门的敞开，宗徒以此表示听者信德的开端。因为\BibleQuote{同时也为我们祈求，使天主给我们开传道的门，}\BibleRef{哥 4:3} 还有什么意思，难道不是极其明显的证明，即使信德的开端本身也是天主的恩赐吗？因为若不认为是由祂所赐，就不会在祈祷中向祂祈求。这天上的恩宠之赐曾降给那卖紫布的妇人（\BibleRef{宗 16:14}），正如圣经在\BibleBook{宗徒}{大事录}中所说：\BibleQuote{主开了她的心，使她留心听保禄所讲的话；}因为她蒙召正是为了让她相信。因为天主在人的心中随己意行事，或藉援助，或藉审判；以致甚至藉着他们，祂的手和旨意所预定要成就的事得以完成。\par

% structure: p0086 source=h2 target=subsection reason=relative-heading-rank; base=h2

\subsection{第四十二章——旧约的证言}

因此，反对者声称我们从\WorkTitle{列王纪}和\WorkTitle{编年纪}的圣经见证所证明的与我们正在讨论的主题无关，这也是徒然的；例如，当天主愿意做那只能由愿意的人来做的事时，他们的心就倾向愿意做这事——也就是说，被祂的能力所倾向，祂以一种奇妙且不可言喻的方式，也在我们内推动意愿。这若不是什么也没说，却还要反驳，又是什么呢？除非或许他们已向你们提出了他们观点的某种理由，而你们却宁愿在你们的信函中对此保持沉默。但我不知道那理由会是什么。是否可能，既然我们已经表明天主如此作用于人的心，并诱导祂所愿意之人的意志至此，使\Person{撒乌耳}或\Person{达味}被立为王——他们难道不认为这些实例适用于这个主题，因为在此世暂时为王与在天主内永远为王并非同一回事吗？因此，他们是否以为天主使祂所愿意之人的意志倾向获得地上的王国，却并不使他们倾向获得天国？但我认为，正是关乎天国而非地上的王国，才说了这些话：\Quote{请使我的心倾向祢的训言；}或\Quote{人的脚步由上主支配，祂愿意祂的道路；}或\Quote{人的意志由上主预备；}或\BibleQuote{愿我们的上主与我们同在，如同与我们的祖先同在一样；不要离弃我们，也不要转脸不顾我们；愿祂使我们的心倾向祂，使我们遵行祂的一切道路；}\BibleRef{列上 8:57}或\BibleQuote{我要赐给他们一颗认识我的心，和有耳能听的耳；}\BibleRef{巴 2:31}或\BibleQuote{我要赐给他们一颗新心，在他们五内注入一种新的精神。}\BibleRef{则 11:19}让他们也听听这话：\BibleQuote{我要将我的圣神放在你们内，使你们遵行我的义德；你们要遵守我的诫命，且一一实行。}\BibleRef{则 36:27}让他们听：\BibleQuote{人的脚步由上主支配，人怎能了解自己的道路？}\BibleRef{箴 20:24}让他们听：\BibleQuote{人的道路看来正直，但上主却衡量人心。}\BibleRef{箴 21:2}让他们听：\BibleQuote{凡预定得永生的人都信了。}\BibleRef{宗 13:48}让他们听听这些经文，以及我未提及的其他此类经文，其中宣告天主预备并转变人的意志，甚至是为天国和永生。然后想一想，相信天主为人建立地上的王国而运作人的意志，但人却靠自己意志的努力去获得天国，这是何等的事情。\par

% structure: p0088 source=h2 target=subsection reason=relative-heading-rank; base=h2

\subsection{第四十三章——结论}

我已经说了很多，或许我早已能说服你们我所希望的，而我还在对这样聪慧的人说话，好像他们是迟钝的，对他们而言，即使太多也不够。但请他们宽恕我，因为一个新问题迫使我如此。因为，虽然在我先前的小作品中，我已经用足够合适的证据证明信德也是天主的恩赐，但还是出现了这个反驳的理由，即：那些见证对于证明信德的增长是天主的恩赐是有用的，但信德的开端，即人最初信从基督，乃出于人自己，并非天主的恩赐——但天主要求这个，以便当它先行之后，其他恩赐可以随之而来，仿佛是建立在这功劳的基础上，而这些恩赐是天主的恩赐；并且其中没有一样是白白赐予的，尽管在这些恩赐中，天主的恩宠被宣告，而这恩宠除非是无偿的，否则就不是恩宠。你们看这一切是多么荒谬。因此我决定，尽我所能地阐明，这开端本身也是天主的恩赐。如果我做得比那些为此缘故而希望我这样做的人所愿的更为详尽，我准备好受他们的责备，只要他们仍然承认，尽管比他们希望的更长，尽管令明白人感到厌恶和疲倦，我还是做了我所做的：也就是说，我教导了，就连信德的开端，如同贞洁、忍耐、义德、孝爱，以及其他他们并不争议的美德，也是天主的恩赐。因此，让这成为本篇论文的结束，免得篇幅过长引起反感。\par

\SourceRecord{Translated by Peter Holmes and Robert Ernest Wallis, and revised by Benjamin B. Warfield.；Nicene and Post-Nicene Fathers, First Series；Vol. 5.；Edited by Philip Schaff.；Buffalo, NY: Christian Literature Publishing Co.,；1887.；<http://www.newadvent.org/fathers/15121.htm>.}

% structure: p0001 source=h1 target=section reason=page-title

\section{\WorkTitle{论圣人的预定（第二卷）}}

\Emph{在第一部分中，他证明使人坚持到底而成为基督的恒心是天主的恩赐；因为向上主祈求不认为是由他给予的事，那是嘲弄。此外，在主祷文中，几乎只祈求恒心，根据殉道者西彼廉的诠释，藉此诠释，这恩宠的敌人在他们出生之前就已定罪。他教导说，恒心的恩宠不是按领受者的功德赐予的，而是有些蒙天主的仁慈获得，有些则因他的公义审判而未得。至于为何在成年人中，一个被召叫而另一个不被召叫，这是不可探究的；同样，在两个婴儿中，一个被接去而另一个被留下，也是不可探究的。但更不可探究的是，在两个虔诚的人中，为何一个被赐予恒心，而另一个没有被赐予；但这是最确定的：前者是预定的，后者不是。他注意到，预定的奥秘在吾主关于提洛和漆冬人的话中显明，这些人如果曾行过在苛辣匝因所行的奇迹，他们就会悔改。他表明婴儿的例子有力地证实了老年人之预定和恩宠的真理；他答复了他的第三卷论自由意志中被他对手在此问题上错误引用的段落。随后，在本作品的第二部分，他反驳他们所说的，即预定的定义与劝诫和谴责的用处相抵触。他另一方面断言，宣讲预定是有益的，以便人不夸耀自己，而夸耀主。至于他们反对预定的异议，他表明同样的异议几乎可以同样地扭转来反对天主的预知或反对他们一致认为其它美事所必需的恩宠（除了信德的开端和坚持的完成）。因为圣人的预定无非就是天主对他恩惠的预知和预备，藉此，凡被救的人都确实被救。但他吩咐应以和谐的方式宣讲预定，而不是以那样的方式，以致于让不熟练的群众看来好像因其宣讲本身而被反驳。最后，他向我们将置于眼前的耶稣基督推荐为预定最卓越的实例。}\par

% structure: p0003 source=h2 target=subsection reason=relative-heading-rank; base=h2

\subsection{第一章 [I.]——论此处所讨论的恒心的本质。}

我现在需要以更大的关怀来考虑恒心的问题；因为在先前那卷书中，我在讨论信德的开端时也说过一些这方面的话。因此，我断言，使我们坚持到底而成为基督的恒心是天主的恩赐；我称那终点为终结，在此生命中唯独有跌倒的危险。因此，只要一个人还活着，就不确定他是否领受了这恩赐。因为如果他在死前跌倒了，当然就称他不曾坚持；这样说非常正确。那么，一个没有坚持到底的人，怎能被称为领受了或拥有过恒心呢？因为如果一个人拥有节德，却从这美德堕落而变得不节制——或者，同样，如果他拥有正义、忍耐，甚至信德，却跌倒了，那么他正确地说曾经拥有这些德行而不再拥有；因为他曾节德，或他曾正义，或他曾忍耐，或他曾有信德，只要他如此；但当他不再如此时，他就不再是以前的自己了。但一个没有坚持到底的人，何曾有过坚持？因为只有通过坚持，一个人才能显明自己是坚持的——而他没有这样做。但恐怕有人会反驳并说：如果从一个人成为信徒的时候算起，他活了——为了论证——十年，而在其中从信德中跌倒，他岂没有坚持了五年？我不在字面上争论。如果这也应被称为恒心，仅就其持续的时间而言，那么他当然一点也没有拥有我们现在所讨论的恒心，即使人坚持到底而成为基督的恒心。一个一年的信徒，或者一个可以想象得更短时期的信徒，如果他忠信地活到死，他就比一个信了许多年却在其死前不久从信德的坚定中堕落的人，拥有更多的这恒心。\par

% structure: p0005 source=h2 target=subsection reason=relative-heading-rank; base=h2

\subsection{第二章 [II.]——信德是基督徒的开端。为主殉道是最佳结局。}

此事既已确定，让我们看看这恒心——经上所说\BibleQuote{恒心至终的，必然得救}\BibleRef{玛 10:22}——是否是天主的恩赐。若不然，宗徒的这话又怎会真实呢：\BibleQuote{因为基督赐给了你们，不仅为相信祂，而且也为祂受苦}\BibleRef{斐 2:29}？当然，这些事中，一个关乎开始，另一个关乎终结。然而，两者皆是天主的恩赐，因为两者都被说成是赐给的；正如我上面已经说过的。因为对基督徒来说，还有什么比相信基督更真实地是开始呢？还有什么结局比为基督受苦更好呢？但就相信基督而言，无论发现了何种反对意见，即认为不是信德的开始，而是信德的增长才应称为天主的恩赐——对此意见，靠着天主的恩赐，我已经答复得足够，甚至绰绰有余。但有什么理由可以证明，恒心至终不应在基督内赐给那已获得为基督受苦之恩赐的人，或者更明确地说，那已获得为基督而死之恩赐的人呢？因为宗徒伯多禄表明这是天主的恩赐，他说：\BibleQuote{如果天主的旨意是这样，为行善而受苦，比为作恶而受苦更好}\BibleRef{伯前 3:17}。当他说\BibleQuote{如果天主的旨意是这样}时，他说明这是由天主赐予的，然而并非赐给所有圣人，使他们为基督受苦。因为那些天主不愿使之获得受苦经验与荣耀的人，如果他们恒心在基督内至终，也绝不会错过天主的国。可是，谁能说，那些因身体虚弱或因任何意外而在基督内死去的人，这恒心不是赐给他们的呢？尽管一种更为艰难的恒心赐给了那些甚至为基督而忍受死亡本身的人。因为当迫害者竭力阻止人恒心时，恒心要艰难得多；因此，他在恒心至死中得以支持。因此，拥有前一种恒心更难——拥有后一种更容易；但对于无所不能的天主，赐予两者都是容易的。因为天主曾应许这事说：\BibleQuote{我要把我的敬畏放在他们心里，使他们不离开我}\BibleRef{耶 32:40}。这若不是指：\Quote{我要将如此巨大敬畏放进他们心里，使他们恒心依恋我}又是什么呢？\par

% structure: p0007 source=h2 target=subsection reason=relative-heading-rank; base=h2

\subsection{第三章　人要求恒心，因为它是天主的恩赐。}

但为何这恒心要从天主那里被祈求，如果它不是由天主赐予的呢？难道这祈求也是嘲弄，即向天主祈求明知祂不会赐予、却是在人能力之内的事？正如感恩的谢恩是嘲弄一样，如果为天主未曾赐予或未曾做成的事而感谢祂。但我在那里所说的，我在这里也要再说一遍：\BibleQuote{不要受骗，天主是嘲笑不得的}\BibleRef{迦 6:6}。人啊，天主不仅是你的话语的见证，也是你思想的见证。如果你诚心诚意地向如此富足的一位求任何事，你要相信，你从你所祈求的那位那里领受了你所求的。你要停止用嘴唇尊敬祂，而在心里自高自大，认为你凭自己有你所假装向祂祈求的。难道这恒心不是从祂那里祈求的吗？说这话的人不应以任何论据来责备，而应以圣人们的祈祷来压倒。难道有哪位圣人不是向天主为自己祈求恒心于祂吗？就在那被称为\Quote{主祷文}的祈祷中——因为主教导了它——当圣人们祈祷时，几乎除了恒心之外，还能理解为求什么别的呢？\par

% structure: p0009 source=h2 target=subsection reason=relative-heading-rank; base=h2

\subsection{第四章　白拉奇主义的三个主要论点。}

请稍加留意真福殉道者西彼廉所著论及此事的论文\WorkTitle{论天主经}中的阐述，看看多少年前，以及何等样的解毒剂，已被预备来对抗白拉奇主义者日后将使用的毒药。因为如您所知，天主教会主要针对他们维护三点。其一：天主的恩宠不是按我们的功绩赐予的，因为义人的每一项功绩都是天主的恩赐，由天主的恩宠赋予。其二：没有人能在可朽的肉体内生活而无任何罪愆，无论他多么正义。其三：人生来便因第一人的罪而应受惩罚，被定罪锁链束缚，除非由生育所沾染的罪债借重生得以解除。在这三点中，我最后提及的那一点，在上述荣耀殉道者的著作中并未论及；但另外两点在那里讲述得如此清晰，以致这些异端者——恩宠基督的现代仇敌——在他们出生之前很久就被判有罪了。因此，在这些圣人的功绩中——它们唯有作为天主的恩赐才算功绩——他说坚忍也是天主的恩赐，用这些话：\Quote{我们说：\InnerQuote{愿你的名被尊为圣；}并非我们祈求天主因我们的祈祷而成圣，而是我们恳求祂，使祂的名在我们内被尊为圣。但谁能使天主成圣呢？既然祂自己成圣？然而，因为祂说：\InnerQuote{你们应当是圣的，因为我是圣的}，我们就祈求与恳求，使我们这些在圣洗中成圣的人，能坚忍在我们已经开始的状态中。}稍后，在继续论述同一主题时，教导我们向主恳求坚忍——若这并非祂的恩赐，我们决不可能正当地、真实地这样做——他说：\Quote{我们祈求这圣化常在我们内；因为我们的主和审判者警告那被祂治愈和复活的人不要再犯罪，免得遭遇更坏的事，我们在不断的祈祷中做此祈求；我们日夜恳求，那从恩宠所领受的圣化和复活，能藉祂的保护得以保存。}因此，那位导师懂得，当我们已成圣的人说\Quote{愿你的名被尊为圣}时，我们是在向祂求坚忍于圣化，即我们应在圣化中坚忍。因为，除了祈求已领受之物也赐予我们不懈怠地持有它之外，还能是什么呢？所以，正如圣人祈求天主使他成圣时，实际上是祈求他继续成圣；同样，贞洁者祈求他贞洁，节制者祈求他节制，义人祈求他正义，虔敬者祈求他虔敬，等等——这些事物，反对白拉奇主义者时我们坚持是天主的恩赐——无疑是在祈求他们坚忍于他们承认已领受的善。如果他们领受了这个，无疑他们也领受了坚忍本身，这伟大的恩赐，天主的其他恩赐藉以得到保存。\par

% structure: p0011 source=h2 target=subsection reason=relative-heading-rank; base=h2

\subsection{第五章 天主经中的第二项祈求}

当我们说\Quote{愿你的国来临}时，我们祈求的若非是那无疑会来临于所有圣人之物也来临于我们，还能是什么？因此，在此处，那些已为圣的人祈祷，除求他们能坚忍于所赐予的圣德外，还求什么呢？因为天主的国是不会以别的方式临于他们的；那确实只临于那些坚忍到底的人，而非他人。\par

% structure: p0013 source=h2 target=subsection reason=relative-heading-rank; base=h2

\subsection{第六章 [III.] 第三项祈求 如何在天主经中理解天和地}

第三个祈求是：\Quote{愿祢的旨意奉行在天地之间}；或在许多抄本中读作，并且祈求者更常使用的是：\Quote{如同在天上，也在地上同样实行}，许多人理解为：\Quote{如同圣天使，也让我们奉行祢的旨意}。这位导师兼殉道者将天和地理解为灵与肉，并说我们祈求以二者的完全和谐奉行天主的旨意。他在这些话中也看到了另一个意义，与最正确的信仰相吻合，这个意义我已在上文谈到——即，为非信徒祈祷，他们如今仍是\Emph{地}，仅在第一次出生中带有属地之人，信徒被理解为在祈祷，他们披戴了属天之人，被无理由地称为\Emph{天}；他清楚表明信德的开端也是天主的恩赐，因为神圣教会不仅为信徒祈祷，求信德增长或在他们内持守，而且还为非信徒祈祷，求他们开始拥有他们根本未曾拥有、并且反而怀有敌对情感的事物。然而，现在我并非讨论信德的开端，这在上一卷中已多次谈到，而是讨论必须持守到底的恒心——这恒心连奉行天主旨意的圣人们也在祈祷中说\Quote{愿祢的旨意奉行}时寻求。因为，既然这旨意已在他们内实现，为何他们仍求它奉行，除非是为持守他们已经开始成为的？不过，这里可以说，圣人们不是求天主的旨意奉行在天上，而是求在地上如同在天上——也就是说，让地模仿天，即，让人模仿天使，或让非信徒模仿信徒；可见圣人们是求那尚未存在的事物得以存在，而不是求那已存在的事物继续存在。因为，无论人因何种圣德而卓越，他们尚未等同于天主的众天使；因此，天主的旨意尚未在他们内如同在天上那样奉行。如果真是这样，那么在我们祈求非信徒成为信徒的那一部分，所祈求的似乎不是恒心，而是开端；但在我们祈求人在奉行天主旨意上等同于天主的众天使的部分——当圣人们为此祈祷时，他们是在祈求恒心；因为除非一个人在地上所接受的圣德上持守到底，否则没有人能达到天上国度中的至高福乐。\par

% structure: p0015 source=h2 target=subsection reason=relative-heading-rank; base=h2

\subsection{第七章【四】——第四项祈求}

第四项祈求是：\Quote{求祢今天赐给我们日用的食粮}，\BibleRef{玛 6:11}。其中真福西彼廉指出了如何在这里也理解为祈求恒心。因为他论及其他事情时说：\Quote{我们祈求这食粮每日赐给我们，使我们这些在基督内、每日领受圣体作为救恩食粮的人，不致因某种大罪的介入，被禁止领受而无法分食天上的食粮，从而与基督的身体分离。}这位天主圣人的话表明，圣人们以这个意向祈求\Quote{求祢今天赐给我们日用的食粮}，就是直接从天主祈求恒心，使他们不致与基督的身体分离，而能持续在圣德之中，不容许任何足以使他们因此与基督的身体分离的罪行。\par

% structure: p0017 source=h2 target=subsection reason=relative-heading-rank; base=h2

\subsection{第八章【五】——第五项祈求。培拉基派人认为义人无罪，是错误的。}

在祈祷的第五句中，我们说：\Quote{宽免我们的罪债，如同我们也宽免我们的亏欠者}，\BibleRef{玛 6:12}。在这个祈求中，似乎没有祈求恒心。因为我们求赦免的罪是过去的，但拯救我们进入永生的恒心，在今世的生命中是必要的；但不是为过去的时间，而是为直到结束的剩余时间。然而，略加思考是值得的，即在这个祈求中，后来才出现的异端者如何被西彼廉的舌头刺穿，如同被真理最无坚不摧的标枪所刺。因为培拉基派人甚至敢这样说：义人在今生毫无罪恶，而且在这种人身上，甚至现今已经有一个没有瑕疵、没有皱纹、或任何类似之物的教会（\BibleRef{弗 5:27}），她就是基督独一无二的新娘；仿佛她不是那位新娘，她遍布全地，说着从祂所学来的话：\Quote{宽免我们的罪债。}但是请看，最光荣的西彼廉如何摧毁这些说法。当他解释主祷文的那一节时，他论及其他事情这样说：\Quote{我们受劝告是多么必要、多么明智、多么有益，因为我们被迫恳求宽恕我们的罪；当从天主祈求赦免时，灵魂回忆起自己的良心。以免有人自诩无辜，因自高而更深地灭亡，他被教导和指教，他每日犯罪，因为他每日被命令恳求赦免自己的罪。同样，若望在他的书信中也警告我们，说（\BibleRef{若一 1:8}）：\BibleQuote{如果我们说我们没有罪，便是欺骗自己，真理不在我们内。}}其余的内容，若在此插入则过于冗长。\par

% structure: p0019 source=h2 target=subsection reason=relative-heading-rank; base=h2

\subsection{第九章——当恩赐给人恒心时，他必能恒心。}

此外，当圣人们说：\BibleQuote{不要让我们陷入诱惑，但救我们脱离凶恶}，\BibleRef{玛 6:13}，他们祈求的不是别的，乃是能持守于圣德之中。因为，确实，当天主的这个恩赐赐给他们时——这恩赐已相当清楚地显示是天主的恩赐，因为它是向祂祈求的——那么，既然天主的恩赐赐给他们，使他们不致陷入诱惑，就没有一个圣人不能持守圣德直到终末。因为除非他首先被引入诱惑，否则没有人会停止持守基督徒的目标。因此，如果按照他的祈祷，恩赐给他不致被诱惑，那么他确实因天主的恩赐而持守于他所领受的圣化之中。\par

% structure: p0021 source=h2 target=subsection reason=relative-heading-rank; base=h2

\subsection{第十章 [VI.]——恒心之恩可由祈祷获得。}

但是您写道：\Quote{这些弟兄不愿意如此宣讲恒心，似乎恒心不能藉祈祷获得，也不能因顽固而失去。}在此他们对自己所说的不甚留意。因为我们所谈论的恒心，是使人坚持到底的恒心；如果这恒心被赐予，人就确实坚持到底；但若人没有坚持到底，那么恒心就没有被赐予，这一点我在上面已经充分讨论过了。因此，除非终点本身已经到来，并且被赐予恒心的人被发现坚持到了终点，否则任何人都不应说恒心被赐予了他。诚然，我们称我们所认识的一位贞洁者贞洁，无论他是否在同样的贞洁中继续或停止；如果他拥有任何其他可以保持和失去的神圣恩赐，我们在他拥有时就说他拥有它；如果他失去了它，我们就说他曾经拥有过它。但是，既然只有坚持到底的人才有坚持到底的恒心，那么许多人可能拥有它，但没有人能失去它。因为不必担心，当一个人坚持到底之后，或许会有某种邪恶的意志在他内兴起，以致他不再坚持到底。因此，天主的这一恩赐可以通过祈祷获得，但一旦被赐予，就不能因倔强而失去。因为当一个人已经坚持到底时，他既不能失去这个恩赐，也不能失去他在终点之前可能失去的其他恩赐。那么，那使人连可能失去的也不失去的恩赐，又如何能被失去呢？\par

% structure: p0023 source=h2 target=subsection reason=relative-heading-rank; base=h2

\subsection{第十一章——为恒心祈祷的效果。}

但是，恐怕有人会说，坚持到底的恒心一旦被赐予——即当一个人坚持到底时——确实不会失去，但在某种意义上，当一个人因倔强而行事，以致无法达到恒心时，恒心就被失去了；正如我们说，一个没有坚持到底的人失去了永生或天主的国，不是因为他已经领受并实际拥有它，而是因为他本会领受并拥有它，如果他坚持到底的话。——让我们撇开言辞的争论，承认有些东西即使尚未拥有，只是希望拥有，也可能失去。谁敢告诉我，天主不能赐予祂命令人向祂祈求的事物？肯定地说，这样断言的人，我不说他是愚昧人，而是疯狂。但天主命令祂的圣徒在祈祷中对祂说：\Quote{不要让我们陷入诱惑。}因此，凡在祈求此事时蒙垂听的人，就不会陷入倔强的诱惑，从而使自己配得或将要失去圣洁的恒心。\par

% structure: p0025 source=h2 target=subsection reason=relative-heading-rank; base=h2

\subsection{第十二章——人出于自己的意志离弃天主，故应被天主离弃。}

但是，另一方面，\Quote{人出于自己的意志离弃天主，因此应被天主离弃。}谁能否认这一点呢？正是为此，我们祈求不要陷入诱惑，以免这事发生。如果我们蒙垂听，这事当然不会发生，因为天主不允许它发生。因为除了祂自己所做的事，或祂所允许发生的事，没有别的事发生。因此，祂既能将意志从邪恶转向良善，又能使倾向跌倒的人回转，或引导他们走上祂所喜悦的道路。因为向祂说：\Quote{天主啊，求祢回转，使我们复苏}并非徒然；说\Quote{求祢不让我的脚动摇}并非徒然；说\Quote{上主，求祢不要将我由我的愿望交给罪人}并非徒然；最后，不一列举许多经文，因为可能您会想起更多，说\BibleQuote{不要让我们陷入诱惑}\BibleRef{玛 6:13} 并非徒然。因为凡不陷入诱惑的人，当然不会陷入自己邪恶意志的诱惑；而凡不陷入自己邪恶意志的诱惑的人，就绝对没有陷入任何诱惑。因为\BibleQuote{每人受诱惑，都是被自己的私欲所勾引、所饵诱}\BibleRef{雅 1:14}，\BibleQuote{天主不诱惑人}\BibleRef{雅 1:13}——也就是说，不以有害的诱惑诱惑人。因为还有一种有益的诱惑，借此我们不被欺骗或压倒，而是被考验，正如经上所说：\Quote{上主，求祢考验我，察验我。}因此，对于宗徒所说的那种有害的诱惑：\BibleQuote{怕那诱惑者诱惑了你们，以致我们的劳苦白费了}\BibleRef{得前 3:5}，\BibleQuote{天主不诱惑人}，如我所说——也就是说，祂不引人或带人进入诱惑。因为受诱惑而不陷入诱惑并非邪恶——甚至是有益的；因为这就是被考验。因此，当我们对天主说\Quote{不要让我们陷入诱惑}时，我们说的不就是\Quote{不容许我们被引入}吗？为此，有些人如此祈祷，在许多抄本中也读到如此，最有福的\Person{息彼连}也这样使用：\Quote{求祢不容许我们被引入诱惑。}然而，在希腊福音中，我所见的总是\Quote{不要让我们陷入诱惑}。因此，如果我们把一切完全交托给天主，而不把自己一部分交托给祂、一部分留给自己，我们就活得更安全，正如那位可敬的殉道者所看到的。因为当他解释祈祷的同一句时，他在别处说：\Quote{当我们祈求不要陷入诱惑时，我们由此祈求而记起自己的软弱和不足，免得有人傲慢自夸——免得有人骄傲自负——免得有人将宣认或受苦的光荣归于自己；因为主亲自教导谦卑说：\InnerQuote{你们要醒寤祈祷，免得陷入诱惑；心神固然愿意，肉体却软弱。}因此，当谦卑顺从的宣认首先发出，并将一切归于天主时，凡以敬畏天主之心恳切祈求的，都可由祂的慈爱赐予。}\par

% structure: p0027 source=h2 target=subsection reason=relative-heading-rank; base=h2

\subsection{第十三章 [VII.]——诱惑是人的境况。}

倘无其他证明，仅此主祷文便足以为我所辩护的恩宠辩护；因为它使我们毫无可自夸之处，表明我们不离开天主并非由我们自己获得，而是必须向天主祈求。因为凡不入诱惑者，必不离开天主。这绝非现今这样自由意志的力量所能及；但在人堕落之前，它曾经有过。然而，这种自由意志在最初状态的卓越中能成就多少，在天使身上显明了：当魔鬼及其使者堕落时，他们屹立于真理中，并获致那永不跌倒的永恒安全，我们确知他们现今已安于此境。但在人堕落之后，天主只愿人接近祂属于祂的恩宠；也只愿人离开祂不属于祂的恩宠。\par

% structure: p0029 source=h2 target=subsection reason=relative-heading-rank; base=h2

\subsection{第十四章——使人归向祂且不离开祂的，是天主的恩宠。}

祂将这恩宠置于\Quote{祂内，因为我们是在祂内蒙拣选，按照那安排万者的预旨而预定}的。\BibleRef{弗 1:11} 因此，正如祂运作使我们归向祂，同样祂运作使我们不离开。为此，借先知之口曾对祂说：\Quote{愿祢的手在祢右边的人身上，在祢为自己坚固的人子身上，我们就不会离开祢。}这当然不是那使我们离开祂的第一个亚当，而是那第二个亚当，祂的手按在祂身上，使我们不离开祂。因为基督与其肢体完全——为了教会，祂的身体——是祂的圆满。因此，当天主的手在祂身上，使我们不离开天主时，天主的作为无疑临到我们（因为这就是天主的手）；借着天主的这个作为，我们得以在基督内与天主常相存留——不像在亚当内那样离开天主。因为\Quote{在基督内我们蒙拣选，按那安排万者的旨意预定}。因此，使我们不离开天主的，是祂的手，而非我们自己的手。我说，这是那说过\Quote{我要把敬畏我的心放在他们心里，使他们不离开我}之祂的手。\BibleRef{耶 32:40}\par

% structure: p0031 source=h2 target=subsection reason=relative-heading-rank; base=h2

\subsection{第十五章——为何天主愿意我们为那本可不经祈祷就赐下的事而祈祷。}

所以祂也愿意我们祈求不被引入诱惑，因为我们若不被引，就绝不离开祂。这事纵使我们不祈求，祂也可以赐予；但祂愿借我们的祈祷提醒我们，这些恩惠是从谁领受的。因为我们不从祂领受，又从谁领受呢？合理的祈求当向祂献上。在这事上，教会无需寻求劳苦的争辩，只需思考自己的日常祈祷。教会祈祷，使不信者得以相信；因此天主使之皈依信仰。教会祈祷，使信徒得以持守；因此天主赐予恒心至终。天主预知祂必行此事。这正是圣人们的预定，\Quote{祂在创世以前，在基督内拣选了我们，为使我们成为圣洁无瑕，在祂面前因爱而成为圣洁无瑕；按祂旨意的喜悦，预定我们借着耶稣基督成为祂的义子，归于祂自己，为使祂恩宠的光荣受赞美，在祂的爱子内赐予我们恩惠；在祂内我们借着祂的血得蒙救赎，罪过的赦免，是按祂恩宠的丰富，这恩宠曾以各种智慧和明达充盈我们；使祂向我们显示祂旨意的奥秘，按祂在祂内所预定的喜悦，在时期圆满的经纶中，使天上和地上的万物在基督内总归于一；在祂内我们也蒙拣选，按那安排万者的旨意预定。}\BibleRef{弗 1:4} 面对如此清晰的真理号角，哪个清醒而警醒的信徒会接受人的论证呢？\par

% structure: p0033 source=h2 target=subsection reason=relative-heading-rank; base=h2

\subsection{第十六章 [VIII.]——为何恩宠不按功德施予？}

但有人说：\Quote{为何天主的恩宠不按人的功德施予？}我答：因为天主是仁慈的。又问：\Quote{那么为何不施予众人？}我于此答：因为天主是审判者。\BibleRef{罗 9:20} 如此，恩宠由祂白白施予；而借着祂公义的判决，在有些人身上显明恩宠赐给那些领受者的是什么。因此，我们不要忘恩负义：按祂旨意的喜悦，仁慈的天主从如此应得的灭亡中拯救如此多人，为使祂恩宠的光荣受赞美；倘若祂不拯救一人，祂也不是不义的。所以，得救者应爱祂的恩宠，不得救者应承认自己的罪有应得。若宽免债务显出良善，若追讨债务显出公义——在天主那里绝不会有不义。\par

% structure: p0035 source=h2 target=subsection reason=relative-heading-rank; base=h2

\subsection{第十七章——拣选一人而遗弃另一人的区别之难解。}

\Quote{但，}它说，\Quote{同样的情况，不仅对于婴儿，甚至对于双胞胎，为何审判如此不同？} 这不也是一个类似的问题吗：\Quote{为何在另一种情况下审判相同？} 那么，让我们回想那些在葡萄园里工作了一整天的工人和那些只工作了一小时的工人。诚然，劳动量不同，但支付工资的审判却相同。难道发牢骚的人从家主那里听到了别的，只有：\Quote{这是我的意愿}吗？诚然，他对一些人的慷慨，使他对他人没有不义。这两种人确实都在善人之列。然而，就正义和恩宠而言，可以真实地对被定罪的罪人和被拯救的罪人说：\BibleQuote{拿你的，去吧；}\BibleRef{玛 20:14}\Quote{我要给这人那不欠的；}\Quote{难道我不可以随意做我所愿意的吗？因为我善良，你就眼红吗？} 如果他问：\Quote{为何不也给我？} 他会听到，而且合情合理地听到：\BibleQuote{你是谁，人啊，竟敢向天主抗辩？}\BibleRef{罗 9:20} 虽然在你身上你看到一位最仁慈的恩主，而在你自己身上你看到一位最公义的索债者，但在两者身上你都看不到一位不公义的天主。因为即使他惩罚两人，他也是公义的，得救者有理由感恩，被定罪者却没有理由抱怨。\par

% structure: p0037 source=h2 target=subsection reason=relative-heading-rank; base=h2

\subsection{第十八章——但为何一人受罚多于另一人？}

\Quote{但是，}它说，\Quote{既然有必要，尽管并非所有人都被定罪，他仍应显示所有人的应得，从而将他的恩宠更自由地赐予怜悯之器；为何在同样的情况下，他会惩罚我多于另一个人，或拯救他多于我？} 我不这样说。如果你问为什么；因为我承认我无法作答。如果你进一步问为什么，那是因为在此事上，正如他的义怒是公义的和他的慈悲是伟大的，同样，他的判断是不可测度的。\par

% structure: p0039 source=h2 target=subsection reason=relative-heading-rank; base=h2

\subsection{第十九章——为何天主将那些将会坚持到底者与那些不会者混合在一起？}

让寻问者继续说：\Quote{为什么天主没有将坚持到底的恩赐赐给那些诚心敬拜他的人？} 除了因为那说话不虚假者说：\BibleQuote{他们从我们中间出去，但不是属于我们的；因为他们若是属于我们的，就必仍旧与我们同在。}\BibleRef{若一 2:19} 难道有两种人性吗？决不是。如果有两种人性，就没有恩宠了，因为如果恩宠是偿还给本性的债务，就没有无偿的拯救。但人们认为，所有看起来是善的信徒都应当接受坚持到底的恩赐。但天主判断为更好的是，将他的一些圣徒中混合一些不会坚持的人，这样那些不希望在此生不被试探的人就不会有安全感。因为宗徒说的话制止了许多人的有害骄傲：\BibleQuote{所以，那自以为站得稳的，要小心，免得跌倒。}\BibleRef{格前 10:12} 但跌倒者是因自己的意志跌倒，站立者因天主的意志站立。\BibleQuote{因为天主能使他站住；}\BibleRef{罗 14:4} 因此，他自己不能使自己站立，只有天主能。然而，不好高骛远而存有敬畏是好的。此外，每个人是在自己的思想中或跌倒或站立。正如宗徒所说，我在之前的论文中也提到：\BibleQuote{我们并不足以凭自己思想什么，我们的充足乃出于天主。}\BibleRef{格后 3:5} 跟随他的还有真福\Person{安波罗削}，他敢于说：\Quote{因为我们的心不在我们自己的能力中，我们的思想也是如此。} 每一个谦卑而真正虔诚的人都感到这是最真实的。\par

% structure: p0041 source=h2 target=subsection reason=relative-heading-rank; base=h2

\subsection{第二十章——\Person{安波罗削}论天主对人思想的掌控}

当\Person{安波罗削}说这话时，他是在他所写的\WorkTitle{论逃离世界}一文中说的，其中他教导说，逃离这个世界不是靠身体，而是靠心，他指出这除非有天主的帮助，否则无法做到。因为他说：\Quote{我们常听到关于逃离这个世界的谈论，但愿我们的心思能像谈论那么容易一样谨慎而关切；但更糟的是，世俗欲望的诱惑不断潜入，虚妄的倾注占据了心灵；以至于你所想避免的，你却在心中思考和考虑。这很难防范，但不可能摆脱。最后，先知证明这是一件意愿而非成就的事，他说：\InnerQuote{使我的心倾向你的法度，不要倾向贪婪。}\BibleRef{咏 119:36} 因为我们的心与思想不在我们自己的能力中，它们出乎意料地涌出，扰乱我们的心神和灵魂，将我们拉向我们自己预设的不同方向；它们使我们回忆世俗之事，插入时间之事，暗示感官之事，编织诱惑之事，就在我们试图提升心灵的那一刻，我们多半被虚妄的念头充满，被抛到世俗之事中。} 因此，人得以成为天主的子女，不在于人的能力，而在于天主的能力。\BibleRef{若 1:12} 因为他们从他那里领受了赐予人心虔诚思想的那一位，由此产生以爱德行事的信德；\BibleRef{迦 5:6} 对于领受和持守这一恩惠，并坚持不懈直至结束，我们并不足以凭自己思想什么，我们的充足乃出于天主，\BibleRef{格后 3:5} 我们的心和我们的思想都在他的权能之下。\par

% structure: p0043 source=h2 target=subsection reason=relative-heading-rank; base=h2

\subsection{第二十一章 [IX.]——天主不可测度审判的实例}

因此，两个同样受原罪束缚的婴儿，为何一个被接纳而另一个被弃；两个已届成熟之年的恶人，为何这一个如此蒙召，以至于跟随那召叫者，而那一个要么根本没有被召叫，要么没有被以这样的方式召叫——天主的判断是不可探测的。但是，在两个虔诚的人中，为什么一个被赐予坚持到底，而另一个却没有被赐予，天主的判断更是不可探测。然而，对信徒来说，这应当是一个最确定的事实：前者是蒙预定的人，后者不是。\Quote{因为他们如果属于我们，}一位蒙预定的人说，他曾在主怀中饮过这秘密，\Quote{他们必定会留在我们中间。}\BibleRef{若一 2:19}我问，\Quote{他们不属于我们，因为他们如果属于我们，必定会留在我们中间}这句话是什么意思？难道二者不都是天主所造——都是亚当所生——都出于尘土，并由那说\Quote{我创造了所有气息}的那一位赐予，灵魂具有同一本性吗？最后，二者不都蒙了召叫，并跟从了那召叫他们的吗？二者不都从恶人成了义人，并藉着重生的圣洗得以更新吗？但是，若是那毫无疑问知道自己在说什么的人听到这话，他或许会回答并说：这些话是真的。就这一切而言，他们是属于我们的。然而，就某种别的区别而言，他们不属于我们，因为他们如果属于我们，必定会留在我们中间。那么，这是什么区别呢？天主的书卷是敞开的，让我们不要转移视线；神圣的圣经在大声呼喊，让我们聆听。他们不属于他们，因为他们没有\Quote{按着旨意蒙召}；他们没有在创世以前在基督里被拣选；他们没有在祂里面获得基业；他们没有按着那运行万有的祂的旨意被预定。因为如果他们是这样，他们就会属于他们，并且毫无疑问会与他们同在。\par

% structure: p0045 source=h2 target=subsection reason=relative-heading-rank; base=h2

\subsection{第二十二章 说死人将因他们若活着会犯的罪而受审判，是荒谬的。}

暂且不说天主如何可能转变那些厌恶并反对祂信仰之人的意志，并运作他们的心，使他们不屈服于任何逆境，也未被任何试探胜过而离开祂——因为祂也能做使徒所说的事，即不容许他们受试探超过他们所能受的——且先不说这一点，天主预知他们会跌倒，当然能在他们跌倒之前就使他们离开此生。我们是否要回到那个论点，仍然争论说死人甚至为那些天主预知他们若活着会犯的罪而受审判，是多么荒谬？这如此违背基督徒，甚至人的情感，以至于人羞于反驳它。为什么不该说，连福音本身也藉着圣人们如此多的劳苦和苦难而宣讲了，是枉然的，甚至今天仍在枉然宣讲，如果人们即使没有听到福音，也可以根据天主预知他们若听到会有的悖逆或顺从而受审判？\Place{提洛}和\Place{漆冬}不会受罚，即使比那些虽不信却由主基督行了奇事的城市受罚更轻；因为如果那些奇事行在他们中间，他们早已在灰尘和灰烬中悔改了，正如真理之言所宣告的，主耶稣在这些话中向我们显示了预定更深奥的奥秘。\par

% structure: p0047 source=h2 target=subsection reason=relative-heading-rank; base=h2

\subsection{第二十三章 为何在那些将要相信的\Place{提洛}和\Place{漆冬}人民中没有行那些在不信的其他地方所行的奇迹。}

因为如果我们被问，为何这样的奇迹行在那些人当中，他们看见了却不肯相信；而为何不行在那些若看见了就会相信的人当中，我们将如何回答呢？难道我们要说出我在那本我回答了外邦人六个问题的书中所说的，但不损及其他智者可以探究的问题吗？我确实说了这些，如你所知，当时被问及为何基督过了如此长久的时间才降临时：\Quote{在那时那地，他的福音未被宣讲，他预知所有人面对他的宣讲，都会如同他肉身临在时的许多人一样——也就是那些即使他使死人复活也不肯相信他的人。} 此外，在同一本书中稍后，同一问题上，我说：\Quote{有什么稀奇呢，如果基督在以前的世代就知道世界充满了不信者，以致他有理由不愿将他的福音传给那些他预知既不信他的话也不信他的奇迹的人？} 当然我们不能对\Place{提洛}和\Place{漆冬}这么说；在他们身上，我们承认那些神圣的审判涉及预定的那些原因，而我不损及这些隐藏的原因时说过，我当时是在回答诸如此类的问题。当然，指责犹太人的不信是容易的，这不信源于他们的自由意志，因为他们拒绝相信在他们当中所行的如此伟大的奇迹。主斥责他们时也宣告了这一点，他说：\Quote{\Place{苛辣匝因}，你是有祸的！\Place{贝特赛达}，你是有祸的！因为在你们当中所行的异能，如果行在\Place{提洛}和\Place{漆冬}，他们早已披麻蒙灰悔改了。} \BibleRef{路 10:13} 但是，我们能说连\Place{提洛}人和\Place{漆冬}人也会拒绝相信在他们当中所行的如此神奇的作为吗？或者，如果这些作为行了，他们仍然不会相信吗？主亲自为他们作证，如果那些神圣权能的标记行在他们当中，他们会极其谦卑地悔改。然而在审判之日，他们仍要受罚；虽然比那些不相信行在它们当中的神奇作为的城市所受的惩罚要轻。因为主接着说：\Quote{但我告诉你们，在审判之日，\Place{提洛}和\Place{漆冬}所受的要比你们容易忍受。} \BibleRef{玛 11:22} 因此，前者要受更重的惩罚，后者受较轻的；但他们仍要受罚。再者，如果死人还要按他们若活着就会做的事受审判，那么既然这些人如果福音以如此大的奇迹传给他们，他们就会成为信徒，他们当然不应受罚；但他们却要受罚。因此，死人还按他们活着时若福音临到就会做的事受审判，这是假的。如果这是假的，那么就没有理由说，对于因未领受圣洗而死亡的婴儿，他们如此遭遇是理所当然的，因为天主预知他们若活着并且福音传给他们，他们会以不信来听。因此，他们唯独被原罪所束缚，并因此被定罪；我们看到在其他同样情况的人身上，这罪除非藉着天主在重生中的白白恩宠，否则不被赦免；并且，因着他隐秘而公义的审判——因为天主没有不义——有些人甚至在领洗后仍会因邪恶生活而丧亡，却仍被存留在此世直到丧亡，而如果身体的死亡抢先在他们堕入罪恶之前降临，从而拯救他们，他们本不会丧亡。因为没有一个死人按他若不死就会做的善恶之事受审判，否则\Place{提洛}人和\Place{漆冬}人就不会按他们所做的受罚；而是按他们若那些福音的异能行在他们当中就会做的事，他们本会藉着极大的悔改和基督的信德而得救恩。\par

% structure: p0049 source=h2 target=subsection reason=relative-heading-rank; base=h2

\subsection{第二十四章 [X.]——可能会被反对：\Place{提洛}和\Place{漆冬}的人民如果听了，可能会相信，随后又从他们的信仰中堕落了。}

一位声誉不凡的天主教论辩者如此解释这段福音经文，说：主预知提洛人和漆冬人后来会背弃信仰，尽管他们曾相信在他们中所行的奇迹；祂出于仁慈没有在那里行那些奇迹，因为如果他们背弃了曾经持守的信仰，就会比从未持守信仰遭受更严厉的惩罚。关于这位博学且非常敏锐之人的意见，我现在为何还要关心去说还有什么合理的问题可问呢？因为即使这个意见也为我所追求的目的服务。因为如果主出于仁慈没有在他们中间行大能的事，因为通过这些事他们可能会成为信徒，以免当他们后来成为不信者时受到更严厉的惩罚——正如祂预知他们会那样——这便充分而清楚地表明，没有人会因祂预知他会犯的罪而受审判，如果他没有以某种方式被帮助不去犯那些罪；正如如果那个意见是真实的，那么基督就被说成是来帮助提洛人和漆冬人，祂宁愿他们根本不来到信仰中，也不愿他们因更大的邪恶背弃信仰，正如祂预见如果他们来到信仰就会背弃那样。尽管如果有人说：\Quote{为什么没有安排使他们反而相信，并赐予他们这份恩赐，使他们在背弃信仰之前离开此生？}我不知道能作何回答。因为那说对那些将要背弃信仰的人，作为恩惠，他们会不被赐予开始拥有那更严重的背弃行为的人，充分表明一个人不会因祂预知他会作恶的事而受审判，如果藉某种恩惠他可以避免行恶。因此，对于被接去的人，免得邪恶改变他的心思，这也是有益的。但为什么这益处没有赐给提洛人和漆冬人，使他们相信并被接去，免得邪恶改变他们的心思——那喜欢以这种方式解决上述问题的人或许可以回答；但就我所讨论的而言，我认为这已经足够了：即使按照那个意见，人也显明不会因他们未曾行的事而受审判，即使它们被预知为肯定会行。然而，正如我所说，我们甚至当耻于反驳这个意见，因为它假定罪会在死去或已死的人身上受罚，只因为他们被预知如果活着一定会行那些罪；免得我们似乎也认为这个意见有些重要性，尽管我们宁愿用论证抑制它，而不是默然放过。\par

% structure: p0051 source=h2 target=subsection reason=relative-heading-rank; base=h2

\subsection{第25章 [XI.]——天主的道路，在仁慈和审判中，都是难寻的。}

因此，正如宗徒所说：\BibleQuote{所以，不在乎那定意的，也不在乎那奔跑的，只在乎那施怜悯的天主。}\BibleRef{罗 9:16}祂随意帮助那些婴儿，虽然他们既不定意也不奔跑，因为祂在创世以前在基督内拣选了他们，要白白赐给他们恩宠——也就是说，没有任何他们的功劳，无论是信德还是善行，预先存在；祂也不帮助那些较成熟的人，尽管祂预见了如果奇迹行在他们中间他们就会相信，因为祂不愿帮助他们，因为在祂的预定中，祂隐秘却又公义地对他们另有决定。因为\BibleQuote{在天主，没有不义；}\BibleRef{罗 9:14}但\Quote{祂的判决是何等不可测度，祂的道路是何等难寻！上主的一切道路都是仁慈和真理。}因此，仁慈是难寻的，祂随意怜悯谁，没有他自身的功劳预先存在；真理是不可测度的，祂随意使谁心硬，尽管他的功劳可能预先存在，但那些功劳多半是与祂所怜悯之人共有的。如两个双胞胎，一个被取去，一个被留下，结局不平等，而功过相同，其中一人如此被天主的伟大良善拯救，而另一人被天主的公义定罪——难道天主有不义吗？断乎不可！但祂的道路是难寻的。因此，让我们毫不迟疑地相信祂在得救者身上的仁慈，以及在受罚者身上的真理；不要试图探查那不可测度的，也不要追寻那不能寻见的。因为从婴孩和吃奶者口中，祂完善了赞美，以至我们在那些得救者身上所见到的——他们的得救没有他们任何善功预先存在——以及在那些受罚者身上所见到的——他们的定罪只有原罪预先存在，原罪是两者共有的——我们绝不畏缩承认这也发生在成年人身上，也就是说，我们不认为恩宠是照各人的功劳赐给的，也不认为任何人受罚不是因为自己的功劳，无论是得救者和受罚者有相似的功过，还是具有不平等的邪恶程度；这样，那自以为站得稳的，要小心免得跌倒；那夸耀的，要夸耀主。\par

% structure: p0053 source=h2 target=subsection reason=relative-heading-rank; base=h2

\subsection{第26章——摩尼教徒不接受旧约的全部书卷，新约也只接受他们所选的部分。}

但是，为何那些毫不犹豫地反对白拉奇派、肯定原罪——这原罪由一人进入世界，且从一人众人归于定罪——的人，如你所写，说\Quote{婴孩的案例不可引用}，又\Quote{作为长者的例子}呢？摩尼派也不接受这一点，他们不仅拒绝所有旧约圣经的权威，甚至接受新约圣经的方式也是如此，即每个人凭自己所谓的特权，更确切地说是凭自己的亵渎，取其所好，弃其所恶。我曾在我的\WorkTitle{论自由意志}著作中针对他们进行论述，他们以为据此有了反对我的理由。我不愿坦率处理那些非常繁难的问题，以免我的作品变得太长，因为在对抗这种乖僻之人时，我无法得到圣经权威的帮助。然而我能够——事实上我确实做到了，无论神圣见证是否真实，鉴于我并未明确将其引入论证——通过某种推理得出结论：天主在一切事上都应受赞美，而无需相信他们想要我们相信的，即存在着两种同永恒的、混乱的善恶实体。\par

% structure: p0055 source=h2 target=subsection reason=relative-heading-rank; base=h2

\subsection{第二十七章——引用\WorkTitle{更正篇}}

最后，在我尚未读完的\WorkTitle{更正篇}第一卷中，当我重新审视那几卷书，即\WorkTitle{论自由意志}时，我这样说：\Quote{在这些书中，}我说，\Quote{许多问题如此讨论，以致于出现了一些我无法阐明或需要同时进行长篇讨论的问题时，它们被如此推延，以至于从这些问题的一面或多面来看，虽然最符合真理的尚未显现，但我的推理仍可为此作结，即无论其中哪个为真，天主都可能被相信，甚至被证明是值得赞美的。因为那场讨论是为了那些否认恶的来源出自自由意志选择、并坚持认为天主——若祂是万有的创造者——应受责备的人而进行的；他们按照其不敬之心的错误（因为他们是摩尼派），想要引入某种与天主同永恒的、不变的恶的本性。}稍后，在另一处我又说：\Quote{于是便说，从这最公义地加于罪人的苦难中，天主的恩宠拯救人，因为人自愿地，即藉自由意志，能够堕落，却不能也重新站起。每一人自出生之初就经历的愚昧与困难，属于这公义定罪的苦难，没有人能脱离这恶，除非藉天主的恩宠。白拉奇派不愿这苦难源于公义的定罪，因为他们否认原罪；然而，即使愚昧与困难是人的自然起点，天主也不应因此受责备，反而应受赞美，正如我在同书第三卷中所论证的。这一论证应视为针对摩尼派，他们不接受旧约圣经（其中记载了原罪），并且他们争辩说，凡从旧约读到的在宗徒书信中的内容，都是圣经篡改者以可恶的无耻插入的，以为宗徒并未说过。但针对白拉奇派，必须坚持圣经所赞扬的内容，因为他们自称接受圣经。}这些是我在\WorkTitle{更正篇}第一卷中，重新审视\WorkTitle{论自由意志}诸书时所说的话。诚然，我在那里关于这些书所说的话并未尽录于此，还有许多其他的话，我认为若将其插入此信给你会显得冗长且不必要；我想当你读到全部内容时，也会如此判断。因此，尽管我在\WorkTitle{论自由意志}第三卷中如此论证婴孩之事，以致于即使白拉奇派所说是真的——即愚昧和困难（无人不带它们出生）是人性固有的要素，而非惩罚——摩尼派仍会被驳倒，因为他们坚持善恶两种本性同永恒。那么，难道应质疑或放弃天主教会针对白拉奇派所持守的信仰吗？教会坚称有原罪，其罪债由生育而感染，必须藉再生而赦免。若他们与我们一同承认这一点，以便我们在白拉奇派这个问题上立即摧毁错误，那么他们为何认为必须怀疑天主能拯救甚至那些藉圣洗圣事领受祂恩宠的婴孩，使他们脱离黑暗的权势，而迁入祂爱子的国度？\BibleRef{哥 1:13}因此，既然祂将这恩宠赐予一些人而不赐予另一些人，他们为何不向主歌颂祂的慈爱与审判？然而，为何赐给这些人而非那些人？谁知道主的心意？谁能洞察不可测度之事？谁能追寻那无法追查之事？\par

% structure: p0057 source=h2 target=subsection reason=relative-heading-rank; base=h2

\subsection{第二十八章（第十二节）——天主的良善与正义在万物中彰显}

因此，已经确定：天主的恩宠并非按照领受者的功过赐予，而是按照祂旨意的喜悦，为赞美和光荣祂自己的恩宠；这样，夸耀的人绝不可夸耀自己，而应夸耀主。祂愿意给谁就给谁，因为祂是仁慈的；但如果祂不给，祂是公义的。祂不给那些祂不愿意给的人，为将祂荣耀的丰富彰显于怜悯的器皿上。\BibleRef{罗 9:23} 因为，通过给一些人他们不配得的，祂确实愿意祂的恩宠是白白的，因而是真正的恩宠；通过不给所有人，祂显明了所有人都应得的。祂在恩待一些人时是良善的，在惩罚另一些人时是公义的；祂在所有方面都是良善的，因为给予应得的总是好的；祂在所有方面都是公义的，因为不按功过赐予而不冤枉任何人。\par

% structure: p0059 source=h2 target=subsection reason=relative-heading-rank; base=h2

\subsection{第二十九章——即使没有原罪，如同白拉奇派所坚持的，天主的真正恩宠仍可维护。}

但天主的恩宠，即没有功过的真正恩宠，仍得到维护，因为即使按照白拉奇派的观点，婴儿受洗时并未被从黑暗的权势中救出，因为他们被认为没有犯任何罪，而只是被转入天主的国：因为即便如此，没有任何善功，天国赐给那些被赐予的人；也没有任何恶行，天国就不赐给那些没有被赐予的人。我们习惯于这样反驳白拉奇派，当他们反对我们说我们把天主的恩宠归于命运，因为我们说恩宠的赐予不取决于我们的功过。因为他们自己在婴儿的事上倒把天主的恩宠归于命运，如果他们声称没有功过就是命运。当然，即使按照白拉奇派自己的观点，在婴儿身上找不到任何功过，使得一些婴儿得以进入天国，而另一些婴儿被排斥在天国之外。但如今，正如为了显示天主的恩宠并非按我们的功过赐予，我宁愿根据两种意见来维护这个真理——既根据我们自己的意见，即认为婴儿受原罪的束缚，也根据白拉奇派的意见，即否认原罪的存在，但我不能因此怀疑婴儿拥有那拯救祂的子民脱离罪恶的祂所能宽恕的东西：同样，在\WorkTitle{论自由意志}第三卷中，我根据两种观点驳斥了摩尼教徒，即无知和困难究竟是惩罚还是无人天生就有的自然元素；但我坚持其中一种观点。此外，在那里我已足够清楚地说明，那不是人被造时的本性，而是他被定罪后的惩罚。\par

% structure: p0061 source=h2 target=subsection reason=relative-heading-rank; base=h2

\subsection{第三十章——\Person{奥斯定}主张在知识上成长的权利。}

因此，有人从我那本旧书中引经据典地向我规定，我不可按照在婴儿问题上应有的方式来论证；也不可借此用显明真理的光照说服我的对手，即天主的恩宠并非按人的功过赐予。因为，如果我以平信徒的身份开始写作\WorkTitle{论自由意志}，并以司铎的身份完成时，我仍对未重生的婴儿的定罪和已重生的婴儿的得救存有疑虑，那么，我想，没有人会如此不公和嫉妒，竟阻止我前进，并断定我必须停留在这不确定性中。但更正确的理解是：应当相信我在那件事上没有怀疑，因为我的反驳目标似乎是如此被驳倒的，即婴儿身上是否有原罪的惩罚，按真理来说，或没有，如某些错误之人所想，无论如何都不应该相信像摩尼教徒的谬误所引入的那种善恶二性的混淆。因此，我们切不可如此放弃婴儿的案件，以至于自认为不确定：他们在基督内重生后，若死于婴幼时期，是否进入永生；而未重生者，则进入第二次死亡。因为那所记载的：\BibleQuote{罪借着一人进入了世界，死亡借着罪也进入了世界；这样死亡就传给了众人，}\BibleRef{罗 5:12} 除非如此理解，否则不能正确明白；也无人，无论大小，能从那最公义地回报罪的永死中救出任何人，唯有祂，为了赦免我们的罪，包括原罪和本罪，自己没有任何原罪或本罪而死。但为什么是有些人而不是另一些人？我们一再地说，并不退缩：\BibleQuote{人啊！你是谁，竟敢反驳天主？}\BibleRef{罗 9:20} \BibleQuote{他的判断是多么不可测量！他的道路是多么不可探察！}\BibleRef{罗 11:33} 我们再添上：\BibleQuote{不要探寻高于你的事，也不要求索超过你能力的事。}\BibleRef{德 3:21}\par

% structure: p0063 source=h2 target=subsection reason=relative-heading-rank; base=h2

\subsection{第三十一章——婴儿不因他们被预知若活着可能做的事而受审判。}

因为你们看，亲爱的诸位，说婴孩去世时应按他们若活着将做的事——这些事已被预知——来受审判，这是多么荒谬，多么远离健全的信德与真理的真诚。因为这种意见——当然，一切人的情感，即使仅基于很少的理由，尤其是基督徒的情感，都会对它感到厌恶——他们被迫走向它，因为他们选择如此远离白拉奇人的错误，却仍认为必须相信并更在辩论中宣称：天主的恩宠，是藉着我们的主耶稣基督而来的，唯独藉着这恩宠，在首先的人（我们都在他内跌倒了）跌倒之后，我们才得到援助，而这恩宠是按照我们的功过赐予的。这个信念，连白拉奇本人，在东方主教们作为审判者面前，也因害怕自己被定罪而弃绝了。如果这一点不适用于那些已死之人的善或恶行（这些行若他们活着本会去做），因此也不适用于任何行为，以及那些甚至在天主的预知中也永不存在的行——如果这一点因此而不被承认，而你们看到它在多大的错误中被陈述，那么除去争论的黑暗后，除了我们承认天主的恩宠不是按照我们的功过赐予的（这是天主教会反对白拉奇异端所辩护的立场），并且我们在婴孩身上更明显地看到这一点之外，还有什么可留下的呢？因为天主不是被命运强迫去帮助这些婴孩而不帮助那些——因为两者的处境相同。或者我们会认为，在婴孩的事上，人类事务不是由天主眷顾管理，而是由偶然的机会管理，当理性灵魂要么被定罪要么被拯救时，尽管事实是，\BibleQuote{没有一只麻雀会不经过我们在天之父的意愿而落在地上}（见：\BibleRef{玛 10:29}）？或者我们必须如此归咎于父母的疏忽，以致婴孩未受圣洗而死去，天上的审判与它无关，仿佛那些这样死得不好的人是自己选择了他们出生于其中的疏忽父母？当婴孩在可能从圣洗圣事获益之前就断气时，我该说什么呢？因为常常当父母渴望、施行圣洗的职员准备好给婴孩付洗时，它仍然没有被赐予，因为天主不愿意；因为没有让它在世上多留片刻以接受圣洗。再者，有时圣洗能帮助不信者的子女，使他们不致进入丧亡，却不能帮助信徒的子女？在这种情况下，肯定显示出天主不看人的情面；否则祂宁可拯救祂敬拜者的子女，而不是祂仇敌的子女。\par

% structure: p0065 source=h2 target=subsection reason=relative-heading-rank; base=h2

\subsection{第三十二章（第十三）——天主自由旨意的不可测度}

但是，既然我们现在正讨论恒心到底的恩赐，为什么援助赐予那即将去世却未受洗的人，而不赐予那已受洗却即将跌倒的人，使他能够在此之前死去？除非我们仍然听取那荒谬说法：在跌倒前死去对任何人没有益处，因为他将按照天主预知他若活着会做的那些行为受审判。谁有耐心听这种如此猛烈反对信德健全的乖谬？谁能忍受？然而，那些不承认天主的恩宠不是按我们的功过赐予的人，被迫这样说。但是，那些不愿意说任何已死之人是按照天主预知他若活着会做的事受审判的人——考虑到这说法是何等明显的虚假与何等大的荒谬——就再也没有理由说教会曾在白拉奇人身上定罪并让白拉奇自己定罪的（即天主的恩宠是按照我们的功过赐予的），当他们看见一些未重生的婴孩从今生被取去归入永死，而另一些重生的婴孩归入永生；并且那些重生的婴孩中，有些恒心到底离世，有些却在今生被存留直到他们跌倒（这些人若是在跌倒前离世，肯定不会跌倒）；又有些跌倒，但至死未离去，直到他们回头（这些人若是在回头前离去，肯定会丧亡）。\par

% structure: p0067 source=h2 target=subsection reason=relative-heading-rank; base=h2

\subsection{第三十三章——天主按其旨意赐予起始之恩宠与恒心之恩宠}

凡此种种，均足以充分明示：天主的恩宠——既开启人的信德，又使之恒心至终——并非按我们的功过而赐予，乃是按祂那至为隐秘、同时至为公义、明智而仁慈的旨意而赐予；因为祂所预定的人，也召叫了他们，\BibleRef{罗 8:30} 藉着那召叫，经上论及说：\Quote{天主的恩赐和召叫是不会撤回的。}\BibleRef{罗 11:29} 对于这召叫，无人能在有生之年确凿无疑地断言自己属于其中，直到他离世为止；但在人的今生，即在地上受考验的时日，\BibleRef{约 7:1} 那自以为站得稳的，须小心，免得跌倒。\BibleRef{格前 10:12} 因为（如我先前所言）那些不会恒心到底的人，凭借天主极具先见的旨意，与那些将恒心到底的人混杂在一起，原因在于要我们学习不可心高气傲，却要俯就卑微，并且要\Quote{怀着恐惧战栗，成就自己的救恩；因为是天主在你们内工作，使你们愿意，并使你们行事，为成就祂的美意。}\BibleRef{斐 2:12} 因此，我们愿意，但天主在我们内工作使我们愿意。因此，我们行事，但天主在我们内工作使我们行事，为成就祂的美意。这对我们有益，要我们相信并如此宣讲——这是虔诚的，这是真实的，要使我们的宣认谦卑顺服，并将一切归于天主。我们思考才相信，我们思考才说话，我们思考才行事，做我们所做的一切；\BibleRef{格后 3:5} 但是，就虔敬之路和真诚敬拜天主而言，我们凭自己不能思考什么，我们的能力来自天主。因为\Quote{我们的心与心思不在我们自己权下}；说出此话的同一位安波罗削也说：\Quote{然而，谁如此有福，能常在心中向上提升？这若没有神圣的帮助，又如何可能？断乎不能。最后，他说，同一圣经在上文证实：\InnerQuote{上主，那以祢为助佑的人是有福的，他的心上升。}}无疑，安波罗削不仅通过阅读圣经而能说出这话，而且，正如我们理当相信，作为这样的人，他在自己心中也体验到了。因此，正如在信徒的圣事中所说的，我们应当向主举心向上，这是天主的恩赐；为此恩赐，被劝勉的人在这话之后，应藉司铎之口向我们的主天主自己感恩；他们回应说\Quote{这样做是合宜且正当的}。因为，我们的心不在我们自己权下，而是藉神圣的帮助被提升，从而上升并认识那上面的——\BibleRef{哥 3:1} 那里有基督坐在天主的右边，而非地上的事——对于如此伟大的恩赐，除了成就此事的主天主，我们还能向谁感恩呢？祂以如此巨大的仁慈，藉着将我们从这世界的深渊中解救出来而拣选了我们，并在世界奠基之前就预定了我们？\par

% structure: p0069 source=h2 target=subsection reason=relative-heading-rank; base=h2

\subsection{第三十四章 [XIV.]——预定的教义并不反对宣讲的益处。}

但有人说：\Quote{预定的定义与宣讲的益处相违背}——仿佛这与宗徒的宣讲相违背似的！那位外邦人的导师岂非屡次以信德和真理，既推崇预定，又不停止宣讲天主圣言吗？因为他曾说：\Quote{是天主在你们内工作，使你们愿意，并使你们行事，为成就祂的美意}\BibleRef{斐 2:13}，难道他不是也劝勉我们要愿意并行事，以中悦天主吗？或者因为他曾说：\Quote{那在你们身上开始了美好工作的，必予以完成，直到耶稣基督的日子}\BibleRef{斐 1:6}，难道他就因此停止劝人开始并恒心至终吗？无疑，主自己命令人相信，并说：\Quote{你们要信天主，也要信我}：\BibleRef{若 14:1} 然而，祂的见解并不因此为假，祂的定义也并非虚设，当祂说：\Quote{除非父赐给他，没有人能到我这里来}——即没有人信我——\Quote{除非父赐给他。}\BibleRef{若 6:66} 再者，因为这定义是真的，先前的命令也非徒然。那么，为何我们认为预定的定义对宣讲、命令、劝勉、责备——圣经屡次重述所有这些事——是无用的呢？既然同一位圣经也称赞这道理？\par

% structure: p0071 source=h2 target=subsection reason=relative-heading-rank; base=h2

\subsection{第三十五章——什么是预定。}

难道有人敢说天主没有预知祂要赐予谁相信，或把谁赐给祂的圣子，使祂不丧失其中任何一个吗？\BibleRef{若 18:9} 诚然，若祂预知了这些事，祂也同样预知祂的仁慈行为，借此祂屈尊拯救我们。这就是圣人的预定——并非别的；即天主的仁慈之预知和预备，借此他们确实被拯救，无论他们是谁。但其余的人被公义的神圣审判留在何处，岂非留在毁灭的众人中，就如\Place{提洛}人和\Place{漆冬}人被留下一样？再者，谁若看见了基督的奇妙奇迹，也许就会相信。但既然没有赐给他们相信，相信的方法也被拒绝给他们。由此可见，有些人在其理解力本身内有一种天生的智能恩赐，借此他们可以被感动而走向信德，只要他们听见适合他们心智的言语或看见标记；然而，如果在天主更高的审判中，他们没有被恩宠的预定从丧亡的众人中分别出来，那么那些神圣的言语或事迹也不会应用于他们，如果他们只是听见或看见这些事，他们或许会相信。再者，犹太人也留在同一毁灭的众人中，因为他们不能相信在他们眼前所行的如此伟大而显赫的异能。因为福音并没有对为什么他们不能相信保持沉默，它说：\Quote{祂虽然在他们面前行了这样大的奇迹，他们却不信祂；这是为应验依撒意亚先知所说的话：\BibleRef{依 53:1}\InnerQuote{主啊，谁相信了我们的报告？主的膀臂向谁启示了呢？}因此，他们不能相信，因为依撒意亚又说：\BibleRef{依 6:10}\InnerQuote{祂弄瞎了他们的眼睛，硬了他们的心，免得他们眼睛看见，心里明白，回转过来，我就医治他们。}}因此，提洛人和漆冬人的眼睛并没有如此被弄瞎，他们的心也没有如此变硬，因为如果他们看见了如同犹太人看见的那样的大能作为，他们就会相信。但对他们而言，能够相信并没有益处，因为他们没有被祂所预定，祂的判断是多么不可测度，祂的道路是多么难寻。但如果他们如此被预定，以致天主光照那些瞎眼，并愿意从那些硬心者那里除去石心，那么不能相信也不会成为他们的障碍。但主关于提洛人和漆冬人所说的话，也许可以另一种方式理解：即除非是赐给他的，没有人来到基督那里，而那是赐给那些在世界创立以前在祂内蒙拣选的人，凡以心灵之耳、而非肉体的聋耳聆听神圣言语的人，都毫无疑问地承认这一点；然而，这预定——甚至由福音的话也足够清楚地显明出来——并没有阻止主既关于开始说了我前文稍提的话：\Quote{你们信天主，也信我罢。}又关于恒心说了：\Quote{人应当时常祈祷，不可灰心。}\BibleRef{路 18:1} 因为凡听见这些事并实行它们的人，是那些被赐予的人；但那些没有被赐予的人，无论他们听见或不听见，他们都不实行。因为祂说：\Quote{你们，}祂说，\Quote{是给你们知道天国的奥秘，但不是给他们。}\BibleRef{玛 13:11} 这些事中，一个指仁慈，另一个指审判，我们的灵魂向祂呼喊：\Quote{我要向祢歌唱仁慈与公义，主。}\par

% structure: p0073 source=h2 target=subsection reason=relative-heading-rank; base=h2

\subsection{第36章——福音宣讲与预定宣讲：同一信息的两个部分}

因此，预定之宣讲不应阻碍恒心与进步的信德之宣讲；这样，凡被赐予应当服从的人，就能听到必要的事。因为若没有宣讲者，他们怎能听到呢？同样，继续至终的进步信德之宣讲也不应阻碍预定之宣讲，好使那信实而服从地生活的人，不致因那服从本身而自高自大，仿佛那是他自己未曾接受的功劳；而是使他夸耀的，夸耀于主。因为\Quote{我们什么都不应夸耀，因为没有什么是我们自己的。}\Person{西彼廉}最忠实地看到并最无畏地解释了这一点，因此他断言预定是极为确定的。因为若我们什么都不应夸耀，既然没有什么是我们自己的，我们当然不应夸耀那最持久的服从。它也不能被称为我们自己的，仿佛它不是从上面赐给我们的。因此，它是天主的恩赐，所有基督徒都承认，天主预知祂会赐给祂的子民，就是那些蒙召的人，论到这召叫，经上说：\Quote{天主的恩赐和召叫是不可撤回的。}\BibleRef{罗 11:29} 这就是我们忠信而谦逊地宣讲的预定。然而，同一位导师和实行者，他既信了基督，又最恒心地生活在神圣的服从里，甚至为基督受苦，并没有因此停止宣讲福音，劝勉信德和虔诚的行为，以及那至终的恒心，因为他曾说：\Quote{我们什么都不应夸耀，因为没有什么是我们自己的；}在这里，他毫不含糊地宣告了天主的真恩宠，即不是因我们的功德而赐予的恩宠；既然天主预知祂会赐予它，预定就毫无疑问地被西彼廉的这些话所宣告；如果这没有阻止西彼廉宣讲服从，当然也不应阻止我们。\par

% structure: p0075 source=h2 target=subsection reason=relative-heading-rank; base=h2

\subsection{第37章——听从的耳朵即是服从的意愿}

虽然我们因此说服从是天主的恩赐，但我们仍然劝勉人去服从。但那些顺从地听从真理劝勉的人，得到了天主本身的恩赐——即顺从地听从；而那些不如此听的人则没有得到。因为说这话的不仅是某个人，而是基督说：\BibleQuote{除非蒙我父赐给他，谁也不能到我这里来} \BibleRef{若 6:66}；以及\BibleQuote{天国的奥秘是给你们知道的，但不是给他们知道的} \BibleRef{玛 13:11}。关于贞洁，祂说：\BibleQuote{这话不是人人都能领悟，惟独那些得了恩赐的人才能领悟} \BibleRef{玛 19:11}。当宗徒劝勉已婚者保持婚姻贞洁时，他说：\BibleQuote{我切愿众人都像我一样；但每人各自有从天主得来的恩赐，有人这样，有人那样} \BibleRef{格前 7:7}；这里他清楚表明不仅贞洁是天主的恩赐，连已婚者的贞洁也是。虽然这些是真实的，我们仍尽力劝勉，按我们每一个人所能劝勉的，因为这也是祂的恩赐，我们自身和我们的言语都在祂手中。因此宗徒也说：\BibleQuote{按照天主赐给我的恩宠，我好像一个明智的建筑师，奠立了根基} \BibleRef{格前 3:10}。在另一处他说：\BibleQuote{正如主赐给了每人，我栽种，阿颇罗浇灌，然而使之生长的乃是天主。所以，栽种的不算什么，浇灌的也不算甚么，惟独使之生长的天主才算甚么} \BibleRef{格前 3:5}。因此，只有那得了这恩赐的人才能正确宣讲和劝勉，同样，那顺从地听从正确劝勉和宣讲的人，正是得了这恩赐的人。由此主说了那句话：当祂对那些肉耳本是开启的人说话时，仍告诉他们：\BibleQuote{有耳听的，就听吧} \BibleRef{路 8:8}；这无疑祂知道并非所有人都有。而他们从谁那里得到耳朵，无论谁是那有耳之人，主亲自显示说：\BibleQuote{我要赐给他们一颗认识我的心，和能听的耳朵} \BibleRef{巴 2:31}。因此，有耳朵本身就是顺从的恩赐，以便那些有此恩赐的人能到祂那里去——\Quote{除非蒙祂父赐给他，没有人能到祂那里}。因此我们劝勉和宣讲，但那些有耳顺从听的人顺从地听我们，而在那些没有耳的人身上，应验了所写的：他们听了却不听见——即用肉耳听了，却没有心中的赞同。但为什么这些人有耳可听，那些人却没有——即为什么有些人由圣父赐予他们到圣子那里，而有些人却没有——谁知道主的心意？谁作过祂的参谋？或者，你这个人哪，你是谁，竟敢向天主抗议？难道那明显的事要被否认，因为那隐藏的事不能被理解吗？我们是否要说，我们所见为是如此的事，其实不然，因为我们无法发现它为何如此？\par

% structure: p0077 source=h2 target=subsection reason=relative-heading-rank; base=h2

\subsection{第三十八章 [XV.]——反对预定宣讲的同样异议也可用以反对预定}

但他们说，如你所说：\Quote{如果人在教会的集会中对众多听者这样说：关于预定，天主旨意的确定含义是这样，你们中有些人将接受服从的意愿，并从不信进入信德，或者接受恒心并居住于信德；但另一些人逗留在罪恶的喜乐中，尚未兴起，因为慈悲恩宠的帮助尚未确实振奋你们。然而，如果有任何人因祂的恩宠被预定为蒙选者，尚未蒙召，你们将接受那恩宠，使你们得以愿意并被选；并且如果有任何人服从，如果你们被预定被弃绝，服从的力量将从你们那里被撤回，以使你们停止服从；那么没有人能被责备的激励所唤起。} 虽然可以说这些话，但它们不应阻止我们承认天主的真恩宠——即不因我们的功绩而赐予我们的恩宠——以及按此承认圣人的预定，正如我们不因承认天主的预知而受阻，尽管有人可能对百姓这样说：\Quote{无论你现在是以义人或恶人的身份生活，你将不久成为主所预知的你将会是的——或是好的，如果祂预知你是好的；或是坏的，如果祂预知你是坏的。} 因为如果因听到此话，有些人转为麻木和懒惰，从奋斗转向放纵，沉溺于自己的欲望，难道就要认为关于天主预知所说的话是假的吗？如果天主预知他们将成为善人，无论他们现在陷入多深的邪恶，岂不仍会成为善人？如果祂预知他们为恶人，无论现在在他们身上看出多少善，岂不仍会成为恶人？在我们的隐修院里有一个人，当弟兄们因他做一些不该做的事、不做一些该做的事而责备他时，他回答说：\Quote{无论我现在如何，我将如天主所预知的那样成为。} 这人确实说了真话，但并没有从这真理中受益而向善，反而如此在邪恶中前进，以至于离开了隐修院团体，变成一只返回呕吐物的狗；然而，他将来会变成什么还不确定。那么，为了这类灵魂的缘故，关于天主预知所说的真理是应当被否认还是被保留不提呢？——例如，当如果不讲，就会引起其他错误的时候？\par

% structure: p0079 source=h2 target=subsection reason=relative-heading-rank; base=h2

\subsection{第三十九章 [XVI]——祈祷与劝勉}

此外，有些人根本不祈祷，或祈祷冷淡，因为从主的话中，他们得知天主在我们求祂之前，就知道我们需要什么。难道要放弃这话的真理，或者认为因着这类人，这话应从福音中删除吗？不，既然明显天主预备了一些东西甚至赐给那些不为它们祈祷的人，例如信德的开端，而另一些东西只赐给那些为它们祈祷的人，例如坚持到底的恒心，那么肯定，那认为自己从自身拥有后者的人，就不会为拥有它而祈祷。因此我们必须小心，免得当我们担心劝诫变得冷淡时，祈祷被窒息，傲慢被激发。\par

% structure: p0081 source=h2 target=subsection reason=relative-heading-rank; base=h2

\subsection{第四十章．何时应说出真理，何时应缄默。}

因此，让真理被说出，尤其是当任何问题促使我们宣告它时；让能接受的人接受它，免得，当我们因那些不能接受的人而沉默时，那些能接受真理从而避免错谬的人，不仅被剥夺真理，反被错谬俘虏。因为有些真理应该因那些不能领会的人而保留，是容易的，不，且是有益的。因为主的那句话从何而来：\BibleQuote{我还有许多事要告诉你们，但你们现在不能担负。}？\BibleRef{若16:12}以及宗徒的那句话：\BibleQuote{我不能向你们讲论，如同对属神的人讲论，只能对属肉体的人，如同对基督内的婴孩讲论。我给你们喝的是奶，并非食物，因为那时你们还不能吃，就是如今你们也不能吃。}？\BibleRef{格前3:1}虽然，从某种意义上说，所说的话既可以是给婴孩的奶，也可以是给成年人的肉。正如\BibleQuote{在起初已有圣言，圣言与天主同在，圣言就是天主，}\BibleRef{若1:1}哪个基督徒能隐藏它呢？谁能接受它呢？或者在纯正道理中有什么比这更全面的呢？然而，这并不对婴孩或成年人隐藏，也不被成熟的人向婴孩隐藏。但保留真理的理由是一回事，说出真理的必要是另一回事。要探究或列出所有保留真理的理由是繁琐的工作；然而，其中有一个理由——免得我们使那些不理解的人变得更糟，同时又希望使那些理解的人更有学问；尽管这些后者在我们保留任何这类事时并不会变得更有学问，但也不会变得更糟。然而，当一个真理的性质是如此，以致不能接受的人因我们说出它而变得更糟，而能接受的人因我们沉默而变得更糟，我们想该怎么办呢？难道我们不该说出真理，使能接受的人得以接受，而不是保持沉默，以致不仅两人都不能接受，而且甚至那更聪明的人自己反而变得更糟吗？因为如果他听到并接受了，借着他，许多人也能学习。因为他越有能力学习，就越适合教导他人。恩宠的仇敌在各种方式上逼迫和催促，使我们相信恩宠是按我们的功德赐给的，这样恩宠就不再是恩宠了；我们难道不愿说出我们能用圣经的见证所说的话吗？我们难道真的害怕因我们的说话而冒犯那不能接受真理的人吗？我们难道不怕因我们的沉默，那能接受真理的人可能陷入错谬吗？\par

% structure: p0083 source=h2 target=subsection reason=relative-heading-rank; base=h2

\subsection{第四十一章．预定界定为唯一天主在其预知中安排事件。}

因为，必须按圣经明确宣告的方式和程度宣讲预定，使被预定者身上天主的恩赐和召叫不会反悔；否则就必须承认天主的恩宠是按照我们的功绩赐予的，——这是白拉奇异端者的意见；尽管他们的这种意见，正如我已多次说过的，可以在东方主教会议纪要中读到，是被白拉奇本人亲口谴责的。此外，我所谈论的那些人，仅就他们脱离了白拉奇异端者的邪曲异端而言，虽然他们不愿承认那些因天主恩宠而成为服从并如此持守的人是被预定的，但他们仍然承认，恩宠先于它所赐予之人的意志；当然，这样说是为了让恩宠不被认为是白白赐予的，正如真理所宣告的，而是按照先行意志的功绩，正如白拉奇的错误与真理相悖所说的那样。因此，恩宠也先于信德；否则，如果信德先于恩宠，那么毫无疑问意志也先于恩宠，因为信德不能没有意志。但如果恩宠先于信德，因为它先于意志，那么它当然先于一切服从；它也先于爱德，唯有藉着爱德，天主才得到真正而愉快的服从。恩宠在它所赐予的人身上，并在他之内先行于这一切，成就这一切。[XVII.] 在这些恩惠中，还有恒心至终，如果主不藉着他的恩宠在所听祈祷的人身上成就它，那么天天向主祈求恒心至终便是徒劳的。请看，否认恒心至终甚至到今生结束是天主的恩赐，这与真理是何等相悖；因为主自己在祂所愿意的时候结束此生，如果祂在迫在眉睫的堕落之前结束生命，祂就使人恒心至终。但更奇妙且更显明于信徒的是天主美善的慷慨，就是这恩宠甚至赐给婴孩，尽管在那个年龄没有可以赐予的服从。因此，天主将祂的恩赐赐给谁，毫无疑问祂预知祂将赐予他们，并在祂的预知中为他们预备了这些恩赐。因此，祂所预定的人，祂也以那种召叫召叫了他们，这种召叫我常乐于提及，关于这召叫经上说：\BibleQuote{天主的恩赐和召叫是不会反悔的。}\BibleRef{罗 11:24} 因为祂在不可欺骗、不可改变的预知中对其未来工作的安排是绝对的，并且不是别的，正是预定。但是，正如天主预知为贞洁的人，尽管他可能认为此事不确定，仍如此行动以成为贞洁；同样，天主预定为贞洁的人，尽管他可能认为此事不确定，却不会因为他听到自己将因天主的恩赐而成为将是的样式，就不努力去成为贞洁。反而，他的爱德喜乐，并且他不自高自大，仿佛他没有领受一样。因此，他不仅不被预定的宣讲阻碍这工作，反而得到助益，以致他若夸耀，就夸耀于主。\par

% structure: p0085 source=h2 target=subsection reason=relative-heading-rank; base=h2

\subsection{四十二、对手不能否认他们自己所承认的那些恩宠恩赐的预定，尽管如此，他们的劝言并不受这预定的阻碍。}

我关于贞洁所说的，也可以同样用于信德、虔敬、爱德、恒心至终，而且，不必一一列举诸德，可以极其真实地用于所有对天主所行的服从。但是那些只将信德的开始和恒心至终放在我们的权力之内，而不视之为天主的恩赐，也不认为天主在我们的思念和意志上工作，使我们能够拥有并持守它们的人，却承认祂赐予其他事物——因为这些都是由信徒的信德从祂那里获得的。他们为何不害怕对这些其他事物的劝言和对这些其他事物的宣讲会受到预定定义的阻碍？或者，也许他们要说这类事物不是被预定的？那么它们就不是由天主赐予的，或者祂不知道祂要赐予它们。因为，如果它们是被赐予的，而且祂预知祂要赐予它们，那么祂当然预定了它们。因此，正如他们自己也劝勉贞洁、爱德、虔敬和其他他们承认是天主恩赐的事物，并且不能否认它们也被祂预知，因而也是被预定的；他们也不说他们的劝言受到天主预定之宣讲的阻碍，即受到天主对那些未来恩赐之预知的宣讲的阻碍：这样他们就可以看到，他们对信德或恒心至终的劝言也没有受到阻碍，即使那些事物本身，正如真理那样，可以说是天主的恩赐，并且那些事物是被预知，即被预定要赐予的；反而让他们看到，藉着这预定的宣讲，只有那种最有害的错误被阻碍和推翻，即说天主的恩宠是按照我们的功绩赐予的，以致那夸耀的人不是夸耀于主，而是夸耀于自己。\par

% structure: p0087 source=h2 target=subsection reason=relative-heading-rank; base=h2

\subsection{四十三、对上述论点的进一步阐述。}

为了向那些反应较慢的人更清楚地阐明这一点，请那些拥有敏捷悟性的人容忍我的迟缓。宗徒雅各伯说：\Quote{你们中若有人缺乏智慧，就该向那慷慨施恩于众人而从不责斥的天主祈求，智慧必会赐给他。}\BibleRef{雅 1:5} 在撒罗满的《箴言》中也记载着：\Quote{因为主赐予智慧。}\BibleRef{箴 2:6} 关于节制，在《智慧篇》中也有记载，这部书的权威被我们之前许多伟大的学者在注释神圣启示时引用过；那里写道：\Quote{我明知除非天主赐予，无人能节制；而知道这恩赐的来源，正是智慧的所在。}\BibleRef{智 8:21} 因此，这些都是天主的恩赐——且不说其他的，智慧与节制便是。那些人也该认同，因为他们并非白拉奇主义者，不会以顽固的异端邪说对抗如此明显的真理。\Quote{但是，}他们说，\Quote{这些由天主赐予我们的事物，是通过信德获得的，而信德源于我们自身。}他们认为拥有这种信德并坚持到底，都是我们自己的作为，仿佛不是从主领受的。这无疑与宗徒的话相悖，他说：\Quote{你有什么不是领受的呢？}\BibleRef{格前 4:7} 这也与殉道者西彼廉的话相悖：\Quote{我们不可夸耀任何事物，因为没有任何事物是属于我们自己的。}当我们说了这些以及其他许多令人厌倦的话，并表明信德的开始和坚持到底都是天主的恩赐，并且天主不可能预知任何他未来的恩赐——无论是所赐之物还是受赐之人——从而那些他所拯救和加冕的人都是他所预定的，他们便以为可以这样回应：\Quote{预定的断言与宣讲的利益相悖，因为人听到后，便不会被责罚的激励所感动。}他们说这话时，\Quote{不愿向人宣告：来到信德中并持守信德都是天主的恩赐，以免人们认为这暗示绝望而非鼓励，因为听者以为人的无知无法确知天主是将这些恩赐赐予何人，或是不赐予何人。}那么，为何他们自己却与我们一同宣讲智慧和节制是天主的恩赐呢？如果声明这些是天主的恩赐，并不妨碍我们劝勉人成为智慧和节制的努力，那又有什么理由认为，当我们劝勉人来相信并坚持到底时，即便这些也被证明是天主的恩赐（正如作为天主见证的圣经所证明的），这种劝勉会受到阻碍呢？\par

% structure: p0089 source=h2 target=subsection reason=relative-heading-rank; base=h2

\subsection{第44章——劝勉寻求智慧，纵使智慧是天主的恩赐。}

现在，暂且不谈节制，只在此讨论智慧。上文提到的雅各伯宗徒说：\Quote{至于从上而来的智慧，它首先是纯洁的，其次是和平的、宽仁的、柔顺的、满有怜悯和善果的、不偏不倚的、没有伪善的。}\BibleRef{雅 3:17} 我恳请你们看看，这智慧如何自光明之父降下，满载诸多伟大的恩惠？因为正如同一宗徒所说：\Quote{一切美好的赠与和完美的恩赐，都是从上而来，自光明之父降下的。}\BibleRef{雅 1:17} 那么，为何我们要责斥那些不洁和好争论的人，尽管我们向他们宣讲天主的恩赐是纯洁和平的智慧，却不担心他们因对天主旨意的不确定而在这宣讲中感到绝望多于鼓励，并且不会被责罚的激励所感动，反而反对我们，因为我们在他们尚未拥有那些我们所声称不是出于人的意愿、而是由天主慷慨赐予的事物时责斥他们？最后，为何这恩宠的宣讲没有阻止雅各伯宗徒责斥不安的灵魂，他说：\Quote{你们心里若怀有苦毒的嫉妒和纷争，就不可自夸，不可说谎话抵挡真理。这种智慧不是从上而来的，而是属地的、属血气的、属魔鬼的；因为哪里有嫉妒和纷争，哪里就有扰乱和各样的坏事。}\BibleRef{雅 3:14} 因此，不安静的人理应受到责斥，既有神圣宣告的证词，也有他们与我们共有的那些冲动；既然我们宣布那能纠正和医治纷争的和平智慧是天主的恩赐，这并不成为反对这种责斥的理由；不信的人同样应当受到责斥，如那些不持守信德的人，关于恩宠的宣讲并不妨碍这种责斥，尽管该宣讲称赞恩宠和持守恩宠是天主的恩赐。因为智慧是从信德获得的，正如雅各伯自己所说：\Quote{你们中若有人缺乏智慧，就该向那慷慨施恩于众人而从不责斥的天主祈求，智慧必会赐给他}\BibleRef{雅 1:5} 时，他立刻补充说：\Quote{不过他该怀着信心祈求，一点也不要疑惑。}然而，信德在求取者祈求之前就被赐予了，所以不能说信德不是天主的恩赐，而是出于我们自己，因为它是没有经过我们祈求就赐予我们的？因为宗徒非常明确地说：\Quote{愿平安、慈爱、信德，从天主父和主耶稣基督，赐与众弟兄。}\BibleRef{弗 6:23} 因此，平安和慈爱是从他而来，信德也是如此；所以我们向他祈求，不仅为使拥有信德者增加，也为使没有信德者获得信德。\par

% structure: p0091 source=h2 target=subsection reason=relative-heading-rank; base=h2

\subsection{第45章——同样劝勉寻求天主的其他恩赐。}

我说这些话所针对的那些人，他们呼喊说劝勉被预定和恩宠的宣讲所抑制，他们只劝勉那些他们争辩说不是由天主所赐、而是出自我们自己的恩赐，例如信德的开端和坚持到底。他们当然应该这样做，即只劝勉不信者相信，相信者继续相信。但那些他们与我们一同不否认是天主恩赐的事物，以便与我们一同摧毁白拉奇派的错误，例如谦逊、节制、忍耐以及其它与圣洁生活有关、并藉信德从主获得的德行，他们应当表明这些需要被祈求，并且只为自己或他人祈求；但他们不应劝勉任何人去追求和持守它们。然而，当他们根据自己的能力劝勉这些事物，并承认人应当被劝勉时，他们显然充分表明，劝勉并不因那种宣讲而受阻碍，无论是对信德的劝勉还是对坚持到底的劝勉，因为我们也宣讲这些事物是天主的恩赐，不是由任何人给予自己，而是由天主给予。\par

% structure: p0093 source=h2 target=subsection reason=relative-heading-rank; base=h2

\subsection{第四十六章．一个不坚持到底的人，因自己的过错而失败。}

但有人说：\\Quote{任何人若失足离开信德，是因自己的过错，当他屈服并同意那导致他离开信德的诱惑时。}谁否认这一点呢？但正因如此，不能说坚持到底的信德不是天主的恩赐。因为人每日祈求的正是此事，当他说：\\BibleQuote{不要让我们陷入诱惑，} \\BibleRef{玛 6:13} 若他被垂听，他所领受的正是此事。因此，当他每日祈求坚持时，他肯定将自己的希望不放在自己身上，而放在天主身上。然而，我不愿夸大其词，而宁可让他们思考并看看他们说服自己相信的是什么——即\\Quote{借着预定的宣讲，听众心中留下的更多是绝望而非劝勉。}因为这就是说，当一个人学会不将希望放在自己身上而放在天主身上时，他就对自己的救恩绝望了，尽管先知呼喊：\\BibleQuote{凡信赖世人的人，是可咒骂的。} \\BibleRef{耶 17:5}\par

% structure: p0095 source=h2 target=subsection reason=relative-heading-rank; base=h2

\subsection{第四十七章．预定有时以预知的名义被表达。}

因此，这些天主的恩赐，即赐给那些按天主的旨意蒙召的选民，其中包含信德的开端和坚持到底到生命终结，正如我通过理性与权威的一致见证所证明的——我说，这些天主的恩赐，如果不存在我所主张的预定，就不会被天主预知。但它们是被预知的。因此，这就是我所主张的预定。[XVIII.] 所以，有时同样的预定也以预知的名义被表达；正如宗徒所说：\\BibleQuote{天主并没有弃绝祂所预知的子民。} \\BibleRef{罗 11:2} 此处，当他说\\Quote{祂预知}时，除非理解为\\Quote{祂预定}，否则其含义不能正确理解，正如经文本身的上下文所示。因为他所谈论的是得救的犹太余民，其余的都丧亡了。因为上文他曾说先知对以色列宣告：\\Quote{我整天向一个悖逆抗辩的民族伸出我的手。} 然后，好像有人回答说：那么，天主对以色列的应许成了什么？他接着补充说：\\Quote{那么我说：天主难道弃绝了祂的子民吗？绝对没有！因为我也是一以色列人，出于亚巴郎的后裔，本雅明支派。} 然后他加上了我现在正在处理的话：\\BibleQuote{天主并没有弃绝祂所预知的子民。} \\BibleRef{罗 11:2} 为了显示余民是因天主的恩宠而留下，而非因他们行为的功绩，他继续补充说：\\Quote{难道你们不知道圣经在厄里亚篇中是如何说的吗？他怎样向天主控告以色列？}等等。\\Quote{但是，}他说，\\Quote{天主的回答对他说了什么？\\InnerQuote{我为自己保留了七千人，他们没有向巴耳屈膝。}} \\BibleRef{罗 11:5} 因为祂不是说：\\Quote{给我留下了}，或\\Quote{他们为我保留了自己，}而是\\BibleQuote{我为自己保留了。} \\BibleQuote{同样，现在这时刻也按恩宠的拣选留下了余民。如果是出于恩宠，就不再是因行为；否则恩宠就不再是恩宠了。} 将他上面引用的话与此连接起来：\\BibleQuote{那么怎么样？} 作为对这问题的回答，他说：\\BibleQuote{以色列所求的没有获得，但选民获得了，其余的被蒙蔽了。} \\BibleRef{罗 11:7} 因此，在拣选中，以及在这因恩宠的拣选而成为如此的余民中，他愿意使人理解天主没有弃绝的子民，因为祂预知了他们。这就是那拣选，藉此祂在基督内、在创世以前拣选了那些祂所愿意的人，使他们在祂面前成为圣洁无瑕，在爱中预定他们获得义子的地位。因此，凡是理解这些事的人，都不允许怀疑当宗徒说\\BibleQuote{天主并没有弃绝祂所预知的子民}时，他意指预定。因为祂预知了那余民，祂将按恩宠的拣选使他们如此。因此，祂预定了他们；因为毫无疑问，祂若预定，祂便预知；但预定就是预知祂将要做什么。\par

% structure: p0097 source=h2 target=subsection reason=relative-heading-rank; base=h2

\subsection{第四十八章 [XIX.]——西彼廉与盎博罗削的实践。}

那么，当我们读到一些圣经注释家论及天主的预知，并讨论选民的召叫时，有什么阻止我们将之理解为同一预定呢？因为他们或许更愿在此事上使用这个词，这个词更易理解，并且与那关于恩宠预定所宣告的真理并不矛盾，实则相符。我知道，没有人能够错误地驳斥我所依据圣经坚持的那预定。然而，我认为那些寻求注释家在此事上的意见者，应当满足于像\Person{西彼廉}和\Person{盎博罗削}这样圣洁且在各处信仰与基督道理中备受称赞的人，我已给出他们如此清楚的证词。并且为了这两种道理——即他们应当绝对相信并到处宣讲天主的恩宠是白白赐予的，正如我们必须相信并宣布的那样；并且他们不应以为这宣讲与那劝勉怠惰者或责备恶者的宣讲相矛盾；因为这些可称赞的人，虽然宣讲天主的恩宠到如此地步，以至于其中一位说：\Quote{我们不可夸耀任何事，因为没有任何事物属于我们自己；}另一位说：\Quote{我们的心和我们的思想不在我们权下；}却仍不停劝勉和责备，为使天主的诫命得以遵行。他们也不怕人对他们说：\Quote{你为何劝勉我们，为何责备我们，如果我们所有的一切美善都不出于我们，我们的心也不在我们权下？}这些圣人们绝不会害怕别人对他们说这样的话，因为他们明白：只有极少数人得以不经任何人宣讲，而由天主自身或天上的天使获得救恩的教导；但许多人得以通过人的媒介而相信天主。然而，无论天主的言语以何种方式向人说出，毫无疑问，人能听从它以致顺服，这是天主的恩赐。\par

% structure: p0099 source=h2 target=subsection reason=relative-heading-rank; base=h2

\subsection{第49章——对\Person{西彼廉}和\Person{盎博罗削}的进一步引用。}

因此，上述那些最卓越的神圣宣言注释家们，既宣讲了天主的真实恩宠，正应如此宣讲——即作为不带任何前人功劳的恩宠，——又迫切劝勉人行天主的诫命，好使那些拥有顺服恩赐的人听到他们应当服从什么命令。因为如果我们的任何功劳先于恩宠，那必定是某种行为、言语或思想的功劳，其中也包含了善意志本身。但那非常简要地总结了所有功劳种类的人说：\Quote{我们不可荣耀任何事，因为没有任何事物属于我们自己。}而那位说\Quote{我们的心和我们的思想不在我们权下}的人，也没有忽略行为和言语，因为人的任何行为或言语无不出自心与思想。然而，那位最光荣的殉道者和最光耀的圣师\Person{西彼廉}对此事还能说出什么呢？他告诫我们在\Emph{天主经}中甚至要为基督信仰的仇敌祈祷，借此表明他对信仰的开端也认为是天主的恩赐，并指出基督的教会每日为坚持到底而祈祷，因为只有天主赐予那坚持者以坚持到底之恩。此外，真福\Person{盎博罗削}在解释圣史\Person{路加}所说\Quote{我也好像合宜}的段落时，\BibleRef{路 1:3}说：\Quote{他宣称自己好像合宜，不能只是他一人觉得合宜。因为不是凭人的意志独自觉得合宜，而是按那在我内说话的\Person{基督}所喜悦的，祂成就那善事，也使我们觉得合宜：因为祂怜悯谁，也召叫谁。因此，那跟随基督的人，当被问及为何愿意成为基督徒时，可以回答：\InnerQuote{我也好像合宜。}他说这话时，并不否认这是天主所喜悦的；因为人的意志是由天主预备的。天主因圣人受尊崇，这正是天主的恩宠。}此外，在同一著作中——即同一福音的解释——当他讲到撒玛黎雅人不肯接待主，因为祂的面容朝向耶路撒冷去时，他说：\Quote{同时要知道，祂不被那些心地不单纯而悔改的人所接待。因为祂若愿意，就能使不虔诚的人变为虔诚。为什么他们不肯接待祂，圣史本人已提及，说：\InnerQuote{因为祂的面容朝向耶路撒冷去。}\BibleRef{路 9:53}但门徒们切望被撒玛黎雅人接待。然而，天主召叫那些祂使之堪当的人，并使祂所愿意的人成为虔诚的。}我们若乐于从圣经注释家那里听到圣经中明确的事，还能要求比这更清楚、更显明的吗？但除了这两位，我们还应加上第三位，即圣\Person{额我略}，他见证相信天主和宣认我们所信的都是天主的恩赐，他说：\Quote{我恳求你们宣认一位神性的天主圣三；但若你们愿意，也可说同一性体，并恳求天主使圣神赐予你们声音；}也就是说，要恳求天主允许赐给你们声音，用以宣认你们所相信的。\Quote{因为祂必定会赐予，我确定。那赐予第一的，也会赐予第二的。}那赐予相信的，也将赐予宣认。\par

% structure: p0101 source=h2 target=subsection reason=relative-heading-rank; base=h2

\subsection{第50章——宣讲天主的恩赐不阻碍顺服。}

这样的教父，而且如此伟大，当他们说：没有什么可让我们自夸为属于我们自己而非天主所赐的，甚至连我们的心和我们的思想也不在我们自己的能力之内；当他们将一切归于天主，并承认我们从祂领受了归向祂的能力，乃至能持续——那美好的事也显为美好，并使我们渴望它——我们尊崇天主、接受基督——从不虔敬之人变为虔敬和热心宗教之人——我们相信天主圣三本身，并也用口舌宣认我们所信的：——无疑，他们将这一切归于天主的恩宠，承认它们是祂的恩赐，并见证这些来自祂而非出于我们自己。但有人会说，他们如此承认天主的恩宠，竟敢否认祂的预知吗？不仅学者，就连未受教育的人也承认预知。再者，如果他们知道天主赐予这些恩赐，而且他们也知道祂预知祂将赐予这些，并且祂不可能不知道赐予谁：那么毫无疑问，他们已深知那预定，就是宗徒所宣讲、我们辛苦而勤勉地维护以反对现代异端者的预定。若有人对他们说：\Quote{如果你们不希望我们被激励去服从的热情在我们心中冷却，就不要向我们宣讲天主的恩宠——你们承认天主赐予那你们劝勉我们要去行的服从。}——这样说，绝无公义可言。\par

% structure: p0103 source=h2 target=subsection reason=relative-heading-rank; base=h2

\subsection{第51章 [XX.]——预定必须被宣讲。}

因此，如果宗徒和继承他们并效法他们的教会教师都做到了这两件事——即，既真正宣讲了不按我们功绩而赐予的天主恩宠，又通过健全的诫命谆谆教诲虔敬的服从——那么，我们时代这些人认为他们完全有义务因着真理不可战胜的力量而说：\Quote{即使关于天主恩惠的预定所说的是真的，也不应向民众宣讲}？它绝对应当被宣讲，使那有耳可听的，就听。谁有耳朵呢？除非他从那说\BibleQuote{我必赐给他们认识我的心，和能听的耳朵？}的祂领受。\BibleRef{巴 2:31} 无疑，那没有领受的，可以拒绝；而那领受的，可以取和喝，可以喝而活。因为正如虔敬必须被宣讲，使有耳可听的人能正确地敬拜天主；端庄必须被宣讲，使有耳可听的人不会因肉体本性而做出非法行为；爱德必须被宣讲，使有耳可听的人能爱天主和邻人；——同样，天主的恩惠的预定也必须如此宣讲，使有耳可听的人能夸耀，不是夸耀自己，而是夸耀主。\par

% structure: p0105 source=h2 target=subsection reason=relative-heading-rank; base=h2

\subsection{第52章——先前的著作预先驳斥了白拉奇异端。}

但是对于他们说：\Quote{没有必要让这么多才智不高的人因这种争执的不确定性而感到不安，因为公教信仰已在如此多年中被捍卫，不无益处，无需这预定的定义，既反对其他人，尤反对白拉奇派，在众多先前的著作中，包括公教徒和他人以及我们自己的；}——我感到十分惊奇，他们竟这样说，且不注意到——暂且不说其他著作——我那些专著正是在白拉奇派开始出现之前就撰写并出版的；他们没有看到那些专著中有多少地方，我无意间砍倒了未来的白拉奇异端，通过宣讲那恩宠——天主凭其无偿的仁慈，在我们任何先功绩之外，从那邪恶的谬误和我们的习气中拯救我们。在我就任主教初期，写给那位有福纪念的米兰教会主教辛普利齐亚诺的辩论中，我更充分地领会了这一点，那时我也领悟并断言：信德的开端是天主的恩赐。\par

% structure: p0107 source=h2 target=subsection reason=relative-heading-rank; base=h2

\subsection{第53章——奥古斯丁的\WorkTitle{忏悔录}。}

我哪一部较小的著作能比我的\WorkTitle{忏悔录}诸卷更普遍、更愉快地被认识呢？尽管我在白拉奇异端出现之前就已出版它们，但在其中我确实对我的天主说，而且多次说：\Quote{赐予你所命的，命令你所愿的。}白拉奇在罗马，当一位兄弟，也是我的同僚主教，在他面前提起我的这些话时，他无法忍受；且过于激动地反驳，几乎与提起这些话的人发生争吵。但天主首要且主要命令的，难道不就是我们相信祂吗？因此，如果人对祂说\Quote{赐予你所命的。}这本身也是祂所赐的。此外，在同一部书中，关于我归化的叙述，当天主使我归向那信德时——我正以最可怜、最狂怒的喋喋不休摧毁这信德——你难道不记得那是如何叙述的吗？我表明我蒙恩于我母亲忠信且每日的眼泪，使我不致丧亡。在那里我明确宣告，天主以祂的恩宠使人的意志转向真信德，这意志不仅回避它，甚至敌对它。此外，我如何恳求天主关于我恒心直到末后的增长，你知道，如果你愿意，可以查阅。因此，天主预知祂将赐予我在那部作品中曾祈求或赞美的一切恩赐，并且祂绝不可能不知道祂将赐予的对象是谁，谁敢——我不说否认——甚至怀疑呢？这就是圣人确凿无疑的预定，后来我因必要更仔细、更费力地加以辩护，当我已与白拉奇派辩论时。因为我学到，每个特殊的异端都将其独特的问题引入教会，为应对这些，圣经得以比没有这种必要迫使辩护时更仔细地被辩护。而正是白拉奇派说天主的恩宠是按我们的功绩赐予的，这若不是对恩宠的绝对否认，又是什么呢？是什么迫使我在劳苦中更充分、更明确地辩护那些赞美预定的圣经段落呢？\par

% structure: p0109 source=h2 target=subsection reason=relative-heading-rank; base=h2

\subsection{第五十四章 [XXI.]——信德的起始与终末属于天主。}

因此，为使这种不讨天主喜悦、敌对使我们得救的那些白白恩惠的见解被摧毁，我主张：按圣经——我已引用许多——信德的起始和其中的恒心，甚至直到末后，都是天主的恩赐。因为如果我们说信德的起始出于我们自己，从而我们配得接受天主其他恩赐，白拉奇派就会得出结论：天主的恩宠是按我们的功绩赐予的。公教信仰如此畏惧这一点，以致白拉奇本人，因害怕被定罪，而谴责了它。此外，如果我们说我们的恒心出于我们自己，而非出于天主，他们就会回答说：我们信德的起始出于我们自己，正如其终结一样，从而论证说：如果我们有持续直到末后的能力出于我们自己，那么我们的起始更出于我们自己，因为完成比开始大得多；于是他们反复得出结论：天主的恩宠是按我们的功绩赐予的。但如果两者都是天主的恩赐，并且天主预知祂将赐予这些祂的恩赐（谁能否认这一点呢？），就必须宣讲预定，好使天主的真恩宠，即不按功绩赐予的恩宠，得以不可战胜地辩护。\par

% structure: p0111 source=h2 target=subsection reason=relative-heading-rank; base=h2

\subsection{第五十五章——他先前著作与书信的见证。}

事实上，在那本标题为\WorkTitle{论责斥与恩宠}的论文中——该论文不能令所有爱我的人满足——我认为我已如此确立即使坚持到底也是天主的恩赐，正如我在此之前要么从未、要么几乎从未如此明确和突出地用文字阐述过这一点，除非我的记忆欺骗了我。但我现在以一种无人曾说过的方式说出了这一点。诚然，真福西彼廉在\WorkTitle{天主经}中，如我早已指出的，这样解释我们的祈求，以至于他说我们在该祈祷的第一祈求中就是在求坚持，因为我们祈求时所说的\Quote{愿祢的名被尊为圣}\BibleRef{玛 6:9}，尽管我们在圣洗中已得圣化——好使我们能坚持我们已经开始成为的。但是，那些因着对我的爱而我不应忘恩负义的人——他们声称，除了辩论中出现的内容之外，他们接受我所有的观点，正如你所写的——让他们看看，在我主教任期之初、佩拉基亚异端出现之前写给米兰主教辛普利西安的那两卷书的第一卷的后半部分，是否还留下任何可以质疑天主的恩宠不是按照我们的功劳而赐予的东西；我是否没有在那里充分论证甚至信仰的开端也是天主的恩赐；以及从那里说的话是否必然得出结论，虽然没有明说，但即使坚持到底，除非由那位预定我们到祂的王国和荣耀的天主赐予，否则也不会赐予。那么，我不是多年前就发表了我写给诺拉主教圣保利诺的那封反对佩拉基亚人的信吗？他们最近开始反驳它了。让他们也看看我写给罗马教会司铎西斯笃的那封信，当时我们与佩拉基亚人进行非常激烈的斗争，他们会发现那封信与给保利诺的信一样。由此他们可以推断出，类似的事情在几年前就已经说过和写过，反对佩拉基亚异端，而现在这些事竟令他们不悦，这实在令人惊讶；尽管我希望没有人会如此接受我所有的观点，以至追随我，除非在他应该看到我没有出错的事情上。因为我现在正在写论文，在其中我承诺要修订我较小的作品，以证明即使我自己也不是在所有事情上都追随自己；但我认为，在天主的仁慈下，我是渐进地写作，而不是从完美开始；因为，的确，如果现在我说在我这个年纪我终于达到了完美，在我所写的东西中没有任何错误，那我说得比事实更傲慢。但区别在于错误的大小和主题，以及任何人纠正错误的容易程度，或为错误辩护的固执程度。当然，对于那个今生最后一天发现他正在进步、以致于在他进步中所欠缺的一切都得以补足、并且他被判定为更适合被完善而非受惩罚的人，是有美好希望的。\par

% structure: p0113 source=h2 target=subsection reason=relative-heading-rank; base=h2

\subsection{第五十六章  天主赐予方法也赐予目的}

因此，如果我不愿对我所爱的人忘恩负义，因为我的辛劳的某些益处在他们爱我之前就已达到他们，那么我更加不愿对天主忘恩负义，我们不应该爱祂，除非祂先爱了我们并使我们爱祂；因为爱出于祂\BibleRef{若一 4:7}，正如那些祂使其不仅成为祂伟大的爱人，而且成为祂伟大的宣讲者的人所说的。还有什么比否认天主的恩宠本身，说恩宠是按照我们的功劳赐予我们的，更忘恩负义呢？公教信仰对佩拉基亚人的这一点感到战栗，并以此作为重罪归咎于佩拉基亚本人；而佩拉基亚本人也谴责了这一点，并非出于对天主真理的爱，而是出于对自己被定罪的恐惧。但是，无论谁作为忠诚的公教徒，都厌恶说天主的恩宠是按照我们的功劳赐予我们的，让他不要从使他获得怜悯而成为忠信的天主恩宠中除去信仰本身；这样，他也将坚持到底归因于天主的恩宠，他由此获得他每日祈求的怜悯，不要被引入诱惑。但是在信仰的开端和坚持的圆满之间，有那些我们赖以正义生活的方法，他们自己也同意这些方法是在信德的祈祷中由天主赐予我们的。而这一切——即信仰的开端，以及祂的其他恩赐直到终末——天主预知祂将赐予祂所召叫的人。因此，反驳预定或对预定产生怀疑，是过于好争辩了。\par

% structure: p0115 source=h2 target=subsection reason=relative-heading-rank; base=h2

\subsection{第五十七章 [第二十二章]——如何宣讲预定而不引起反感}

然而，不可如此向会众宣讲此教理，以免给不熟练的群众或较迟钝的百姓造成仿佛该宣讲本身在某种程度上被驳斥的印象。正如天主的预知——人们当然无法否认——若对他们说\Quote{无论你奔跑或睡觉，你都将成为那位不能受欺骗者所预知你将是的}，似乎就被驳倒了。即使是有用的药剂，若由一个欺骗性的或不熟练的医生调配，也可能要么无益，要么有害。但必须说：\Quote{你们要这样跑，好能获得奖品}\BibleRef{格前 9:24}；并由此，通过你们自己的奔跑，你们可以知道自己是那些应合法奔跑而被预知的人；以及任何其他可以如此宣讲天主的预知的方式，以便击退人的懒惰。\par

% structure: p0117 source=h2 target=subsection reason=relative-heading-rank; base=h2

\subsection{第五十八章——教理应有区别地应用}

因此，如今，关于预定的天主旨意的明确决定是这样：有些人从不信中获得服从的意愿，转而信德或恒心于信德；而另一些人停留在可咒之罪的快乐中，即使他们已被预定，却仍未兴起，因为怜悯恩宠的帮助尚未扶起他们。因为若有人尚未蒙召，而天主已藉其恩宠预定他们蒙选，他们将获得那恩宠，从而愿意蒙选，也得以如此；若有人服从，但未被预定进入天主的国和荣耀，他们只是暂时的，不会在同样的服从里坚持到底。然而，尽管这些是真实的，却不应对听众如此宣讲，以免他们以为这话是针对自己，也不应对他们说你在信中所写、我上文引用的那些话：\Quote{关于预定的天主旨意的明确决定是这样：你们中有些人将从不信中获得服从的意愿，并来到信德。}有何必要说\Quote{你们中有些}？因为我们若是向天主的教会说话，向信徒说话，为何说\Quote{他们中有些}已来到信德，似乎亏待了其余的人？我们更适宜的说法是：关于预定的天主旨意的明确决定是这样：你们将从不信中获得服从的意愿，并来到信德，并接受恒心，坚持到底。\par

% structure: p0119 source=h2 target=subsection reason=relative-heading-rank; base=h2

\subsection{第五十九章 应避免的冒犯}

以下话语也绝不可说——即：\Quote{但你们中其他停留在罪之快乐中的人尚未兴起，因为怜悯恩宠的帮助尚未扶起你们；}而可以而且应当合宜地这样说：\Quote{但你们中若有人仍耽延于可咒之罪的愉悦中，请握住最有益的纪律；然而，当你这样做之后，不要自高，仿佛出于自己的工作，也不要夸口，仿佛不是领受的。因为是天主在你们内运作，为成就祂的美意，使你们愿意并实行\BibleRef{斐 1:13}，并且你们的脚步由主引导，使你选择祂的道路。但对于你们自己的良善和正义的进程，要谨慎学习，它应归于神圣恩宠的预定。}\par

% structure: p0121 source=h2 target=subsection reason=relative-heading-rank; base=h2

\subsection{第六十章 对普世教会的应用}

此外，接着所说的：\Quote{但你们中若有人尚未蒙召，而天主藉其恩宠已预定他们蒙召，你们将获得那恩宠，藉此你们将愿意蒙选，并得以蒙选，}比我们考虑到我们并非对一般人说话，而是对基督的教会说话时所能说的更为严厉。因为为何不可以这样说：\Quote{如果你们中有人尚未蒙召，让我们为他们祈祷，使他们得以蒙召。因为他们或许如此被预定，以至于我们的祈祷得到应允，并获得那恩宠，藉此他们愿意并成为蒙选的人？}因为天主，祂完成了所有祂预定的事，也愿意我们为信德的仇敌祈祷，好使我们由此明白，祂自己也赐予不信者信德的恩赐，并使不愿者成为愿意者。\par

% structure: p0123 source=h2 target=subsection reason=relative-heading-rank; base=h2

\subsection{第六十一章 使用第三人称而非第二人称}

但现在我惊奇，若基督教聚会中任何软弱的弟兄可以耐心地听与这些话相关联的内容，当对他们说：\Quote{如果你们中有人服从，若你们被预定遭弃绝，服从的能力将从你们那里撤回，使你们停止服从。}这话说出来除了诅咒，或在某种程度上预言灾祸之外，还有什么？然而，若有关那些不恒心的人需要说些什么，为何不能至少以我刚才所说的方式来说——首先，这话不是针对会众中听的人说的，而是对他们论及他人；即，不应说：\Quote{如果你们中有人服从，若你们被预定遭弃绝，}而应说：\Quote{若有人服从，}其余部分使用动词的第三人称而非第二人称？因为对听众当面辱骂，不是可说的，而是可憎的，并且极其粗暴可恶，当说话者对他们说：\Quote{如果你们中有人服从，并被预定遭弃绝，服从的能力将从你们那里撤回，使你们停止服从。}如果这样表达：\Quote{但若有人服从，却未被预定进入祂的国和荣耀，他们只是暂时的，不会在那种服从里坚持到底，}教义缺少了什么？不是同样既更真实又更合宜吗？这样我们似乎不是好像为他们渴望那么多，而是述说他人所憎恶、认为不属于自己的邪恶，因希望和祈求更好的事。但在他们认为必须那样说的方式中，几乎用同样的话也可以宣告关于天主预知相同的判断，他们当然不能否认，以至于说：\Quote{如果你们中有人服从，若你们被预知要遭弃绝，你们将停止服从。}毫无疑问，这非常真实，的确如此；但它非常丑陋，非常轻率，非常不合适，不是因为虚假的宣告，而是因为其宣告没有有益地应用于人类软弱的健康。\par

% structure: p0125 source=h2 target=subsection reason=relative-heading-rank; base=h2

\subsection{第六十二章 然而，应强调祈祷}

但是我认为，我所提及的在对会众宣讲预定时所应采用的方式，并不足以满足传道者的需要，除非他再加上以下这番话，或类似的说法：\Quote{因此，你们也应当从光明之父那里盼望在顺服中的恒心，一切美好的施与和完美的恩赐都从祂而来，\BibleRef{雅 1:17}并要在每日的祈祷中向祂求赐；你们在这样做时，应当相信你们并不被排除在祂子民的预定之外，因为正是祂亲自赐予你们行此事的能力。你们万不可对自己绝望，因为你们受命要寄望于祂，而非寄望于自己。凡寄望于人者，皆受诅咒；\BibleRef{耶 17:5}依靠主总比依靠人更好，因为所有投靠祂的人都有福了。怀着这样的希望，你们当以敬畏事奉主，并以战兢向祂欢跃。因为没有人能确知那不说谎的天主在永恒之前应许给恩许之子的永生——没有人，除非他在世上的试炼生涯结束之时。\BibleRef{约 7:1}但祂必使我们恒心持守于祂内，直到那生命的终结，因为我们每日向祂说：\InnerQuote{不要让我们陷入诱惑。}} \BibleRef{玛 6:13} 当这些事情以及诸如此类的事情被宣讲时，无论是对少数基督徒还是对教会的众多信友，我们为何害怕宣讲圣人的预定和天主的真实恩宠——即那并非按照我们的功劳而赐予的恩宠——正如圣经所宣告的那样？或者说，难道要担心一个人会在自己的希望被显示为寄托于天主时而对自己绝望，而不是该担心他若因过度骄傲和不幸而将希望寄托于自身时，反而会对自己绝望吗？\par

% structure: p0127 source=h2 target=subsection reason=relative-heading-rank; base=h2

\subsection{第63章 [XXIII.]——整个教会祈祷的见证。}

而且我希望那些心灵迟钝和软弱的人，他们不能，或尚且不能理解圣经或圣经的解释，他们不论听到或没有听到我们对这个问题的论证，都能更仔细地思考自己的祈祷，这些祈祷是教会从起初直到今世的终结一直使用并将继续使用的。因为关于这件事——我现在不仅被迫提及，甚至要保护和辩护以对抗这些新的异端者——教会在其祈祷中从未沉默，尽管在其论述中，它认为没有必要提出，因为没有对手迫使它这样做。因为在教会中，何时没有为不信者和反对者祈祷，使他们能相信呢？何时曾有信友拥有不信的朋友、邻居或妻子，却没有为他们向主恳求一颗顺服于基督信仰的心呢？又有谁没有为自己祈祷过，求主使自己能常住于主内呢？又有谁敢不仅用言语，甚至在心里，指责那位代表信友呼求主的司铎，若他曾说：\Quote{主啊，求祢赐给他们恒心持守在祢内直到终结！}而众人岂不更应以宣认的口和相信的心回应这样的祝福，说\Quote{阿们}吗？既然在主祷文中，信友们不为别的事祈祷，尤其当他们说出那句恳求：\Quote{不要让我们陷入诱惑}，除了求能从圣洁的顺服中恒心持守。因此，教会既是在这些祈祷中诞生并成长及已经成长，它也是在这样的信德中诞生并成长及已经成长；凭此信德，人们相信天主的恩宠并非按照领受者的功劳而赐予。因为，当然，教会不会为不信者祈祷求赐信德，除非它相信天主将人的悖逆和敌对意志都转向祂自己。教会也不会祈祷求能在基督的信德中恒心持守，不被世界的诱惑欺骗或胜过，除非它相信主将我们的心置于祂的权能之下，以至于那仅凭我们自己意志就不能保持的善，除非祂也在我们内工作，使我们愿意，否则我们也不能保持。因为如果教会确实向祂祈求这些事，却认为这些事是凭自己给自己的，那么它所使用的祈祷就不是真实的，而是敷衍的——但愿这事远离我们！因为如果一个人认为他所领受的是从自己而来，而非从主而来，他又怎能真正叹息、渴望从主那里领受他所祈求的呢？\par

% structure: p0129 source=h2 target=subsection reason=relative-heading-rank; base=h2

\subsection{第64章——圣神以何种方式为我们代祷，呼喊：\Quote{阿爸，父啊！}}

尤其因为\BibleQuote{我们不知道怎样祈求才算合适，}宗徒说，\BibleQuote{但是圣神亲自以不可言喻的叹息为我们转求；那洞悉人心的知道圣神的意愿，因为祂按照天主为圣徒转求。}\BibleRef{罗 8:26}\Quote{圣神亲自转求}是什么意思，不就是\Quote{促使转求}，\Quote{以不可言喻的叹息}不就是\Quote{真实的话}，既然圣神是真理？因为祂就是那位，宗徒在别处谈到祂说：天主派遣了祂圣子的圣神到我们心中，\BibleQuote{呼喊：\InnerQuote{阿爸，父啊！}}\BibleRef{迦 4:6}这里\Quote{呼喊}是什么意思？不就是\Quote{使呼喊}吗？用那种修辞手法，我们称使人快乐的日子为愉快的一天。他在别处也阐明这一点，说：\BibleQuote{因为你们没有再次领受奴隶的神而恐惧，而是领受了义子的圣神，在祂内我们呼喊：\InnerQuote{阿爸，父啊！}}\BibleRef{罗 8:15}他那里说\Quote{呼喊}，这里说\Quote{在祂内我们呼喊}；就是说，阐明他说\Quote{呼喊}的含义，也就是如我已解释过的\Quote{使呼喊}，当我们明白这本身也是天主的恩赐，即我们以真诚的心和属神的方式向天主呼喊。因此，让他们看到自己是如何错误，认为我们的寻找、祈求、叩门是出于我们自己，而非赐予我们；并说这是因为恩宠以我们的功绩为先；当我们祈求就领受、寻找就找到、叩门就为我们开门时，恩宠跟随功绩。他们不明白，我们祈祷本身也是天主的恩赐，即我们祈求、寻找、叩门。因为我们领受了义子的圣神，在祂内我们呼喊：\Quote{阿爸，父啊！}\Person{圣安博}也这样说过。因为他曾说：\Quote{向天主祈祷也是属神灵恩的工作，正如经上记载：\BibleQuote{除非在圣神内，没有人能说\InnerQuote{耶稣是主}。}}\par

% structure: p0131 source=h2 target=subsection reason=relative-heading-rank; base=h2

\subsection{第65章——教会的祈祷隐含教会的信德。}

因此，教会向主所求的这些事，从她开始存在时就一直祈求，天主预知祂会赐给祂所召选的人，以至于在预定本身中已经赐给了他们；正如宗徒毫不含糊地宣告。因为他在写给弟茂德的信中说：\BibleQuote{你要按照天主的能力，为福音同我一起劳苦，祂拯救了我们，以圣召召叫了我们，不是按照我们的行为，而是按照祂自己的旨意和恩宠，这恩宠在万世以前，在基督耶稣内赐给了我们，如今借着我们的救主耶稣基督的显现，显示了出来。}\BibleRef{弟后 1:8}因此，让那敢于说教会任何时候都不曾相信这预定和恩宠的真理（这真理如今正以更警醒的态度被持守，以对抗最近的异端者）的人，让他说：教会任何时候都不曾祈求，或没有真诚地祈求，既为不信者能相信，也为信徒能坚持到底。如果教会一直为这些益处祈求，那么它就一直相信它们确实是天主的恩赐；它也从未有权否认天主预知了它们。因此，基督的教会从未未能持守这预定的信德，这信德如今正以新的热忱被捍卫，以对抗这些现代的异端者。\par

% structure: p0133 source=h2 target=subsection reason=relative-heading-rank; base=h2

\subsection{第66章[XXIV]——摘要与劝勉}

但我还能说什么呢？我认为我已经充分——或者说过于充分——教导了，在主内的信德之始和在主内坚持到底，都是天主的恩赐。其他属于善生、借此正确敬拜天主的美善之事，就连那些我为撰写此书而写给他们的人，也承认是天主的恩赐。此外，他们不能否认天主预知了祂所有的恩赐，以及祂将要赐予的对象。因此，正如必须宣讲其他事，好使宣讲者能被人听从，同样也必须宣讲预定，好使听从这些事的人，不是在人——因此也不是在自身——而是在主内夸耀；因为这也是天主的诫命，而听从这条诫命——即\BibleQuote{夸耀的应当夸耀主}\BibleRef{格前 1:31}——，与其他诫命一样，是天主的恩赐。没有这恩赐的人，我也不怕说，无论他拥有其他什么恩赐，都是徒然。我们祈求白拉奇主义者能获得这恩赐，也祈求我们的弟兄更丰盛地获得它。因此，让我们不要精于争辩而疏于祈祷。让我们祈祷，亲爱的，让我们祈祷，愿恩宠的天主甚至赐予我们的仇敌，尤其赐予我们的弟兄和亲爱的，使他们明白并承认：在那巨大且不可言喻的毁灭中，我们所有人都一同堕落了，除了天主的恩宠，没有人得救；并且恩宠不是按照领受者的功绩作为负债而偿还，而是作为真正的恩宠白白赐予，没有任何先行的功绩。\par

% structure: p0135 source=h2 target=subsection reason=relative-heading-rank; base=h2

\subsection{第67章——预定的最高典型是耶稣基督}

但再没有比耶稣基督自己更卓越的\OriginalTerm{预定}{predestination}例子了，我在前一篇论文中已就此争论过；在本篇结尾，我选择着重强调这一点。我说，再没有比中保自己更卓越的预定例子了。若任何信徒希望彻底理解此教义，就让他默观祂，在他内，他也会找到自己。我说的是信徒，那在祂内相信并承认我们本性的真实人性者，尽管由于被天主圣言摄取而独一无二地提升为天主的独生子，以致那摄取者与被摄取者成了圣三中的一位。因为从人的摄取并未导致\OriginalTerm{四位一体}{Quaternity}，而是保持为圣三，因为那摄取以不可言喻的方式使得天主与人合为一位的真理。因为我们说基督不仅是天主，如\OriginalTerm{摩尼异端者}{Manichean heretics}所主张的；也不仅是人，如\OriginalTerm{佛提诺异端者}{Photinian heretics}所宣称的；也不是那样的人，以致缺少任何确实属于人性的东西——无论是灵魂，还是在灵魂本身内的理性心智，或是并非取自女人，而是由圣言转化改变而成的肉身——这三种虚假空洞的观念造成了\OriginalTerm{亚波利拿里异端者}{Apollinarian heretics}的三个不同派别；但我们说基督是真天主，由天主圣父所生，没有任何时间起点；祂也是真实的人，由人世母亲所生，在确定的时期满足时；祂的人性，由此祂小于圣父，并不减损祂的神性，由此祂与圣父相等。因为两者是一位基督——而且，祂就天主性而言最真实地说：\BibleQuote{我与父原是一体；}\BibleRef{若 10:30}，就人性而言最真实地说：\BibleQuote{父比我大。}\BibleRef{若 14:28}因此，祂从达味的后裔中造了这义人——他永不会不义——而没有祂任何先行意志的功绩，也正是祂从恶人中造义人，没有他们任何先行意志的功绩，使祂成为头，他们成为肢体。因此，祂造了那人，没有祂先前的功绩，既没有从祂的根源继承，也没有因祂的意志犯任何罪需要赦免，也照样造了信祂的人，没有他们先前的功绩，祂赦免他们所有的罪。祂造了祂，使祂从未有过也不会有邪恶的意志，也照样在祂的肢体中将邪恶的意志变为善良的意志。因此，祂预定了祂和我们，因为无论就祂作为我们的头，还是就我们作为祂的身体，祂都预知我们的功绩不会先行，而是祂的行动先行。\par

% structure: p0137 source=h2 target=subsection reason=relative-heading-rank; base=h2

\subsection{第六十八章 结语}

阅读此书的，若理解，就感谢天主；若不理解，就应祈祷，使能有内在导师，从祂临在而来知识明达。\BibleRef{箴 2:6}但那些以为我犯了错误的人，要再三仔细考虑这里所说的话，以免他们自己可能错了。当通过阅读我作品的人，我不仅变得更明智，甚至更成全时，我承认天主对我的恩惠；我特别寄望于教会的导师们，如果我的作品落到他们手中，而他们屈尊肯定它的话。\par

\SourceRecord{Translated by Peter Holmes and Robert Ernest Wallis, and revised by Benjamin B. Warfield.；Nicene and Post-Nicene Fathers, First Series；Vol. 5.；Edited by Philip Schaff.；Buffalo, NY: Christian Literature Publishing Co.,；1887.；<http://www.newadvent.org/fathers/15122.htm>.}
